
\Data\ID{1003020802}
\Updated{2020-07-15 03:07:12}
\Level{A}
\SubArea{Points and vectors}
\LANGen{
\Question{
\( S \) is the midpoint of \( AB \). \( A = [5; -7] \), \( S = [0; -9] \). Find the coordinates of \( B \).}
{\( B = [-5; -11] \)}{\( B = [5; -18]	\)}{\( B = \left[-\frac52; -\frac{11}2\right] \)}{\( B = [5; 11] \)}{}{}}
\LANGpl{
\Question{
\( S \) to środek  odcinka  \( AB \).  \( A = [5; -7] \), \( S = [0; -9] \). Określ współrzędne punktu \( B \).}
{\( B = [-5; -11] \)}{\( B = [5; -18]	\)}{\( B = \left[-\frac52; -\frac{11}2\right] \)}{\( B = [5; 11] \)}{}{}}
\LANGcs{
\Question{
\( S \) je středem úsečky \( AB \). \( A = [5; -7] \), \( S = [0; -9] \). Určete souřadnice bodu \( B \).}
{\( B = [-5; -11] \)}{\( B = [5; -18]	\)}{\( B = \left[-\frac52; -\frac{11}2\right] \)}{\( B = [5; 11] \)}{}{}}
\LANGsk{
\Question{
\( S \) je stredom úsečky \( AB \). \( A = [5; -7] \), \( S = [0; -9] \). Určte súradnice bodu \( B \).}
{\( B = [-5; -11] \)}{\( B = [5; -18]	\)}{\( B = \left[-\frac52; -\frac{11}2\right] \)}{\( B = [5; 11] \)}{}{}}
\LANGes{
\Question{
\( S \) es el centro de \( AB \). \( A = [5; -7] \), \( S = [0; -9] \). Determina las coordenadas de \( B \).}
{\( B = [-5; -11] \)}{\( B = [5; -18]	\)}{\( B = \left[-\frac52; -\frac{11}2\right] \)}{\( B = [5; 11] \)}{}{}}
\End

\Data\ID{1003020806}
\Updated{2020-07-15 03:07:12}
\Level{A}
\SubArea{Points and vectors}
\LANGen{
\Question{
Let \( ABC \) be a triangle with \( A = [1 ; 2 ; -3] \), \( B = [-3 ; 3 ; -2] \) and \( C = [-1 ; 1 ; -1] \). Find its perimeter.}
{\( 6+3\sqrt2 \)}{\( 18 \)}{\( 6 + 2\sqrt3 \)}{\( 9\sqrt2 \)}{}{}}
\LANGpl{
\Question{
Dany jest trójkąt \( ABC \), gdzie  \( A = [1 ; 2 ; -3] \), \( B = [-3 ; 3 ; -2] \) i  \( C = [-1 ; 1 ; -1] \). Oblicz obwód trójkąta.}
{\( 6+3\sqrt2 \)}{\( 18 \)}{\( 6 + 2\sqrt3 \)}{\( 9\sqrt2 \)}{}{}}
\LANGcs{
\Question{
Je dán trojúhelník \( ABC \), kde \( A = [1 ; 2 ; -3] \), \( B = [-3 ; 3 ; -2] \) a \( C = [-1 ; 1 ; -1] \). Vypočtěte jeho obvod.}
{\( 6+3\sqrt2 \)}{\( 18 \)}{\( 6 + 2\sqrt3 \)}{\( 9\sqrt2 \)}{}{}}
\LANGsk{
\Question{
Je daný trojuholník \( ABC \), kde \( A = [1 ; 2 ; -3] \), \( B = [-3 ; 3 ; -2] \) a \( C = [-1 ; 1 ; -1] \). Vypočítajte jeho obvod.}
{\( 6+3\sqrt2 \)}{\( 18 \)}{\( 6 + 2\sqrt3 \)}{\( 9\sqrt2 \)}{}{}}
\LANGes{
\Question{
Sea \( ABC \) un triángulo con  \( A = [1 ; 2 ; -3] \), \( B = [-3 ; 3 ; -2] \) y \( C = [-1 ; 1 ; -1] \). Calcula su perímetro.}
{\( 6+3\sqrt2 \)}{\( 18 \)}{\( 6 + 2\sqrt3 \)}{\( 9\sqrt2 \)}{}{}}
\End

\Data\ID{1003020807}
\Updated{2020-07-15 03:07:12}
\Level{A}
\SubArea{Points and vectors}
\LANGen{
\Question{
Find the set of all values of a parameter \( a\in\mathbb{R} \) such that  the distance of the points \( [a;3;-2] \) and \( [1;1;4] \)   equals  \( 7 \).}
{\( \{4;-2\} \)}{\( \{4;2\} \)}{\( \emptyset \)}{\( \{0\} \)}{}{}}
\LANGpl{
\Question{
Wskaż zbiór wszystkich wartości parametru \( a\in\mathbb{R} \) tak, aby odległość punktów \( [a;3;-2] \) i \( [1;1;4] \)  była równa  \( 7 \).}
{\( \{4;-2\} \)}{\( \{4;2\} \)}{\( \emptyset \)}{\( \{0\} \)}{}{}}
\LANGcs{
\Question{
Určete množinu všech takových hodnot parametru \( a\in\mathbb{R} \), pro které je vzdálenost bodů \( [a;3;-2] \) a \( [1;1;4] \) rovna \( 7 \).}
{\( \{4;-2\} \)}{\( \{4;2\} \)}{\( \emptyset \)}{\( \{0\} \)}{}{}}
\LANGsk{
\Question{
Určte množinu všetkých takých hodnôt parametra \( a\in\mathbb{R} \), pre ktoré je vzdialenosť bodov \( [a;3;-2] \) a \( [1;1;4] \) rovná \( 7 \).}
{\( \{4;-2\} \)}{\( \{4;2\} \)}{\( \emptyset \)}{\( \{0\} \)}{}{}}
\LANGes{
\Question{
Determina el conjunto de todos los valores de un parámetro \( a\in\mathbb{R} \) suponiendo que la distancia de los puntos \( [a;3;-2] \) y \( [1;1;4] \)  equivale a \( 7 \).}
{\( \{4;-2\} \)}{\( \{4;2\} \)}{\( \emptyset \)}{\( \{0\} \)}{}{}}
\End

\Data\ID{1003024306}
\Updated{2020-07-15 03:07:12}
\Level{A}
\SubArea{Points and vectors}
\LANGen{
\Question{
We are given the points A = [-4;2;3], B = [-5;6;3], D = [1;1;4]. Find the coordinates of a point \( C \), if: 
\[ \overrightarrow{u} = \overrightarrow{AB}\text{, }\ \overrightarrow{CD} = -\frac12\overrightarrow{u}\]}
{\( C = \left[\frac12; 3; 4\right] \)}{\( C = \left[-\frac12;-3;-4\right] \)}{\( C = \left[\frac32;3;4\right] \)}{\( C = \left[\frac32;-3;-4\right] \)}{}{}}
\LANGpl{
\Question{
Dane są punkty A = [-4;2;3], B = [-5;6;3], D = [1;1;4].  Określ współrzędne punktu \( C \), jeśli:
\[ \overrightarrow{u} = \overrightarrow{AB}\text{, }\ \overrightarrow{CD} = -\frac12\overrightarrow{u}\]}
{\( C = \left[\frac12; 3; 4\right] \)}{\( C = \left[-\frac12;-3;-4\right] \)}{\( C = \left[\frac32;3;4\right] \)}{\( C = \left[\frac32;-3;-4\right] \)}{}{}}
\LANGcs{
\Question{
Jsou dány body A = [-4;2;3], B = [-5;6;3], D = [1;1;4]. Určete souřadnice bodu \( C \) tak, aby platilo: 
\[ \overrightarrow{u} = \overrightarrow{AB}\text{, }\ \overrightarrow{CD} = -\frac12\overrightarrow{u}\]}
{\( C = \left[\frac12; 3; 4\right] \)}{\( C = \left[-\frac12;-3;-4\right] \)}{\( C = \left[\frac32;3;4\right] \)}{\( C = \left[\frac32;-3;-4\right] \)}{}{}}
\LANGsk{
\Question{
Sú dané body A = [-4;2;3], B = [-5;6;3], D = [1;1;4]. Určte súradnice bodov \( C \) tak, aby platilo: 
\[ \overrightarrow{u} = \overrightarrow{AB}\text{, }\ \overrightarrow{CD} = -\frac12\overrightarrow{u}\]}
{\( C = \left[\frac12; 3; 4\right] \)}{\( C = \left[-\frac12;-3;-4\right] \)}{\( C = \left[\frac32;3;4\right] \)}{\( C = \left[\frac32;-3;-4\right] \)}{}{}}
\LANGes{
\Question{
Dados los puntos A = [-4;2;3], B = [-5;6;3], D = [1;1;4]. Determina las coordenadas del punto \( C \), si: 
\[ \overrightarrow{u} = \overrightarrow{AB}\text{, }\ \overrightarrow{CD} = -\frac12\overrightarrow{u}\]}
{\( C = \left[\frac12; 3; 4\right] \)}{\( C = \left[-\frac12;-3;-4\right] \)}{\( C = \left[\frac32;3;4\right] \)}{\( C = \left[\frac32;-3;-4\right] \)}{}{}}
\End

\Data\ID{1003024307}
\Updated{2020-07-15 03:07:12}
\Level{A}
\SubArea{Points and vectors}
\LANGen{
\Question{
Let \( \overrightarrow{a} = (-1;2) \), \( \overrightarrow{b} = (2;1) \), \( \overrightarrow{c} = (-4;3) \). Express vector \( \overrightarrow{c} \) as a linear combination of vectors \( \overrightarrow{a} \) and \( \overrightarrow{b} \).}
{\( \overrightarrow{c} = 2\overrightarrow{a} - \overrightarrow{b} \)}{\( \overrightarrow{c} = 4\overrightarrow{a} - 8\overrightarrow{b} \)}{\( \overrightarrow{c} = 4\overrightarrow{a} - \overrightarrow{b} \)}{\( \overrightarrow{c} = -2\overrightarrow{a} + \overrightarrow{b} \)}{}{}}
\LANGpl{
\Question{
Podano \( \overrightarrow{a} = (-1;2) \), \( \overrightarrow{b} = (2;1) \), \( \overrightarrow{c} = (-4;3) \). Przedstaw wektor \( \overrightarrow{c} \) jako liniową kombinację wektorów \( \overrightarrow{a} \) i \( \overrightarrow{b} \).}
{\( \overrightarrow{c} = 2\overrightarrow{a} - \overrightarrow{b} \)}{\( \overrightarrow{c} = 4\overrightarrow{a} - 8\overrightarrow{b} \)}{\( \overrightarrow{c} = 4\overrightarrow{a} - \overrightarrow{b} \)}{\( \overrightarrow{c} = -2\overrightarrow{a} + \overrightarrow{b} \)}{}{}}
\LANGcs{
\Question{
Jsou dány vektory \( \overrightarrow{a} = (-1;2) \), \( \overrightarrow{b} = (2;1) \), \( \overrightarrow{c} = (-4;3) \). Vyjádřete vektor \( \overrightarrow{c} \) jako lineární kombinaci vektorů \( \overrightarrow{a} \) a \( \overrightarrow{b} \).}
{\( \overrightarrow{c} = 2\overrightarrow{a} - \overrightarrow{b} \)}{\( \overrightarrow{c} = 4\overrightarrow{a} - 8\overrightarrow{b} \)}{\( \overrightarrow{c} = 4\overrightarrow{a} - \overrightarrow{b} \)}{\( \overrightarrow{c} = -2\overrightarrow{a} + \overrightarrow{b} \)}{}{}}
\LANGsk{
\Question{
Sú dané vektory \( \overrightarrow{a} = (-1;2) \), \( \overrightarrow{b} = (2;1) \), \( \overrightarrow{c} = (-4;3) \). Vyjadrite vektor \( \overrightarrow{c} \) ako lineárnu kombináciu vektorov \( \overrightarrow{a} \) a \( \overrightarrow{b} \).}
{\( \overrightarrow{c} = 2\overrightarrow{a} - \overrightarrow{b} \)}{\( \overrightarrow{c} = 4\overrightarrow{a} - 8\overrightarrow{b} \)}{\( \overrightarrow{c} = 4\overrightarrow{a} - \overrightarrow{b} \)}{\( \overrightarrow{c} = -2\overrightarrow{a} + \overrightarrow{b} \)}{}{}}
\LANGes{
\Question{
Sean \( \overrightarrow{a} = (-1;2) \), \( \overrightarrow{b} = (2;1) \), \( \overrightarrow{c} = (-4;3) \). Expresa el vector \( \overrightarrow{c} \)  como la combinación lineal de vectores \( \overrightarrow{a} \) y \( \overrightarrow{b} \).}
{\( \overrightarrow{c} = 2\overrightarrow{a} - \overrightarrow{b} \)}{\( \overrightarrow{c} = 4\overrightarrow{a} - 8\overrightarrow{b} \)}{\( \overrightarrow{c} = 4\overrightarrow{a} - \overrightarrow{b} \)}{\( \overrightarrow{c} = -2\overrightarrow{a} + \overrightarrow{b} \)}{}{}}
\End

\Data\ID{1103020801}
\Updated{2020-07-15 03:07:12}
\Level{A}
\SubArea{Points and vectors}
\IMG{1103020801.pdf}
\LANGen{
\Question{
Find the coordinates of the midpoints of the line segments \( AB \), \( BC \), \( AC \). For coordinates of the points \( A \), \( B \) and \( C \), see the picture.}
{\( S_{AB}=\left[-\frac12;1 \right]\text{, }\ S_{BC}=[4;2 ]\text{, }\ S_{AC}=\left[\frac12; 4\right] \)}{\( S_{AB}=\left[-\frac32;2 \right]\text{, }\ S_{BC}=[1;3 ]\text{, }\ S_{AC}=\left[\frac52; 4\right] \)}{\( S_{AB}=\left[\frac12;1 \right]\text{, }\ S_{BC}=[4;2 ]\text{, }\ S_{AC}=\left[-\frac12; 4\right] \)}{\( S_{AB}=\left[1;-\frac12 \right]\text{, }\ S_{BC}=[2;4 ]\text{, }\ S_{AC}=\left[4;\frac12\right] \)}{}{}}
\LANGpl{
\Question{
Wskaż współrzędne środka odcinków \( AB \), \( BC \), \( AC \). Współrzędne punktów  \( A \), \( B \) i \( C \)  wskazano na rysunku.}
{\( S_{AB}=\left[-\frac12;1 \right]\text{, }\ S_{BC}=[4;2 ]\text{, }\ S_{AC}=\left[\frac12; 4\right] \)}{\( S_{AB}=\left[-\frac32;2 \right]\text{, }\ S_{BC}=[1;3 ]\text{, }\ S_{AC}=\left[\frac52; 4\right] \)}{\( S_{AB}=\left[\frac12;1 \right]\text{, }\ S_{BC}=[4;2 ]\text{, }\ S_{AC}=\left[-\frac12; 4\right] \)}{\( S_{AB}=\left[1;-\frac12 \right]\text{, }\ S_{BC}=[2;4 ]\text{, }\ S_{AC}=\left[4;\frac12\right] \)}{}{}}
\LANGcs{
\Question{
Určete souřadnice středů úseček \( AB \), \( BC \), \( AC \). Souřadnice bodů \( A \), \( B \), \( C \) naleznete v obrázku.}
{\( S_{AB}=\left[-\frac12;1 \right]\text{, }\ S_{BC}=[4;2 ]\text{, }\ S_{AC}=\left[\frac12; 4\right] \)}{\( S_{AB}=\left[-\frac32;2 \right]\text{, }\ S_{BC}=[1;3 ]\text{, }\ S_{AC}=\left[\frac52; 4\right] \)}{\( S_{AB}=\left[\frac12;1 \right]\text{, }\ S_{BC}=[4;2 ]\text{, }\ S_{AC}=\left[-\frac12; 4\right] \)}{\( S_{AB}=\left[1;-\frac12 \right]\text{, }\ S_{BC}=[2;4 ]\text{, }\ S_{AC}=\left[4;\frac12\right] \)}{}{}}
\LANGsk{
\Question{
Určte súradnice stredu úsečiek \( AB \), \( BC \), \( AC \). Súradnice bodov \( A \), \( B \), \( C \) nájdete v obrázku.}
{\( S_{AB}=\left[-\frac12;1 \right]\text{, }\ S_{BC}=[4;2 ]\text{, }\ S_{AC}=\left[\frac12; 4\right] \)}{\( S_{AB}=\left[-\frac32;2 \right]\text{, }\ S_{BC}=[1;3 ]\text{, }\ S_{AC}=\left[\frac52; 4\right] \)}{\( S_{AB}=\left[\frac12;1 \right]\text{, }\ S_{BC}=[4;2 ]\text{, }\ S_{AC}=\left[-\frac12; 4\right] \)}{\( S_{AB}=\left[1;-\frac12 \right]\text{, }\ S_{BC}=[2;4 ]\text{, }\ S_{AC}=\left[4;\frac12\right] \)}{}{}}
\LANGes{
\Question{
Determina las coordenadas de los puntos medios de los segmentos \( AB \), \( BC \), \( AC \). Para encontrar las coordenadas de los puntos \( A \), \( B \) y \( C \), mira la imagen.}
{\( S_{AB}=\left[-\frac12;1 \right]\text{, }\ S_{BC}=[4;2 ]\text{, }\ S_{AC}=\left[\frac12; 4\right] \)}{\( S_{AB}=\left[-\frac32;2 \right]\text{, }\ S_{BC}=[1;3 ]\text{, }\ S_{AC}=\left[\frac52; 4\right] \)}{\( S_{AB}=\left[\frac12;1 \right]\text{, }\ S_{BC}=[4;2 ]\text{, }\ S_{AC}=\left[-\frac12; 4\right] \)}{\( S_{AB}=\left[1;-\frac12 \right]\text{, }\ S_{BC}=[2;4 ]\text{, }\ S_{AC}=\left[4;\frac12\right] \)}{}{}}
\End

\Data\ID{1103020803}
\Updated{2020-07-15 03:07:12}
\Level{A}
\SubArea{Points and vectors}
\IMG{1103020803.pdf}
\LANGen{
\Question{
Find the coordinates of the centroid of the triangle \( ABC \). Triangle \( ABC \) is shown in the picture.}
{\( T=\left[1; \frac13\right] \)}{\( T=[1;0] \)}{\( T=[1;1] \)}{\( T=\left[\frac43; 0\right] \)}{}{}}
\LANGpl{
\Question{
Określ współrzędne środka ciężkości trójkąta  \( ABC \)  pokazanego na rysunku.}
{\( T=\left[1; \frac13\right] \)}{\( T=[1;0] \)}{\( T=[1;1] \)}{\( T=\left[\frac43; 0\right] \)}{}{}}
\LANGcs{
\Question{
Určete souřadnice těžiště trojúhelníka \( ABC \) znázorněného na obrázku.}
{\( T=\left[1; \frac13\right] \)}{\( T=[1;0] \)}{\( T=[1;1] \)}{\( T=\left[\frac43; 0\right] \)}{}{}}
\LANGsk{
\Question{
Určte súradnice ťažiska trojuholníka \( ABC \) znázorneného na obrázku.}
{\( T=\left[1; \frac13\right] \)}{\( T=[1;0] \)}{\( T=[1;1] \)}{\( T=\left[\frac43; 0\right] \)}{}{}}
\LANGes{
\Question{
Determina las coordenadas del baricentro del triángulo \( ABC \). El triángulo \( ABC \) es mostrado en la imagen.}
{\( T=\left[1; \frac13\right] \)}{\( T=[1;0] \)}{\( T=[1;1] \)}{\( T=\left[\frac43; 0\right] \)}{}{}}
\End

\Data\ID{1103020804}
\Updated{2020-07-15 03:07:12}
\Level{A}
\SubArea{Points and vectors}
\IMG{1103020804.pdf}
\LANGen{
\Question{
In the parallelogram \( ABCD \) shown in the picture, \( G \) is the midpoint of \( CD \), \( F \) is the midpoint of \( BC \) and  \( \overrightarrow{u}=\overrightarrow{CG} \), \( \overrightarrow{v}=\overrightarrow{CF} \), \( \overrightarrow{a}=\overrightarrow{AD} \)  and  \( \overrightarrow{b}=\overrightarrow{AC} \). 
Express vectors \( \overrightarrow{a} \)  and \( \overrightarrow{b} \) as a linear combination of vectors \( \overrightarrow{u} \) and \( \overrightarrow{v} \).}
{\( \overrightarrow{a}=-2\overrightarrow{v};\ \overrightarrow{b}=-2\overrightarrow{u}-2\overrightarrow{v} \)}{\( \overrightarrow{a}=\overrightarrow{b}+2\overrightarrow{u};\ \overrightarrow{b}=-2\overrightarrow{u}+2\overrightarrow{v} \)}{\( \overrightarrow{a}=\overrightarrow{b}-2\overrightarrow{u};\ \overrightarrow{b}=-\sqrt2\overrightarrow{u}-\sqrt2\overrightarrow{v} \)}{\( \overrightarrow{a}=-2\overrightarrow{v};\ \overrightarrow{b}=2\overrightarrow{u}+2\overrightarrow{v} \)}{}{}}
\LANGpl{
\Question{
Dany jest równoległobok \( ABCD \), gdzie \( G \) to środek \( CD \), \( F \) to środek \( BC \) i  \( \overrightarrow{u}=\overrightarrow{CG} \), \( \overrightarrow{v}=\overrightarrow{CF} \), \( \overrightarrow{a}=\overrightarrow{AD} \)  i  \( \overrightarrow{b}=\overrightarrow{AC} \). 
Przedstaw wektory \( \overrightarrow{a} \)  i \( \overrightarrow{b} \) jako liniową kombinację wektorów \( \overrightarrow{u} \) i \( \overrightarrow{v} \).}
{\( \overrightarrow{a}=-2\overrightarrow{v};\ \overrightarrow{b}=-2\overrightarrow{u}-2\overrightarrow{v} \)}{\( \overrightarrow{a}=\overrightarrow{b}+2\overrightarrow{u};\ \overrightarrow{b}=-2\overrightarrow{u}+2\overrightarrow{v} \)}{\( \overrightarrow{a}=\overrightarrow{b}-2\overrightarrow{u};\ \overrightarrow{b}=-\sqrt2\overrightarrow{u}-\sqrt2\overrightarrow{v} \)}{\( \overrightarrow{a}=-2\overrightarrow{v};\ \overrightarrow{b}=2\overrightarrow{u}+2\overrightarrow{v} \)}{}{}}
\LANGcs{
\Question{
V rovnoběžníku \( ABCD \) jsou vyznačeny body \( G \) - střed \( CD \), \( F \) - střed \( BC \) a vektory  \( \overrightarrow{u}=\overrightarrow{CG} \), \( \overrightarrow{v}=\overrightarrow{CF} \), \( \overrightarrow{a}=\overrightarrow{AD} \)  a  \( \overrightarrow{b}=\overrightarrow{AC} \). 
Vyjádřete vektory \( \overrightarrow{a} \)  a \( \overrightarrow{b} \) jako lineární kombinaci vektorů \( \overrightarrow{u} \) a \( \overrightarrow{v} \).}
{\( \overrightarrow{a}=-2\overrightarrow{v};\ \overrightarrow{b}=-2\overrightarrow{u}-2\overrightarrow{v} \)}{\( \overrightarrow{a}=\overrightarrow{b}+2\overrightarrow{u};\ \overrightarrow{b}=-2\overrightarrow{u}+2\overrightarrow{v} \)}{\( \overrightarrow{a}=\overrightarrow{b}-2\overrightarrow{u};\ \overrightarrow{b}=-\sqrt2\overrightarrow{u}-\sqrt2\overrightarrow{v} \)}{\( \overrightarrow{a}=-2\overrightarrow{v};\ \overrightarrow{b}=2\overrightarrow{u}+2\overrightarrow{v} \)}{}{}}
\LANGsk{
\Question{
V rovnobežníku \( ABCD \) sú vyznačené body \( G \) - stred \( CD \), \( F \) - stred \( BC \) a vektory  \( \overrightarrow{u}=\overrightarrow{CG} \), \( \overrightarrow{v}=\overrightarrow{CF} \), \( \overrightarrow{a}=\overrightarrow{AD} \)  a  \( \overrightarrow{b}=\overrightarrow{AC} \). 
Vyjadrite vektory \( \overrightarrow{a} \)  a \( \overrightarrow{b} \) ako lineárna kombinácia vektorov \( \overrightarrow{u} \) a \( \overrightarrow{v} \).}
{\( \overrightarrow{a}=-2\overrightarrow{v};\ \overrightarrow{b}=-2\overrightarrow{u}-2\overrightarrow{v} \)}{\( \overrightarrow{a}=\overrightarrow{b}+2\overrightarrow{u};\ \overrightarrow{b}=-2\overrightarrow{u}+2\overrightarrow{v} \)}{\( \overrightarrow{a}=\overrightarrow{b}-2\overrightarrow{u};\ \overrightarrow{b}=-\sqrt2\overrightarrow{u}-\sqrt2\overrightarrow{v} \)}{\( \overrightarrow{a}=-2\overrightarrow{v};\ \overrightarrow{b}=2\overrightarrow{u}+2\overrightarrow{v} \)}{}{}}
\LANGes{
\Question{
En el paralelogramo \( ABCD \) que es mostrado en la imagen, \( G \) es el centro de \( CD \), \( F \) es el centro de \( BC \) y  \( \overrightarrow{u}=\overrightarrow{CG} \), \( \overrightarrow{v}=\overrightarrow{CF} \), \( \overrightarrow{a}=\overrightarrow{AD} \)  y  \( \overrightarrow{b}=\overrightarrow{AC} \). 
Expresa vectores \( \overrightarrow{a} \)  y \( \overrightarrow{b} \) como la combinación lineal de vectores\( \overrightarrow{u} \) y \( \overrightarrow{v} \).}
{\( \overrightarrow{a}=-2\overrightarrow{v};\ \overrightarrow{b}=-2\overrightarrow{u}-2\overrightarrow{v} \)}{\( \overrightarrow{a}=\overrightarrow{b}+2\overrightarrow{u};\ \overrightarrow{b}=-2\overrightarrow{u}+2\overrightarrow{v} \)}{\( \overrightarrow{a}=\overrightarrow{b}-2\overrightarrow{u};\ \overrightarrow{b}=-\sqrt2\overrightarrow{u}-\sqrt2\overrightarrow{v} \)}{\( \overrightarrow{a}=-2\overrightarrow{v};\ \overrightarrow{b}=2\overrightarrow{u}+2\overrightarrow{v} \)}{}{}}
\End

\Data\ID{1103020805}
\Updated{2020-07-15 03:07:12}
\Level{A}
\SubArea{Points and vectors}
\IMG{1103020805.pdf}
\LANGen{
\Question{
We are given the points \( A = [1;1;4] \), \( C = [0;4;7] \) and \( D = [2;0;5] \) as seen in the picture. 
What are the coordinates of a point \( B \), if \( ABCD \) is a parallelogram?}
{\( B = [-1;5;6] \)}{\( B = [3;-3;2] \)}{\( B = [-2;4;3] \)}{\( B = [-3;3;-2] \)}{}{}}
\LANGpl{
\Question{
Dane są punkty \( A = [1;1;4] \), \( C = [0;4;7] \) i \( D = [2;0;5] \). 
Wskaż współrzędne punktu  \( B \) tak, aby \( ABCD \) był równoległobokiem.}
{\( B = [-1;5;6] \)}{\( B = [3;-3;2] \)}{\( B = [-2;4;3] \)}{\( B = [-3;3;-2] \)}{}{}}
\LANGcs{
\Question{
Jsou dány body \( A = [1;1;4] \), \( C = [0;4;7] \) a \( D = [2;0;5] \). 
Určete souřadnice bodu \( B \) tak, aby \( ABCD \) byl rovnoběžník.}
{\( B = [-1;5;6] \)}{\( B = [3;-3;2] \)}{\( B = [-2;4;3] \)}{\( B = [-3;3;-2] \)}{}{}}
\LANGsk{
\Question{
Sú dané body \( A = [1;1;4] \), \( C = [0;4;7] \) a \( D = [2;0;5] \). 
Určte súradnice bodu \( B \) tak, aby \( ABCD \) bol rovnobežník.}
{\( B = [-1;5;6] \)}{\( B = [3;-3;2] \)}{\( B = [-2;4;3] \)}{\( B = [-3;3;-2] \)}{}{}}
\LANGes{
\Question{
Dados los puntos \( A = [1;1;4] \), \( C = [0;4;7] \) y \( D = [2;0;5] \) que se pueden ver en la imagen. 
Cuáles son las coordenadas de un punto \( B \), si \( ABCD \) es el paralelogramo?}
{\( B = [-1;5;6] \)}{\( B = [3;-3;2] \)}{\( B = [-2;4;3] \)}{\( B = [-3;3;-2] \)}{}{}}
\End

\Data\ID{1103020808}
\Updated{2020-07-15 03:07:12}
\Level{A}
\SubArea{Points and vectors}
\IMG{1103020808.pdf}
\LANGen{
\Question{
Let \( ABC  \) be a triangle. In the picture, the midpoint of the side \( BC \) and the centroid of the triangle are indicated. Out of the following vector relations select the one that is not true.}
{\( \overrightarrow{ST}=  \frac12 \overrightarrow{AT} \)}{\( \overrightarrow{AT}= \frac23\overrightarrow{AS} \)}{\( \overrightarrow{ST} = -\frac13\overrightarrow{AS} \)}{\( \overrightarrow{SA}= -3\overrightarrow{TS} \)}{}{}}
\LANGpl{
\Question{
Dano trójkąt \( ABC  \). Na rysunku wskazano środek  boku \( BC \)  i środek ciężkości trójkąta. Z poniższych zależności  wybierz tą, która nie jest prawdziwa.}
{\( \overrightarrow{ST}=  \frac12 \overrightarrow{AT} \)}{\( \overrightarrow{AT}= \frac23\overrightarrow{AS} \)}{\( \overrightarrow{ST} = -\frac13\overrightarrow{AS} \)}{\( \overrightarrow{SA}= -3\overrightarrow{TS} \)}{}{}}
\LANGcs{
\Question{
V trojúhelníku \( ABC  \) na obrázku je vyznačen střed strany \( BC \) a těžiště trojúhelníka. Z následujících vztahů vyberte ten, který neplatí.}
{\( \overrightarrow{ST}=  \frac12 \overrightarrow{AT} \)}{\( \overrightarrow{AT}= \frac23\overrightarrow{AS} \)}{\( \overrightarrow{ST} = -\frac13\overrightarrow{AS} \)}{\( \overrightarrow{SA}= -3\overrightarrow{TS} \)}{}{}}
\LANGsk{
\Question{
V trojuholníku \( ABC  \) na obrázku je vyznačený stred strany \( BC \) a ťažisko trojuholníka. Z nasledujúcich vzťahov vyberte ten, ktorý neplatí.}
{\( \overrightarrow{ST}=  \frac12 \overrightarrow{AT} \)}{\( \overrightarrow{AT}= \frac23\overrightarrow{AS} \)}{\( \overrightarrow{ST} = -\frac13\overrightarrow{AS} \)}{\( \overrightarrow{SA}= -3\overrightarrow{TS} \)}{}{}}
\LANGes{
\Question{
Sea \( ABC  \) un triángulo. En la imagen se puede ver el centro del lado \( BC \) y el bicentro de este triángulo. De las siguientes relaciones vectoriales, selecciona la que no sea verdadera.}
{\( \overrightarrow{ST}=  \frac12 \overrightarrow{AT} \)}{\( \overrightarrow{AT}= \frac23\overrightarrow{AS} \)}{\( \overrightarrow{ST} = -\frac13\overrightarrow{AS} \)}{\( \overrightarrow{SA}= -3\overrightarrow{TS} \)}{}{}}
\End

\Data\ID{1103024301}
\Updated{2020-07-15 03:07:12}
\Level{A}
\SubArea{Points and vectors}
\IMG{1103024301.pdf}
\LANGen{
\Question{
In a triangle \( ABC \), let \( K \), \( L \) and \( M \) be the midpoints of \( AB \), \( BC \) and \( AC \) consecutively and let \( T \) be the centroid of \( ABC \). 
Find the values of coefficients \( k \), \( l \) and \(m \) if:
\[ \overrightarrow{TM} = k\cdot\overrightarrow{BT};\ \overrightarrow{ML} = l\cdot\overrightarrow{BA};\ \overrightarrow{CK} = m\cdot\overrightarrow{TC} \]}
{\( k=\frac12;\ l=-\frac12 ;\ m=-\frac32 \)}{\( k=\frac12;\ l=\frac12;\ m=-\frac32 \)}{\( k=\frac12 ;\ l=-\frac12 ;\ m=-\frac23 \)}{\( k=\frac12;\ l=-\frac12;\ m=\frac32 \)}{}{}}
\LANGpl{
\Question{
W trójkącie  \( ABC \),  \( K \), \( L \) i \( M \) to odpowiednio środki odcinków\( AB \), \( BC \) i \( AC \),  \( T \) to środek ciężkości trójkąta \( ABC \). 
Wskaż wartość współczynników \( k \), \( l \) i \(m \) jeśli:
\[ \overrightarrow{TM} = k\cdot\overrightarrow{BT};\ \overrightarrow{ML} = l\cdot\overrightarrow{BA};\ \overrightarrow{CK} = m\cdot\overrightarrow{TC} \]}
{\( k=\frac12;\ l=-\frac12 ;\ m=-\frac32 \)}{\( k=\frac12;\ l=\frac12;\ m=-\frac32 \)}{\( k=\frac12 ;\ l=-\frac12 ;\ m=-\frac23 \)}{\( k=\frac12;\ l=-\frac12;\ m=\frac32 \)}{}{}}
\LANGcs{
\Question{
V trojúhelníku \( ABC \) jsou body \( K \), \( L \), \( M \) postupně středy stran \( AB \), \( BC \) a \( AC \). Označme \( T \) těžiště trojúhelníka \( ABC \). 
Určete v následujících případech hodnoty koeficientů \( k \), \( l \), \(m \) tak, aby platilo:
\[ \overrightarrow{TM} = k\cdot\overrightarrow{BT};\ \overrightarrow{ML} = l\cdot\overrightarrow{BA};\ \overrightarrow{CK} = m\cdot\overrightarrow{TC} \]}
{\( k=\frac12;\ l=-\frac12 ;\ m=-\frac32 \)}{\( k=\frac12;\ l=\frac12;\ m=-\frac32 \)}{\( k=\frac12 ;\ l=-\frac12 ;\ m=-\frac23 \)}{\( k=\frac12;\ l=-\frac12;\ m=\frac32 \)}{}{}}
\LANGsk{
\Question{
V trojuholníku \( ABC \) sú body \( K \), \( L \), \( M \) postupne stredy strán \( AB \), \( BC \) a \( AC \). Označme \( T \) ťažisko trojuholníka \( ABC \). 
Určte v následujúcich prípadoch hodnoty koeficientov \( k \), \( l \), \(m \) tak, aby platilo:
\[ \overrightarrow{TM} = k\cdot\overrightarrow{BT};\ \overrightarrow{ML} = l\cdot\overrightarrow{BA};\ \overrightarrow{CK} = m\cdot\overrightarrow{TC} \]}
{\( k=\frac12;\ l=-\frac12 ;\ m=-\frac32 \)}{\( k=\frac12;\ l=\frac12;\ m=-\frac32 \)}{\( k=\frac12 ;\ l=-\frac12 ;\ m=-\frac23 \)}{\( k=\frac12;\ l=-\frac12;\ m=\frac32 \)}{}{}}
\LANGes{
\Question{
En el triángulo \( ABC \), sean \( K \), \( L \) y \( M \) los centros de\( AB \), \( BC \) y \( AC \) consecutivamente y \( T \) es el bicentro del triángulo \( ABC \). 
Determina los valores de los coeficientes \( k \), \( l \) y \(m \) si:
\[ \overrightarrow{TM} = k\cdot\overrightarrow{BT};\ \overrightarrow{ML} = l\cdot\overrightarrow{BA};\ \overrightarrow{CK} = m\cdot\overrightarrow{TC} \]}
{\( k=\frac12;\ l=-\frac12 ;\ m=-\frac32 \)}{\( k=\frac12;\ l=\frac12;\ m=-\frac32 \)}{\( k=\frac12 ;\ l=-\frac12 ;\ m=-\frac23 \)}{\( k=\frac12;\ l=-\frac12;\ m=\frac32 \)}{}{}}
\End

\Data\ID{1103024302}
\Updated{2020-07-15 03:07:12}
\Level{A}
\SubArea{Points and vectors}
\IMG{1103024302.pdf}
\LANGen{
\Question{
In a regular hexagon \( ABCDEF \) shown in the picture, let \( \overrightarrow{a} = \overrightarrow{AB} \), \( \overrightarrow{b} = \overrightarrow{BC} \), \( \overrightarrow{c} = \overrightarrow{FD} \) and \( \overrightarrow{d} = \overrightarrow{CD} \). Express vectors \( \overrightarrow{c} \) and \( \overrightarrow{d} \) as a linear combination of vectors \( \overrightarrow{a} \) and \( \overrightarrow{b} \).}
{\( \overrightarrow{c} = \overrightarrow{a} + \overrightarrow{b};\ \overrightarrow{d} = \overrightarrow{b} - \overrightarrow{a}  \)}{\( \overrightarrow{c} = 2\overrightarrow{a} + 2\overrightarrow{b};\ \overrightarrow{d} = 2\overrightarrow{b} - 0.5\overrightarrow{a}  \)}{\( \overrightarrow{c} = 2\overrightarrow{a} + \overrightarrow{b};\ \overrightarrow{d} = \overrightarrow{b} - \overrightarrow{a}  \)}{\( \overrightarrow{c} = \overrightarrow{a} + \overrightarrow{b};\ \overrightarrow{d} = \overrightarrow{a} - \overrightarrow{b}  \)}{}{}}
\LANGpl{
\Question{
Na rysunku przedstawiono prawidłowy sześciokąt \( ABCDEF \), niech \( \overrightarrow{a} = \overrightarrow{AB} \), \( \overrightarrow{b} = \overrightarrow{BC} \), \( \overrightarrow{c} = \overrightarrow{FD} \) i \( \overrightarrow{d} = \overrightarrow{CD} \). Przedstaw wektory \( \overrightarrow{c} \) i \( \overrightarrow{d} \) jako kombinację liniową wektorów \( \overrightarrow{a} \) i \( \overrightarrow{b} \).}
{\( \overrightarrow{c} = \overrightarrow{a} + \overrightarrow{b};\ \overrightarrow{d} = \overrightarrow{b} - \overrightarrow{a}  \)}{\( \overrightarrow{c} = 2\overrightarrow{a} + 2\overrightarrow{b};\ \overrightarrow{d} = 2\overrightarrow{b} - 0{,}5\overrightarrow{a}  \)}{\( \overrightarrow{c} = 2\overrightarrow{a} + \overrightarrow{b};\ \overrightarrow{d} = \overrightarrow{b} - \overrightarrow{a}  \)}{\( \overrightarrow{c} = \overrightarrow{a} + \overrightarrow{b};\ \overrightarrow{d} = \overrightarrow{a} - \overrightarrow{b}  \)}{}{}}
\LANGcs{
\Question{
V pravidelném šestiúhelníku \( ABCDEF \) na obrázku jsou vyznačeny vektory \( \overrightarrow{a} = \overrightarrow{AB} \), \( \overrightarrow{b} = \overrightarrow{BC} \), \( \overrightarrow{c} = \overrightarrow{FD} \) a \( \overrightarrow{d} = \overrightarrow{CD} \). Vyjádřete vektory \( \overrightarrow{c} \) a \( \overrightarrow{d} \) jako lineární kombinaci vektorů \( \overrightarrow{a} \) a \( \overrightarrow{b} \).}
{\( \overrightarrow{c} = \overrightarrow{a} + \overrightarrow{b};\ \overrightarrow{d} = \overrightarrow{b} - \overrightarrow{a}  \)}{\( \overrightarrow{c} = 2\overrightarrow{a} + 2\overrightarrow{b};\ \overrightarrow{d} = 2\overrightarrow{b} - 0{,}5\overrightarrow{a}  \)}{\( \overrightarrow{c} = 2\overrightarrow{a} + \overrightarrow{b};\ \overrightarrow{d} = \overrightarrow{b} - \overrightarrow{a}  \)}{\( \overrightarrow{c} = \overrightarrow{a} + \overrightarrow{b};\ \overrightarrow{d} = \overrightarrow{a} - \overrightarrow{b}  \)}{}{}}
\LANGsk{
\Question{
V pravidelnom šesťuholníku \( ABCDEF \) na obrázku sú vyznačené vektory \( \overrightarrow{a} = \overrightarrow{AB} \), \( \overrightarrow{b} = \overrightarrow{BC} \), \( \overrightarrow{c} = \overrightarrow{FD} \) a \( \overrightarrow{d} = \overrightarrow{CD} \). Vyjadrite vektory \( \overrightarrow{c} \) a \( \overrightarrow{d} \) ako lineárnu kombináciu vektorov \( \overrightarrow{a} \) a \( \overrightarrow{b} \).}
{\( \overrightarrow{c} = \overrightarrow{a} + \overrightarrow{b};\ \overrightarrow{d} = \overrightarrow{b} - \overrightarrow{a}  \)}{\( \overrightarrow{c} = 2\overrightarrow{a} + 2\overrightarrow{b};\ \overrightarrow{d} = 2\overrightarrow{b} - 0{,}5\overrightarrow{a}  \)}{\( \overrightarrow{c} = 2\overrightarrow{a} + \overrightarrow{b};\ \overrightarrow{d} = \overrightarrow{b} - \overrightarrow{a}  \)}{\( \overrightarrow{c} = \overrightarrow{a} + \overrightarrow{b};\ \overrightarrow{d} = \overrightarrow{a} - \overrightarrow{b}  \)}{}{}}
\LANGes{
\Question{
En un hexágono regular  \( ABCDEF \) que es mostrado en la imagen, sean \( \overrightarrow{a} = \overrightarrow{AB} \), \( \overrightarrow{b} = \overrightarrow{BC} \), \( \overrightarrow{c} = \overrightarrow{FD} \) y \( \overrightarrow{d} = \overrightarrow{CD} \). Expresa vectores \( \overrightarrow{c} \) y \( \overrightarrow{d} \) como una combinación lineal de vectores\( \overrightarrow{a} \) y \( \overrightarrow{b} \).}
{\( \overrightarrow{c} = \overrightarrow{a} + \overrightarrow{b};\ \overrightarrow{d} = \overrightarrow{b} - \overrightarrow{a}  \)}{\( \overrightarrow{c} = 2\overrightarrow{a} + 2\overrightarrow{b};\ \overrightarrow{d} = 2\overrightarrow{b} - 0.5\overrightarrow{a}  \)}{\( \overrightarrow{c} = 2\overrightarrow{a} + \overrightarrow{b};\ \overrightarrow{d} = \overrightarrow{b} - \overrightarrow{a}  \)}{\( \overrightarrow{c} = \overrightarrow{a} + \overrightarrow{b};\ \overrightarrow{d} = \overrightarrow{a} - \overrightarrow{b}  \)}{}{}}
\End

\Data\ID{1103024303}
\Updated{2020-07-15 03:07:12}
\Level{A}
\SubArea{Points and vectors}
\IMG{1103024303.pdf}
\LANGen{
\Question{
The picture shows a rectangular cuboid \( ABCDEFGH \) with \( \overrightarrow{a} = \overrightarrow{AB} \), \( \overrightarrow{b} = \overrightarrow{AD} \), \( \overrightarrow{c} = \overrightarrow{AE} \), \( \overrightarrow{x} = \overrightarrow{AK} \) and \( \overrightarrow{y} = \overrightarrow{AL} \). Point \( K \) is the midpoint of \( FG \) and point \( L \) is the centre of face \( BCGF \).  Express vectors \( \overrightarrow{x} \) and \( \overrightarrow{y} \) as a linear combination of vectors \( \overrightarrow{a} \),  \( \overrightarrow{b} \),  \( \overrightarrow{c} \).}
{\(  \overrightarrow{x} =  \overrightarrow{a} +  \frac12\overrightarrow{b} + \overrightarrow{c};\  \overrightarrow{y} =  \overrightarrow{a} + \frac12\overrightarrow{b} + \frac12\overrightarrow{c} \)}{\(  \overrightarrow{x} =  \frac12\overrightarrow{a} +  \overrightarrow{b} + \frac12\overrightarrow{c};\  \overrightarrow{y} = \overrightarrow{a} - \frac12\overrightarrow{b} + \frac12\overrightarrow{c} \)}{\(  \overrightarrow{x} =  \overrightarrow{a} +  \frac12\overrightarrow{b} + \frac12\overrightarrow{c};\  \overrightarrow{y} =   \overrightarrow{a} - \frac12\overrightarrow{b} + \frac12\overrightarrow{c} \)}{\(  \overrightarrow{x} =  \overrightarrow{a} + \frac12\overrightarrow{b} + \frac12\overrightarrow{c};\  \overrightarrow{y} =  \frac12\overrightarrow{a} + \frac12\overrightarrow{b} + \frac12\overrightarrow{c} \)}{}{}}
\LANGpl{
\Question{
Na rysunku przedstawiono prostopadłościan \( ABCDEFGH \), gdzie \( \overrightarrow{a} = \overrightarrow{AB} \), \( \overrightarrow{b} = \overrightarrow{AD} \), \( \overrightarrow{c} = \overrightarrow{AE} \), \( \overrightarrow{x} = \overrightarrow{AK} \) i \( \overrightarrow{y} = \overrightarrow{AL} \). Punkt \( K \) to środek \( FG \), punkt \( L \) to środek ściany  \( BCGF \).  Przedstaw wektory \( \overrightarrow{x} \) i \( \overrightarrow{y} \) jako kombinację liniową wektorów \( \overrightarrow{a} \),  \( \overrightarrow{b} \),  \( \overrightarrow{c} \).}
{\(  \overrightarrow{x} =  \overrightarrow{a} +  \frac12\overrightarrow{b} + \overrightarrow{c};\  \overrightarrow{y} =  \overrightarrow{a} + \frac12\overrightarrow{b} + \frac12\overrightarrow{c} \)}{\(  \overrightarrow{x} =  \frac12\overrightarrow{a} +  \overrightarrow{b} + \frac12\overrightarrow{c};\  \overrightarrow{y} = \overrightarrow{a} - \frac12\overrightarrow{b} + \frac12\overrightarrow{c} \)}{\(  \overrightarrow{x} =  \overrightarrow{a} +  \frac12\overrightarrow{b} + \frac12\overrightarrow{c};\  \overrightarrow{y} =   \overrightarrow{a} - \frac12\overrightarrow{b} + \frac12\overrightarrow{c} \)}{\(  \overrightarrow{x} =  \overrightarrow{a} + \frac12\overrightarrow{b} + \frac12\overrightarrow{c};\  \overrightarrow{y} =  \frac12\overrightarrow{a} + \frac12\overrightarrow{b} + \frac12\overrightarrow{c} \)}{}{}}
\LANGcs{
\Question{
V kvádru \( ABCDEFGH \) na obrázku jsou vyznačeny vektory \( \overrightarrow{a} = \overrightarrow{AB} \), \( \overrightarrow{b} = \overrightarrow{AD} \), \( \overrightarrow{c} = \overrightarrow{AE} \), \( \overrightarrow{x} = \overrightarrow{AK} \) a \( \overrightarrow{y} = \overrightarrow{AL} \). Bod \( K \) je středem hrany \( FG \) a bod \( L \) je středem stěny \( BCGF \).  Vyjádřete vektory \( \overrightarrow{x} \) a \( \overrightarrow{y} \) jako lineární kombinaci vektorů \( \overrightarrow{a} \),  \( \overrightarrow{b} \),  \( \overrightarrow{c} \).}
{\(  \overrightarrow{x} =  \overrightarrow{a} +  \frac12\overrightarrow{b} + \overrightarrow{c};\  \overrightarrow{y} =  \overrightarrow{a} + \frac12\overrightarrow{b} + \frac12\overrightarrow{c} \)}{\(  \overrightarrow{x} =  \frac12\overrightarrow{a} +  \overrightarrow{b} + \frac12\overrightarrow{c};\  \overrightarrow{y} = \overrightarrow{a} - \frac12\overrightarrow{b} + \frac12\overrightarrow{c} \)}{\(  \overrightarrow{x} =  \overrightarrow{a} +  \frac12\overrightarrow{b} + \frac12\overrightarrow{c};\  \overrightarrow{y} =   \overrightarrow{a} - \frac12\overrightarrow{b} + \frac12\overrightarrow{c} \)}{\(  \overrightarrow{x} =  \overrightarrow{a} + \frac12\overrightarrow{b} + \frac12\overrightarrow{c};\  \overrightarrow{y} =  \frac12\overrightarrow{a} + \frac12\overrightarrow{b} + \frac12\overrightarrow{c} \)}{}{}}
\LANGsk{
\Question{
V kvádri \( ABCDEFGH \) na obrázku sú vyznačené vektory \( \overrightarrow{a} = \overrightarrow{AB} \), \( \overrightarrow{b} = \overrightarrow{AD} \), \( \overrightarrow{c} = \overrightarrow{AE} \), \( \overrightarrow{x} = \overrightarrow{AK} \) a \( \overrightarrow{y} = \overrightarrow{AL} \). Bod \( K \) je stredom hrany \( FG \) a bod \( L \) je stredom steny \( BCGF \).  Vyjadrite vektory \( \overrightarrow{x} \) a \( \overrightarrow{y} \) ako lineárnu kombináciu vektorov \( \overrightarrow{a} \),  \( \overrightarrow{b} \),  \( \overrightarrow{c} \).}
{\(  \overrightarrow{x} =  \overrightarrow{a} +  \frac12\overrightarrow{b} + \overrightarrow{c};\  \overrightarrow{y} =  \overrightarrow{a} + \frac12\overrightarrow{b} + \frac12\overrightarrow{c} \)}{\(  \overrightarrow{x} =  \frac12\overrightarrow{a} +  \overrightarrow{b} + \frac12\overrightarrow{c};\  \overrightarrow{y} = \overrightarrow{a} - \frac12\overrightarrow{b} + \frac12\overrightarrow{c} \)}{\(  \overrightarrow{x} =  \overrightarrow{a} +  \frac12\overrightarrow{b} + \frac12\overrightarrow{c};\  \overrightarrow{y} =   \overrightarrow{a} - \frac12\overrightarrow{b} + \frac12\overrightarrow{c} \)}{\(  \overrightarrow{x} =  \overrightarrow{a} + \frac12\overrightarrow{b} + \frac12\overrightarrow{c};\  \overrightarrow{y} =  \frac12\overrightarrow{a} + \frac12\overrightarrow{b} + \frac12\overrightarrow{c} \)}{}{}}
\LANGes{
\Question{
La imagen muestra un ortoedro \( ABCDEFGH \) con \( \overrightarrow{a} = \overrightarrow{AB} \), \( \overrightarrow{b} = \overrightarrow{AD} \), \( \overrightarrow{c} = \overrightarrow{AE} \), \( \overrightarrow{x} = \overrightarrow{AK} \) y \( \overrightarrow{y} = \overrightarrow{AL} \). El punto \( K \) es el centro de \( FG \) y el punto \( L \) es el centro de la cara \( BCGF \).  Expresa vectores \( \overrightarrow{x} \) y \( \overrightarrow{y} \) como una combinación lineal de vectores \( \overrightarrow{a} \),  \( \overrightarrow{b} \),  \( \overrightarrow{c} \).}
{\(  \overrightarrow{x} =  \overrightarrow{a} +  \frac12\overrightarrow{b} + \overrightarrow{c};\  \overrightarrow{y} =  \overrightarrow{a} + \frac12\overrightarrow{b} + \frac12\overrightarrow{c} \)}{\(  \overrightarrow{x} =  \frac12\overrightarrow{a} +  \overrightarrow{b} + \frac12\overrightarrow{c};\  \overrightarrow{y} = \overrightarrow{a} - \frac12\overrightarrow{b} + \frac12\overrightarrow{c} \)}{\(  \overrightarrow{x} =  \overrightarrow{a} +  \frac12\overrightarrow{b} + \frac12\overrightarrow{c};\  \overrightarrow{y} =   \overrightarrow{a} - \frac12\overrightarrow{b} + \frac12\overrightarrow{c} \)}{\(  \overrightarrow{x} =  \overrightarrow{a} + \frac12\overrightarrow{b} + \frac12\overrightarrow{c};\  \overrightarrow{y} =  \frac12\overrightarrow{a} + \frac12\overrightarrow{b} + \frac12\overrightarrow{c} \)}{}{}}
\End

\Data\ID{1103024304}
\Updated{2020-07-15 03:07:12}
\Level{A}
\SubArea{Points and vectors}
\IMG{1103024304_0.pdf}
\LANGen{
\Question{
The picture shows a rectangular cuboid  \( ABCDEFGH \).  In the cuboid find the vector that is the sum of \( \overrightarrow{BC} + \overrightarrow{AE} + \overrightarrow{CF} + \overrightarrow{FA} + \overrightarrow{HG} \).}
{\( \overrightarrow{BF} \)}{\( \overrightarrow{BE} \)}{\( \overrightarrow{BG} \)}{\( \overrightarrow{BH} \)}{}{}}
\LANGpl{
\Question{
Na rysunku przedstawiono prostopadłościan  \( ABCDEFGH \).  Wskaż wektor w prostopadłościanie, który jest sumą \( \overrightarrow{BC} + \overrightarrow{AE} + \overrightarrow{CF} + \overrightarrow{FA} + \overrightarrow{HG} \).}
{\( \overrightarrow{BF} \)}{\( \overrightarrow{BE} \)}{\( \overrightarrow{BG} \)}{\( \overrightarrow{BH} \)}{}{}}
\LANGcs{
\Question{
V kvádru  \( ABCDEFGH \) s vyznačenými vektory určete součet \( \overrightarrow{BC} + \overrightarrow{AE} + \overrightarrow{CF} + \overrightarrow{FA} + \overrightarrow{HG} \).}
{\( \overrightarrow{BF} \)}{\( \overrightarrow{BE} \)}{\( \overrightarrow{BG} \)}{\( \overrightarrow{BH} \)}{}{}}
\LANGsk{
\Question{
V kvádri  \( ABCDEFGH \) s vyznačenými vektormi určte súčet \( \overrightarrow{BC} + \overrightarrow{AE} + \overrightarrow{CF} + \overrightarrow{FA} + \overrightarrow{HG} \).}
{\( \overrightarrow{BF} \)}{\( \overrightarrow{BE} \)}{\( \overrightarrow{BG} \)}{\( \overrightarrow{BH} \)}{}{}}
\LANGes{
\Question{
La imagen muestra un ortoedro \( ABCDEFGH \). En el ortoedro, determina el vector que es la suma de  \( \overrightarrow{BC} + \overrightarrow{AE} + \overrightarrow{CF} + \overrightarrow{FA} + \overrightarrow{HG} \).}
{\( \overrightarrow{BF} \)}{\( \overrightarrow{BE} \)}{\( \overrightarrow{BG} \)}{\( \overrightarrow{BH} \)}{}{}}
\End

\Data\ID{1103024305}
\Updated{2020-07-15 03:07:12}
\Level{A}
\SubArea{Points and vectors}
\IMG{1103024305.pdf}
\LANGen{
\Question{
In a tetrahedron \( ABCD \), let \( \overrightarrow{b} = \overrightarrow{AB} \), \( \overrightarrow{c} = \overrightarrow{AC} \), \( \overrightarrow{d} = \overrightarrow{AD} \), \( \overrightarrow{e} = \overrightarrow{AE} \) and \( \overrightarrow{f} = \overrightarrow{DE} \). Further let \( E \) be the midpoint of \( BC \). Express vectors \( \overrightarrow{e} \) and \( \overrightarrow{f} \) as a linear combination of vectors \( \overrightarrow{b} \), \( \overrightarrow{c} \), \( \overrightarrow{d} \).}
{\( \overrightarrow{e} = \frac12\overrightarrow{b} + \frac12\overrightarrow{c};\ \overrightarrow{f} = \frac12\overrightarrow{b} + \frac12\overrightarrow{c} - \overrightarrow{d} \)}{\( \overrightarrow{e} = \frac12\overrightarrow{b} + \frac12\overrightarrow{d};\ \overrightarrow{f} = \overrightarrow{b} + \overrightarrow{c} + \overrightarrow{d} \)}{\( \overrightarrow{e} = \overrightarrow{b} + \overrightarrow{c};\ \overrightarrow{f} =\frac12\overrightarrow{b} + \frac12\overrightarrow{c} - \overrightarrow{d} \)}{\( \overrightarrow{e} = \frac12\overrightarrow{b} + \frac12\overrightarrow{c};\ \overrightarrow{f} = \frac12\overrightarrow{b} + \frac12\overrightarrow{c} + \overrightarrow{d} \)}{}{}}
\LANGpl{
\Question{
W czworościanie \( ABCD \), gdzie \( \overrightarrow{b} = \overrightarrow{AB} \), \( \overrightarrow{c} = \overrightarrow{AC} \), \( \overrightarrow{d} = \overrightarrow{AD} \), \( \overrightarrow{e} = \overrightarrow{AE} \) i \( \overrightarrow{f} = \overrightarrow{DE} \).  \( E \) to środek \( BC \). Przestaw wektory \( \overrightarrow{e} \) i \( \overrightarrow{f} \) jako kombinację liniową wektorów \( \overrightarrow{b} \), \( \overrightarrow{c} \), \( \overrightarrow{d} \).}
{\( \overrightarrow{e} = \frac12\overrightarrow{b} + \frac12\overrightarrow{c};\ \overrightarrow{f} = \frac12\overrightarrow{b} + \frac12\overrightarrow{c} - \overrightarrow{d} \)}{\( \overrightarrow{e} = \frac12\overrightarrow{b} + \frac12\overrightarrow{d};\ \overrightarrow{f} = \overrightarrow{b} + \overrightarrow{c} + \overrightarrow{d} \)}{\( \overrightarrow{e} = \overrightarrow{b} + \overrightarrow{c};\ \overrightarrow{f} =\frac12\overrightarrow{b} + \frac12\overrightarrow{c} - \overrightarrow{d} \)}{\( \overrightarrow{e} = \frac12\overrightarrow{b} + \frac12\overrightarrow{c};\ \overrightarrow{f} = \frac12\overrightarrow{b} + \frac12\overrightarrow{c} + \overrightarrow{d} \)}{}{}}
\LANGcs{
\Question{
Ve čtyřstěnu \( ABCD \) jsou vyznačeny vektory \( \overrightarrow{b} = \overrightarrow{AB} \), \( \overrightarrow{c} = \overrightarrow{AC} \), \( \overrightarrow{d} = \overrightarrow{AD} \), \( \overrightarrow{e} = \overrightarrow{AE} \) a \( \overrightarrow{f} = \overrightarrow{DE} \), kde \( E \) je střed hrany \( BC \). Vyjádřete vektory \( \overrightarrow{e} \) a \( \overrightarrow{f} \) jako lineární kombinaci vektorů \( \overrightarrow{b} \), \( \overrightarrow{c} \), \( \overrightarrow{d} \).}
{\( \overrightarrow{e} = \frac12\overrightarrow{b} + \frac12\overrightarrow{c};\ \overrightarrow{f} = \frac12\overrightarrow{b} + \frac12\overrightarrow{c} - \overrightarrow{d} \)}{\( \overrightarrow{e} = \frac12\overrightarrow{b} + \frac12\overrightarrow{d};\ \overrightarrow{f} = \overrightarrow{b} + \overrightarrow{c} + \overrightarrow{d} \)}{\( \overrightarrow{e} = \overrightarrow{b} + \overrightarrow{c};\ \overrightarrow{f} =\frac12\overrightarrow{b} + \frac12\overrightarrow{c} - \overrightarrow{d} \)}{\( \overrightarrow{e} = \frac12\overrightarrow{b} + \frac12\overrightarrow{c};\ \overrightarrow{f} = \frac12\overrightarrow{b} + \frac12\overrightarrow{c} + \overrightarrow{d} \)}{}{}}
\LANGsk{
\Question{
V štvorstene \( ABCD \) sú vyznačené vektory \( \overrightarrow{b} = \overrightarrow{AB} \), \( \overrightarrow{c} = \overrightarrow{AC} \), \( \overrightarrow{d} = \overrightarrow{AD} \), \( \overrightarrow{e} = \overrightarrow{AE} \) a \( \overrightarrow{f} = \overrightarrow{DE} \), kde \( E \) je stred hrany \( BC \). Vyjadrite vektory \( \overrightarrow{e} \) a \( \overrightarrow{f} \) ako lineárnu kombináciu vektorov \( \overrightarrow{b} \), \( \overrightarrow{c} \), \( \overrightarrow{d} \).}
{\( \overrightarrow{e} = \frac12\overrightarrow{b} + \frac12\overrightarrow{c};\ \overrightarrow{f} = \frac12\overrightarrow{b} + \frac12\overrightarrow{c} - \overrightarrow{d} \)}{\( \overrightarrow{e} = \frac12\overrightarrow{b} + \frac12\overrightarrow{d};\ \overrightarrow{f} = \overrightarrow{b} + \overrightarrow{c} + \overrightarrow{d} \)}{\( \overrightarrow{e} = \overrightarrow{b} + \overrightarrow{c};\ \overrightarrow{f} =\frac12\overrightarrow{b} + \frac12\overrightarrow{c} - \overrightarrow{d} \)}{\( \overrightarrow{e} = \frac12\overrightarrow{b} + \frac12\overrightarrow{c};\ \overrightarrow{f} = \frac12\overrightarrow{b} + \frac12\overrightarrow{c} + \overrightarrow{d} \)}{}{}}
\LANGes{
\Question{
En un tetraedro  \( ABCD \), sean \( \overrightarrow{b} = \overrightarrow{AB} \), \( \overrightarrow{c} = \overrightarrow{AC} \), \( \overrightarrow{d} = \overrightarrow{AD} \), \( \overrightarrow{e} = \overrightarrow{AE} \) y \( \overrightarrow{f} = \overrightarrow{DE} \). Además, sea \( E \) el centro de  \( BC \). Expresa vectores \( \overrightarrow{e} \) y \( \overrightarrow{f} \) como una combinación lineal de vectores \( \overrightarrow{b} \), \( \overrightarrow{c} \), \( \overrightarrow{d} \).}
{\( \overrightarrow{e} = \frac12\overrightarrow{b} + \frac12\overrightarrow{c};\ \overrightarrow{f} = \frac12\overrightarrow{b} + \frac12\overrightarrow{c} - \overrightarrow{d} \)}{\( \overrightarrow{e} = \frac12\overrightarrow{b} + \frac12\overrightarrow{d};\ \overrightarrow{f} = \overrightarrow{b} + \overrightarrow{c} + \overrightarrow{d} \)}{\( \overrightarrow{e} = \overrightarrow{b} + \overrightarrow{c};\ \overrightarrow{f} =\frac12\overrightarrow{b} + \frac12\overrightarrow{c} - \overrightarrow{d} \)}{\( \overrightarrow{e} = \frac12\overrightarrow{b} + \frac12\overrightarrow{c};\ \overrightarrow{f} = \frac12\overrightarrow{b} + \frac12\overrightarrow{c} + \overrightarrow{d} \)}{}{}}
\End

\Data\ID{1103024308}
\Updated{2020-07-15 03:07:12}
\Level{A}
\SubArea{Points and vectors}
\IMG{1103024308.pdf}
\LANGen{
\Question{
Given the vectors \( \overrightarrow{a} \), \( \overrightarrow{b} \), \( \overrightarrow{c} \) shown in the picture, express a vector \( \overrightarrow{c} \) as a linear combination of vectors \( \overrightarrow{a} \) and \( \overrightarrow{b} \).}
{\( \overrightarrow{c} = -2\overrightarrow{a} + \overrightarrow{b}  \)}{\( \overrightarrow{c} = -\overrightarrow{a} + \frac12\overrightarrow{b}  \)}{\( \overrightarrow{c} = -\frac32\overrightarrow{a} + \overrightarrow{b}  \)}{\( \overrightarrow{c} = -2\overrightarrow{a} + \frac32\overrightarrow{b}  \)}{}{}}
\LANGpl{
\Question{
Na rysunku wskazano wektory \( \overrightarrow{a} \), \( \overrightarrow{b} \), \( \overrightarrow{c} \), przedstaw wektor \( \overrightarrow{c} \) jako kombinację liniową wektorów \( \overrightarrow{a} \) i  \( \overrightarrow{b} \).}
{\( \overrightarrow{c} = -2\overrightarrow{a} + \overrightarrow{b}  \)}{\( \overrightarrow{c} = -\overrightarrow{a} + \frac12\overrightarrow{b}  \)}{\( \overrightarrow{c} = -\frac32\overrightarrow{a} + \overrightarrow{b}  \)}{\( \overrightarrow{c} = -2\overrightarrow{a} + \frac32\overrightarrow{b}  \)}{}{}}
\LANGcs{
\Question{
V obrázku jsou dány vektory \( \overrightarrow{a} \), \( \overrightarrow{b} \), \( \overrightarrow{c} \). Vyjádřete vektor \( \overrightarrow{c} \) jako lineární kombinaci vektorů \( \overrightarrow{a} \) a \( \overrightarrow{b} \).}
{\( \overrightarrow{c} = -2\overrightarrow{a} + \overrightarrow{b}  \)}{\( \overrightarrow{c} = -\overrightarrow{a} + \frac12\overrightarrow{b}  \)}{\( \overrightarrow{c} = -\frac32\overrightarrow{a} + \overrightarrow{b}  \)}{\( \overrightarrow{c} = -2\overrightarrow{a} + \frac32\overrightarrow{b}  \)}{}{}}
\LANGsk{
\Question{
Na obrázku sú dané vektory \( \overrightarrow{a} \), \( \overrightarrow{b} \), \( \overrightarrow{c} \). Vyjadrite vektor \( \overrightarrow{c} \) ako lineárnu kombináciu vektorov \( \overrightarrow{a} \) a \( \overrightarrow{b} \).}
{\( \overrightarrow{c} = -2\overrightarrow{a} + \overrightarrow{b}  \)}{\( \overrightarrow{c} = -\overrightarrow{a} + \frac12\overrightarrow{b}  \)}{\( \overrightarrow{c} = -\frac32\overrightarrow{a} + \overrightarrow{b}  \)}{\( \overrightarrow{c} = -2\overrightarrow{a} + \frac32\overrightarrow{b}  \)}{}{}}
\LANGes{
\Question{
Dados los vectores \( \overrightarrow{a} \), \( \overrightarrow{b} \), \( \overrightarrow{c} \) mostrados en la imagen, expressa el vector \( \overrightarrow{c} \) como una combinación lineal de vectores \( \overrightarrow{a} \) y \( \overrightarrow{b} \).}
{\( \overrightarrow{c} = -2\overrightarrow{a} + \overrightarrow{b}  \)}{\( \overrightarrow{c} = -\overrightarrow{a} + \frac12\overrightarrow{b}  \)}{\( \overrightarrow{c} = -\frac32\overrightarrow{a} + \overrightarrow{b}  \)}{\( \overrightarrow{c} = -2\overrightarrow{a} + \frac32\overrightarrow{b}  \)}{}{}}
\End

\Data\ID{1103024309}
\Updated{2020-07-15 03:07:12}
\Level{A}
\SubArea{Points and vectors}
\IMG{1103024309.pdf}
\LANGen{
\Question{
Given the vectors \( \overrightarrow{a} \), \( \overrightarrow{b} \), \( \overrightarrow{c} \) shown in the picture, express a vector \( \overrightarrow{b} \) as a linear combination of vectors \( \overrightarrow{a} \) and \( \overrightarrow{c} \).}
{\( \overrightarrow{b} = 2\overrightarrow{a} + \overrightarrow{c}  \)}{\( \overrightarrow{b} = 2\overrightarrow{a} - \overrightarrow{c}  \)}{\( \overrightarrow{b} = -2\overrightarrow{a} + \overrightarrow{c}  \)}{\( \overrightarrow{b} = -2\overrightarrow{a} - \overrightarrow{c}  \)}{}{}}
\LANGpl{
\Question{
Na rysunku wskazano wektory \( \overrightarrow{a} \), \( \overrightarrow{b} \), \( \overrightarrow{c} \), przedstaw wektor  \( \overrightarrow{b} \) jako kombinację liniową wektorów \( \overrightarrow{a} \) i \( \overrightarrow{c} \).}
{\( \overrightarrow{b} = 2\overrightarrow{a} + \overrightarrow{c}  \)}{\( \overrightarrow{b} = 2\overrightarrow{a} - \overrightarrow{c}  \)}{\( \overrightarrow{b} = -2\overrightarrow{a} + \overrightarrow{c}  \)}{\( \overrightarrow{b} = -2\overrightarrow{a} - \overrightarrow{c}  \)}{}{}}
\LANGcs{
\Question{
V obrázku jsou dány vektory \( \overrightarrow{a} \), \( \overrightarrow{b} \), \( \overrightarrow{c} \). Vyjádřete vektor \( \overrightarrow{b} \) jako lineární kombinaci vektorů \( \overrightarrow{a} \) a \( \overrightarrow{c} \).}
{\( \overrightarrow{b} = 2\overrightarrow{a} + \overrightarrow{c}  \)}{\( \overrightarrow{b} = 2\overrightarrow{a} - \overrightarrow{c}  \)}{\( \overrightarrow{b} = -2\overrightarrow{a} + \overrightarrow{c}  \)}{\( \overrightarrow{b} = -2\overrightarrow{a} - \overrightarrow{c}  \)}{}{}}
\LANGsk{
\Question{
V obrázku sú dané vektory \( \overrightarrow{a} \), \( \overrightarrow{b} \), \( \overrightarrow{c} \). Vyjadrite vektor \( \overrightarrow{b} \) ako lineárnu kombináciu vektorov \( \overrightarrow{a} \) a \( \overrightarrow{c} \).}
{\( \overrightarrow{b} = 2\overrightarrow{a} + \overrightarrow{c}  \)}{\( \overrightarrow{b} = 2\overrightarrow{a} - \overrightarrow{c}  \)}{\( \overrightarrow{b} = -2\overrightarrow{a} + \overrightarrow{c}  \)}{\( \overrightarrow{b} = -2\overrightarrow{a} - \overrightarrow{c}  \)}{}{}}
\LANGes{
\Question{
Dados los vectores \( \overrightarrow{a} \), \( \overrightarrow{b} \), \( \overrightarrow{c} \) mostrados en la imagen, expresa el vector  \( \overrightarrow{b} \) como una combinación lineal de vectores  \( \overrightarrow{a} \) y \( \overrightarrow{c} \).}
{\( \overrightarrow{b} = 2\overrightarrow{a} + \overrightarrow{c}  \)}{\( \overrightarrow{b} = 2\overrightarrow{a} - \overrightarrow{c}  \)}{\( \overrightarrow{b} = -2\overrightarrow{a} + \overrightarrow{c}  \)}{\( \overrightarrow{b} = -2\overrightarrow{a} - \overrightarrow{c}  \)}{}{}}
\End

\Data\ID{1103024310}
\Updated{2020-07-15 03:07:12}
\Level{A}
\SubArea{Points and vectors}
\IMG{1103024310.pdf}
\LANGen{
\Question{
The picture shows the triangle \( KLM \) with indicated vectors  \( \overrightarrow{a} \), \( \overrightarrow{b} \), \( \overrightarrow{c} \) in a coordinate system. What are the vector coordinates \( \overrightarrow{b} \)? Express \( \overrightarrow{b} \) as a linear combination of \( \overrightarrow{a} \) and \( \overrightarrow{c} \).}
{\( \overrightarrow{b} = \left(1;3;4.5\right);\ \overrightarrow{b} = \frac12\overrightarrow{a} + \frac12\overrightarrow{c}  \)}{\( \overrightarrow{b} = \left(3;1;4.5\right);\ \overrightarrow{b} = \overrightarrow{a} + \overrightarrow{c}  \)}{\( \overrightarrow{b} = \left(1;3;4.5\right);\ \overrightarrow{b} = \overrightarrow{a} + \overrightarrow{c}  \)}{\( \overrightarrow{b} = \left(3;1;4.5\right);\ \overrightarrow{b} = \frac12\overrightarrow{a} + \frac12\overrightarrow{c}  \)}{}{}}
\LANGpl{
\Question{
W układzie współrzędnych podano trójkąt \( KLM \) o wektorach  \( \overrightarrow{a} \), \( \overrightarrow{b} \), \( \overrightarrow{c} \). Określ współrzędne wektora \( \overrightarrow{b} \). Przedstaw\( \overrightarrow{b} \) jako kombinację liniową  wektorów \( \overrightarrow{a} \) i  \( \overrightarrow{c} \).}
{\( \overrightarrow{b} = \left(1;3;4{,}5\right);\ \overrightarrow{b} = \frac12\overrightarrow{a} + \frac12\overrightarrow{c}  \)}{\( \overrightarrow{b} = \left(3;1;4{,}5\right);\ \overrightarrow{b} = \overrightarrow{a} + \overrightarrow{c}  \)}{\( \overrightarrow{b} = \left(1;3;4{,}5\right);\ \overrightarrow{b} = \overrightarrow{a} + \overrightarrow{c}  \)}{\( \overrightarrow{b} = \left(3;1;4{,}5\right);\ \overrightarrow{b} = \frac12\overrightarrow{a} + \frac12\overrightarrow{c}  \)}{}{}}
\LANGcs{
\Question{
V souřadném systému je dán trojúhelník \( KLM \) s vyznačenými vektory \( \overrightarrow{a} \), \( \overrightarrow{b} \), \( \overrightarrow{c} \). Určete souřadnice vektoru \( \overrightarrow{b} \) a vyjádřete jej jako lineární kombinaci vektorů \( \overrightarrow{a} \) a \( \overrightarrow{c} \).}
{\( \overrightarrow{b} = \left(1;3;4{,}5\right);\ \overrightarrow{b} = \frac12\overrightarrow{a} + \frac12\overrightarrow{c}  \)}{\( \overrightarrow{b} = \left(3;1;4{,}5\right);\ \overrightarrow{b} = \overrightarrow{a} + \overrightarrow{c}  \)}{\( \overrightarrow{b} = \left(1;3;4{,}5\right);\ \overrightarrow{b} = \overrightarrow{a} + \overrightarrow{c}  \)}{\( \overrightarrow{b} = \left(3;1;4{,}5\right);\ \overrightarrow{b} = \frac12\overrightarrow{a} + \frac12\overrightarrow{c}  \)}{}{}}
\LANGsk{
\Question{
V súradnicovom systéme je daný trojuholník \( KLM \) s vyznačenými vektormi \( \overrightarrow{a} \), \( \overrightarrow{b} \), \( \overrightarrow{c} \). Určte súradnice vektora \( \overrightarrow{b} \) a vyjadrite ich ako lineárnu kombináciu vektorov \( \overrightarrow{a} \) a \( \overrightarrow{c} \).}
{\( \overrightarrow{b} = \left(1;3;4{,}5\right);\ \overrightarrow{b} = \frac12\overrightarrow{a} + \frac12\overrightarrow{c}  \)}{\( \overrightarrow{b} = \left(3;1;4{,}5\right);\ \overrightarrow{b} = \overrightarrow{a} + \overrightarrow{c}  \)}{\( \overrightarrow{b} = \left(1;3;4{,}5\right);\ \overrightarrow{b} = \overrightarrow{a} + \overrightarrow{c}  \)}{\( \overrightarrow{b} = \left(3;1;4{,}5\right);\ \overrightarrow{b} = \frac12\overrightarrow{a} + \frac12\overrightarrow{c}  \)}{}{}}
\LANGes{
\Question{
La imagen muestra el triángulo \( KLM \) con los vectores indicados  \( \overrightarrow{a} \), \( \overrightarrow{b} \), \( \overrightarrow{c} \) en un sistema de coordenadas. ¿Cuáles son las coordenadas de un vector\( \overrightarrow{b} \)? Expresa \( \overrightarrow{b} \) como una combinación lineal de vectores \( \overrightarrow{a} \) y \( \overrightarrow{c} \).}
{\( \overrightarrow{b} = \left(1;3;4.5\right);\ \overrightarrow{b} = \frac12\overrightarrow{a} + \frac12\overrightarrow{c}  \)}{\( \overrightarrow{b} = \left(3;1;4.5\right);\ \overrightarrow{b} = \overrightarrow{a} + \overrightarrow{c}  \)}{\( \overrightarrow{b} = \left(1;3;4.5\right);\ \overrightarrow{b} = \overrightarrow{a} + \overrightarrow{c}  \)}{\( \overrightarrow{b} = \left(3;1;4.5\right);\ \overrightarrow{b} = \frac12\overrightarrow{a} + \frac12\overrightarrow{c}  \)}{}{}}
\End

\Data\ID{1103030701}
\Updated{2020-07-15 03:07:12}
\Level{A}
\SubArea{Points and vectors}
\IMG{1103030701.pdf}
\LANGen{
\Question{
We are given points \( A = [1;-1;2] \), \( B = [0;5;-3] \), \( S = [2;0;5] \). Point \( S \) is the centre of a parallelogram \( ABCD \). Find the coordinates of vertices \( C \) and \( D \).}
{\( C = [3;1;8]; D = [4;-5;13] \)}{\( C = [4;-5;13]; D = [3;1;8] \)}{\( C = [1;1;3]; D = [2;-5;8] \)}{\( C = [-3;-1;-8]; D = [-4;5;-13] \)}{}{}}
\LANGpl{
\Question{
Dane są punkty \( A = [1;-1;2] \), \( B = [0;5;-3] \), \( S = [2;0;5] \). Punkt \( S \) to środek równoległoboku \( ABCD \). Wyznacz współrzędne wierzchołków \( C \) i \( D \).}
{\( C = [3;1;8]; D = [4;-5;13] \)}{\( C = [4;-5;13]; D = [3;1;8] \)}{\( C = [1;1;3]; D = [2;-5;8] \)}{\( C = [-3;-1;-8]; D = [-4;5;-13] \)}{}{}}
\LANGcs{
\Question{
Jsou dány body \( A = [1;-1;2] \), \( B = [0;5;-3] \), \( S = [2;0;5] \). Bod \( S \) je středem rovnoběžníku \( ABCD \). Určete souřadnice vrcholů \( C \) a \( D \).}
{\( C = [3;1;8]; D = [4;-5;13] \)}{\( C = [4;-5;13]; D = [3;1;8] \)}{\( C = [1;1;3]; D = [2;-5;8] \)}{\( C = [-3;-1;-8]; D = [-4;5;-13] \)}{}{}}
\LANGsk{
\Question{
Sú dané body \( A = [1;-1;2] \), \( B = [0;5;-3] \), \( S = [2;0;5] \). Bod \( S \) je stredom rovnobežníka \( ABCD \). Určte súradnice vrcholov \( C \) a \( D \).}
{\( C = [3;1;8]; D = [4;-5;13] \)}{\( C = [4;-5;13]; D = [3;1;8] \)}{\( C = [1;1;3]; D = [2;-5;8] \)}{\( C = [-3;-1;-8]; D = [-4;5;-13] \)}{}{}}
\LANGes{
\Question{
Dados los puntos \( A = [1;-1;2] \), \( B = [0;5;-3] \), \( S = [2;0;5] \). El punto \( S \) es el centro de un paralelogramo \( ABCD \). Determina las coordenadas de los vértices \( C \)  y \( D \).}
{\( C = [3;1;8]; D = [4;-5;13] \)}{\( C = [4;-5;13]; D = [3;1;8] \)}{\( C = [1;1;3]; D = [2;-5;8] \)}{\( C = [-3;-1;-8]; D = [-4;5;-13] \)}{}{}}
\End

\Data\ID{1103030702}
\Updated{2020-07-15 03:07:12}
\Level{A}
\SubArea{Points and vectors}
\IMG{1103030702.pdf}
\LANGen{
\Question{
We are given points \( A = [1;-1;2] \), \( B = [0;5;-3] \), \( S = [2;0;5] \). Point \( S \) is the centre of a parallelogram \( ABCD \). Evaluate the length of \( AD \).}
{\( |AD|=\sqrt{146} \)}{\( |AD|=\sqrt{114} \)}{\( |AD|=\sqrt{44} \)}{\( |AD|=\sqrt{172} \)}{}{}}
\LANGpl{
\Question{
Dane są punkty \( A = [1;-1;2] \), \( B = [0;5;-3] \), \( S = [2;0;5] \). Punkt \( S \) to środek równoległoboku \( ABCD \). Oblicz długość  \( AD \).}
{\( |AD|=\sqrt{146} \)}{\( |AD|=\sqrt{114} \)}{\( |AD|=\sqrt{44} \)}{\( |AD|=\sqrt{172} \)}{}{}}
\LANGcs{
\Question{
Jsou dány body \( A = [1;-1;2] \), \( B = [0;5;-3] \), \( S = [2;0;5] \). Bod \( S \) je středem rovnoběžníku \( ABCD \). Určete délku jeho strany \( AD \).}
{\( |AD|=\sqrt{146} \)}{\( |AD|=\sqrt{114} \)}{\( |AD|=\sqrt{44} \)}{\( |AD|=\sqrt{172} \)}{}{}}
\LANGsk{
\Question{
Sú dané body \( A = [1;-1;2] \), \( B = [0;5;-3] \), \( S = [2;0;5] \). Bod \( S \) je stredom rovnobežníka \( ABCD \). Určte dĺžku jeho strany \( AD \).}
{\( |AD|=\sqrt{146} \)}{\( |AD|=\sqrt{114} \)}{\( |AD|=\sqrt{44} \)}{\( |AD|=\sqrt{172} \)}{}{}}
\LANGes{
\Question{
Dados los puntos \( A = [1;-1;2] \), \( B = [0;5;-3] \), \( S = [2;0;5] \). El punto \( S \) es el centro de un paralelogramo \( ABCD \). Calcula la longitud de \( AD \).}
{\( |AD|=\sqrt{146} \)}{\( |AD|=\sqrt{114} \)}{\( |AD|=\sqrt{44} \)}{\( |AD|=\sqrt{172} \)}{}{}}
\End

\Data\ID{1103030703}
\Updated{2020-07-15 03:07:12}
\Level{A}
\SubArea{Points and vectors}
\IMG{1103030703.pdf}
\LANGen{
\Question{
We are given points \( A = [2;1] \), \( B = [4;-1] \), and \( T = [6;2] \), where point \( T \) is the centroid of triangle \( ABC \). Find the coordinates of \( C \), which is the vertex of \( ABC \).}
{\( C = [12;6] \)}{\( C = [8;4] \)}{\( C = [9;6] \)}{\( C = [8;5] \)}{}{}}
\LANGpl{
\Question{
Dane są punkty \( A = [2;1] \), \( B = [4;-1] \), i \( T = [6;2] \), punkt  \( T \) to środek ciężkości trójkąta \( ABC \). Określ współrzędne punktu \( C \), który jest wierzchołkiem trójkąta  \( ABC \).}
{\( C = [12;6] \)}{\( C = [8;4] \)}{\( C = [9;6] \)}{\( C = [8;5] \)}{}{}}
\LANGcs{
\Question{
Jsou dány body \( A = [2;1] \), \( B = [4;-1] \) a \( T = [6;2] \). Bod \( T \) je těžištěm trojúhelníku \( ABC \). Určete souřadnice vrcholu \( C \) tohoto trojúhelníku.}
{\( C = [12;6] \)}{\( C = [8;4] \)}{\( C = [9;6] \)}{\( C = [8;5] \)}{}{}}
\LANGsk{
\Question{
Sú dané body \( A = [2;1] \), \( B = [4;-1] \) a \( T = [6;2] \). Bod \( T \) je ťažiskom trojuholníka \( ABC \). Určte súradnice vrcholu \( C \) tohoto trojuholníka.}
{\( C = [12;6] \)}{\( C = [8;4] \)}{\( C = [9;6] \)}{\( C = [8;5] \)}{}{}}
\LANGes{
\Question{
Dados los puntos \( A = [2;1] \), \( B = [4;-1] \), y \( T = [6;2] \), donde el punto \( T \) es el baricentro de un triángulo \( ABC \). Determina las coordenadas de \( C \), que es el vértice del triángulo \( ABC \).}
{\( C = [12;6] \)}{\( C = [8;4] \)}{\( C = [9;6] \)}{\( C = [8;5] \)}{}{}}
\End

\Data\ID{1103030704}
\Updated{2020-07-15 03:07:12}
\Level{A}
\SubArea{Points and vectors}
\IMG{1103030704.pdf}
\LANGen{
\Question{
We are given points \( A = [2;1] \), \( B = [4;-1] \), and \( T = [6;2] \), where point \( T \) is the centroid of triangle \( ABC \). Find the length of the median of triangle \( ABC \) to side \( AC \).}
{\( |t_b|=\frac{\sqrt{117}}2 \)}{\( |t_b|=\frac{\sqrt{45}}2 \)}{\( |t_b|=\frac{\sqrt{153}}2 \)}{\( |t_b|=\sqrt{117} \)}{}{}}
\LANGpl{
\Question{
Dane są punkty \( A = [2;1] \), \( B = [4;-1] \), i \( T = [6;2] \), gdzie punkt \( T \) to środek ciężkości trójkąta \( ABC \). Oblicz długość środkowej    trójkąta  \( ABC \)  na bok \( AC \).}
{\( |t_b|=\frac{\sqrt{117}}2 \)}{\( |t_b|=\frac{\sqrt{45}}2 \)}{\( |t_b|=\frac{\sqrt{153}}2 \)}{\( |t_b|=\sqrt{117} \)}{}{}}
\LANGcs{
\Question{
Jsou dány body \( A = [2;1] \), \( B = [4;-1] \) a \( T = [6;2] \). Bod \( T \) je těžištěm trojúhelníku \( ABC \). Určete délku těžnice na stranu \( AC \) v tomto trojúhelníku.}
{\( |t_b|=\frac{\sqrt{117}}2 \)}{\( |t_b|=\frac{\sqrt{45}}2 \)}{\( |t_b|=\frac{\sqrt{153}}2 \)}{\( |t_b|=\sqrt{117} \)}{}{}}
\LANGsk{
\Question{
Sú dané body \( A = [2;1] \), \( B = [4;-1] \) a \( T = [6;2] \). Bod \( T \) je ťažiskom trojuholníka \( ABC \). Určte dĺžku ťažnice na stranu \( AC \) v tomto trojuholníku.}
{\( |t_b|=\frac{\sqrt{117}}2 \)}{\( |t_b|=\frac{\sqrt{45}}2 \)}{\( |t_b|=\frac{\sqrt{153}}2 \)}{\( |t_b|=\sqrt{117} \)}{}{}}
\LANGes{
\Question{
Dados los puntos \( A = [2;1] \), \( B = [4;-1] \), y \( T = [6;2] \), donde el punto \( T \) es el baricentro de un triángulo \( ABC \). Determina la longitud de la mediana del triángulo \( ABC \) al lado \( AC \).}
{\( |t_b|=\frac{\sqrt{117}}2 \)}{\( |t_b|=\frac{\sqrt{45}}2 \)}{\( |t_b|=\frac{\sqrt{153}}2 \)}{\( |t_b|=\sqrt{117} \)}{}{}}
\End

\Data\ID{1103030705}
\Updated{2020-07-15 03:07:12}
\Level{A}
\SubArea{Points and vectors}
\IMG{1103030705.pdf}
\LANGen{
\Question{
Let there be a triangle KLM and vectors \( \overrightarrow{a} \), \( \overrightarrow{c} \) in the coordinate system. Triangle \( KLM \) and vectors \( \overrightarrow{a} \), \( \overrightarrow{c} \) are
given in the coordinate system shown in the picture. Point T is the centroid of the triangle KLM. Express vector \( \overrightarrow{x} \), where \( \overrightarrow{x}=\overrightarrow{KT} \) as a linear combination of \( \overrightarrow{a} \) and \( \overrightarrow{c} \) and evaluate \( \left|\overrightarrow{x}\right| \).}
{\( \overrightarrow{x}=\frac13 \overrightarrow{a}+\frac13 \overrightarrow{c} \), \( \left|\overrightarrow{x}\right|=5 \)}{\( \overrightarrow{x}=\frac23 \overrightarrow{a}+\frac23 \overrightarrow{c} \), \( \left|\overrightarrow{x}\right|=10 \)}{\( \overrightarrow{x}=\frac12 \overrightarrow{a}+\frac12 \overrightarrow{c} \), \( \left|\overrightarrow{x}\right|=\frac{15}2 \)}{\( \overrightarrow{x}=\frac14 \overrightarrow{a}+\frac14 \overrightarrow{c} \), \( \left|\overrightarrow{x}\right|=\frac{225}{12} \)}{}{}}
\LANGpl{
\Question{
Dany jest trójkąt  KLM  i wektory \( \overrightarrow{a} \), \( \overrightarrow{c} \) w układzie współrzędnych. 
 Punkt  T to środek ciężkości trójkąta KLM. Przedstaw wektor \( \overrightarrow{x} \), gdzie \( \overrightarrow{x}=\overrightarrow{KT} \) jest liniową kombinacją  \( \overrightarrow{a} \) i \( \overrightarrow{c} \) i oblicz \( \left|\overrightarrow{x}\right| \).}
{\( \overrightarrow{x}=\frac13 \overrightarrow{a}+\frac13 \overrightarrow{c} \), \( \left|\overrightarrow{x}\right|=5 \)}{\( \overrightarrow{x}=\frac23 \overrightarrow{a}+\frac23 \overrightarrow{c} \), \( \left|\overrightarrow{x}\right|=10 \)}{\( \overrightarrow{x}=\frac12 \overrightarrow{a}+\frac12 \overrightarrow{c} \), \( \left|\overrightarrow{x}\right|=\frac{15}2 \)}{\( \overrightarrow{x}=\frac14 \overrightarrow{a}+\frac14 \overrightarrow{c} \), \( \left|\overrightarrow{x}\right|=\frac{225}{12} \)}{}{}}
\LANGcs{
\Question{
V souřadném systému je dán trojúhelník KLM a vektory \( \overrightarrow{a} \), \( \overrightarrow{c} \). Vektor \( \overrightarrow{x}=\overrightarrow{KT} \), kde T je těžiště trojúhelníka KLM, vyjádřete jako lineární kombinaci daných vektorů \( \overrightarrow{a} \), \( \overrightarrow{c} \) a vypočtěte \( \left|\overrightarrow{x}\right| \).}
{\( \overrightarrow{x}=\frac13 \overrightarrow{a}+\frac13 \overrightarrow{c} \), \( \left|\overrightarrow{x}\right|=5 \)}{\( \overrightarrow{x}=\frac23 \overrightarrow{a}+\frac23 \overrightarrow{c} \), \( \left|\overrightarrow{x}\right|=10 \)}{\( \overrightarrow{x}=\frac12 \overrightarrow{a}+\frac12 \overrightarrow{c} \), \( \left|\overrightarrow{x}\right|=\frac{15}2 \)}{\( \overrightarrow{x}=\frac14 \overrightarrow{a}+\frac14 \overrightarrow{c} \), \( \left|\overrightarrow{x}\right|=\frac{225}{12} \)}{}{}}
\LANGsk{
\Question{
V súradnicovom systéme je daný trojuholník KLM a vektory \( \overrightarrow{a} \), \( \overrightarrow{c} \). Vektor \( \overrightarrow{x}=\overrightarrow{KT} \), kde T je ťažisko trojuholníka KLM, vyjadrite ako lineárnu kombináciu daných vektorov \( \overrightarrow{a} \), \( \overrightarrow{c} \) a vypočítajte \( \left|\overrightarrow{x}\right| \).}
{\( \overrightarrow{x}=\frac13 \overrightarrow{a}+\frac13 \overrightarrow{c} \), \( \left|\overrightarrow{x}\right|=5 \)}{\( \overrightarrow{x}=\frac23 \overrightarrow{a}+\frac23 \overrightarrow{c} \), \( \left|\overrightarrow{x}\right|=10 \)}{\( \overrightarrow{x}=\frac12 \overrightarrow{a}+\frac12 \overrightarrow{c} \), \( \left|\overrightarrow{x}\right|=\frac{15}2 \)}{\( \overrightarrow{x}=\frac14 \overrightarrow{a}+\frac14 \overrightarrow{c} \), \( \left|\overrightarrow{x}\right|=\frac{225}{12} \)}{}{}}
\LANGes{
\Question{
Dado el triángulo  \( KLM \) y vectores \( \overrightarrow{a} \), \( \overrightarrow{c} \) en un sistema de coordenadas.  El punto \( T \)  es el baricentro del triángulo \( KLM \).  Expresa el vector \( \overrightarrow{x} \), donde \( \overrightarrow{x}=\overrightarrow{KT} \) como una combinación lineal de vectores \( \overrightarrow{a} \) y \( \overrightarrow{c} \) y calcula \( \left|\overrightarrow{x}\right| \).}
{\( \overrightarrow{x}=\frac13 \overrightarrow{a}+\frac13 \overrightarrow{c} \), \( \left|\overrightarrow{x}\right|=5 \)}{\( \overrightarrow{x}=\frac23 \overrightarrow{a}+\frac23 \overrightarrow{c} \), \( \left|\overrightarrow{x}\right|=10 \)}{\( \overrightarrow{x}=\frac12 \overrightarrow{a}+\frac12 \overrightarrow{c} \), \( \left|\overrightarrow{x}\right|=\frac{15}2 \)}{\( \overrightarrow{x}=\frac14 \overrightarrow{a}+\frac14 \overrightarrow{c} \), \( \left|\overrightarrow{x}\right|=\frac{225}{12} \)}{}{}}
\End

\Data\ID{9000100701}
\Updated{2020-07-15 03:07:37}
\Level{A}
\SubArea{Points and vectors}
\IMG{001007_01-figure0.pdf}
\LANGen{
\Question{
Given vectors \(\vec{a}\),
\(\vec{b}\),
\(\vec{c}\),
\(\vec{d}\), find
\(\vec{a} +\vec{ b} +\vec{ c} +\vec{ d}\).}
{\((-1;-2)\)}{\((17;7)\)}{\((6;10)\)}{\((2;-3)\)}{}{}}
\LANGpl{
\Question{
Dane są wektory \(\vec{a}\),
\(\vec{b}\),
\(\vec{c}\),
\(\vec{d}\), wyznacz
\(\vec{a} +\vec{ b} +\vec{ c} +\vec{ d}\).}
{\((-1;-2)\)}{\((17;7)\)}{\((6;10)\)}{\((2;-3)\)}{}{}}
\LANGcs{
\Question{
Na obrázku jsou zobrazeny vektory \(\vec{a}\),
\(\vec{b}\),
\(\vec{c}\),
\(\vec{d}\).
Součtem těchto vektorů je vektor:}
{\((-1;-2)\)}{\((17;7)\)}{\((6;10)\)}{\((2;-3)\)}{}{}}
\LANGsk{
\Question{
Na obrázku sú zobrazené vektory \(\vec{a}\),
\(\vec{b}\),
\(\vec{c}\),
\(\vec{d}\).
Súčtom týchto vektorov je vektor:}
{\((-1;-2)\)}{\((17;7)\)}{\((6;10)\)}{\((2;-3)\)}{}{}}
\LANGes{
\Question{
Dados los vectores \(\vec{a}\),
\(\vec{b}\),
\(\vec{c}\),
\(\vec{d}\), halla
\(\vec{a} +\vec{ b} +\vec{ c} +\vec{ d}\).}
{\((-1;-2)\)}{\((17;7)\)}{\((6;10)\)}{\((2;-3)\)}{}{}}
\End

\Data\ID{9000100702}
\Updated{2020-07-15 03:07:37}
\Level{A}
\SubArea{Points and vectors}
\LANGen{
\Question{
Given vectors \(\vec{a} = (x;-1)\),
\(\vec{b} = (3;y)\), find
\(x\) and
\(y\) such
that \(2\vec{a} - 3\vec{b} = (-5;4)\).}
{\(x = 2\),
\(y = -2\)}{\(x = -2\),
\(y = 2\)}{\(x = 2\),
\(y = 5\)}{\(x = 2\),
\(y = 2\)}{}{}}
\LANGpl{
\Question{
Dane są wektory \(\vec{a} = (x;-1)\),
\(\vec{b} = (3;y)\), wskaż 
\(x\) i
\(y\) jeśli
 \(2\vec{a} - 3\vec{b} = (-5;4)\).}
{\(x = 2\),
\(y = -2\)}{\(x = -2\),
\(y = 2\)}{\(x = 2\),
\(y = 5\)}{\(x = 2\),
\(y = 2\)}{}{}}
\LANGcs{
\Question{
Jsou dány vektory \(\vec{a} = (x;-1)\),
\(\vec{b} = (3;y)\). Určete
souřadnice \(x\) 
a \(y\) tak, aby
platilo \(2\vec{a} - 3\vec{b} = (-5;4)\).}
{\(x = 2,\ y = -2\)}{\(x = -2,\ y = 2\)}{\(x = 2,\ y = 5\)}{\(x = 2,\ y = 2\)}{}{}}
\LANGsk{
\Question{
Sú dané vektory \(\vec{a} = (x;-1)\),
\(\vec{b} = (3;y)\). Určte
súradnice \(x\) 
a \(y\) tak, aby
platilo \(2\vec{a} - 3\vec{b} = (-5;4)\).}
{\(x = 2,\ y = -2\)}{\(x = -2,\ y = 2\)}{\(x = 2,\ y = 5\)}{\(x = 2,\ y = 2\)}{}{}}
\LANGes{
\Question{
Dados los vectores \(\vec{a} = (x;-1)\),
\(\vec{b} = (3;y)\), halla
\(x\) y
\(y\) suponiendo que \(2\vec{a} - 3\vec{b} = (-5;4)\).}
{\(x = 2\),
\(y = -2\)}{\(x = -2\),
\(y = 2\)}{\(x = 2\),
\(y = 5\)}{\(x = 2\),
\(y = 2\)}{}{}}
\End

\Data\ID{9000100703}
\Updated{2020-07-15 03:07:37}
\Level{A}
\SubArea{Points and vectors}
\IMG{001007_03-figure0.pdf}
\LANGen{
\Question{
Given vectors \(\vec{a}\),
\(\vec{b}\),
\(\vec{c}\), find
\(2\vec{a} + 3\vec{b} -\vec{ c}\).}
{\((13;-5)\)}{\((7;-5)\)}{\((7;7)\)}{\((7;0)\)}{}{}}
\LANGpl{
\Question{
Dane są wektory \(\vec{a}\),
\(\vec{b}\),
\(\vec{c}\), wyznacz
\(2\vec{a} + 3\vec{b} -\vec{ c}\).}
{\((13;-5)\)}{\((7;-5)\)}{\((7;7)\)}{\((7;0)\)}{}{}}
\LANGcs{
\Question{
Na obrázku jsou zobrazeny vektory \(\vec{a}\),
\(\vec{b}\),
\(\vec{c}\). Vektor
\(\vec{u} = 2\vec{a} + 3\vec{b} -\vec{ c}\) je roven:}
{\((13;-5)\)}{\((7;-5)\)}{\((7;7)\)}{\((7;0)\)}{}{}}
\LANGsk{
\Question{
Na obrázku sú zobrazené vektory \(\vec{a}\),
\(\vec{b}\),
\(\vec{c}\). Vektor
\(\vec{u} = 2\vec{a} + 3\vec{b} -\vec{ c}\) je rovný:}
{\((13;-5)\)}{\((7;-5)\)}{\((7;7)\)}{\((7;0)\)}{}{}}
\LANGes{
\Question{
Dados los vectores  \(\vec{a}\),
\(\vec{b}\),
\(\vec{c}\), halla
\(2\vec{a} + 3\vec{b} -\vec{ c}\).}
{\((13;-5)\)}{\((7;-5)\)}{\((7;7)\)}{\((7;0)\)}{}{}}
\End

\Data\ID{9000100704}
\Updated{2020-07-15 03:07:37}
\Level{A}
\SubArea{Points and vectors}
\LANGen{
\Question{
Given vectors \(\vec{a} = (2;2;-3)\),
\(\vec{b} = (-1;0;1)\),
\(\vec{c} = (0;-2;1)\), find the length
of the vector \(\vec{u} =\vec{ a} - 2\vec{b} + 3\vec{c}\).}
{\(|\vec{u}| = 6\)}{\(|\vec{u}| = 5\)}{\(|\vec{u}| = 4\)}{\(|\vec{u}| = 3\)}{}{}}
\LANGpl{
\Question{
Dane są wektory \(\vec{a} = (2;2;-3)\),
\(\vec{b} = (-1;0;1)\),
\(\vec{c} = (0;-2;1)\), wyznacz długość wektora \(\vec{u} =\vec{ a} - 2\vec{b} + 3\vec{c}\).}
{\(|\vec{u}| = 6\)}{\(|\vec{u}| = 5\)}{\(|\vec{u}| = 4\)}{\(|\vec{u}| = 3\)}{}{}}
\LANGcs{
\Question{
Jsou dány vektory \(\vec{a} = (2;2;-3)\),
\(\vec{b} = (-1;0;1)\),
\(\vec{c} = (0;-2;1)\). Pro velikost
vektoru \(\vec{u} =\vec{ a} - 2\vec{b} + 3\vec{c}\) 
platí:}
{\(|\vec{u}| = 6\)}{\(|\vec{u}| = 5\)}{\(|\vec{u}| = 4\)}{\(|\vec{u}| = 3\)}{}{}}
\LANGsk{
\Question{
Sú dané vektory \(\vec{a} = (2;2;-3)\),
\(\vec{b} = (-1;0;1)\),
\(\vec{c} = (0;-2;1)\). Pre veľkosť
vektorov \(\vec{u} =\vec{ a} - 2\vec{b} + 3\vec{c}\) 
platí:}
{\(|\vec{u}| = 6\)}{\(|\vec{u}| = 5\)}{\(|\vec{u}| = 4\)}{\(|\vec{u}| = 3\)}{}{}}
\LANGes{
\Question{
Dados los vectores  \(\vec{a} = (2;2;-3)\),
\(\vec{b} = (-1;0;1)\),
\(\vec{c} = (0;-2;1)\), halla la longitud de un vector
\(\vec{u} =\vec{ a} - 2\vec{b} + 3\vec{c}\).}
{\(|\vec{u}| = 6\)}{\(|\vec{u}| = 5\)}{\(|\vec{u}| = 4\)}{\(|\vec{u}| = 3\)}{}{}}
\End

\Data\ID{9000100705}
\Updated{2021-06-18 05:06:11}
\Level{A}
\SubArea{Points and vectors}
\LANGen{
\Question{
Given vectors \(\vec{a} = (1;y;3)\) and 
\(\vec{b} = (2;-1;-2)\), find the value of coordinate 
\(y\) so 
that the vector \(\vec{u} = (-4;-1;12)\) is a
linear combination of \(\vec{a}\) 
and \(\vec{b}\).}
{\(y = -2\)}{\(y = 1\)}{\(y = -1\)}{\(y = 3\)}{}{}}
\LANGpl{
\Question{
Dane są wektory  \(\vec{a} = (1;y;3)\),
\(\vec{b} = (2;-1;-2)\), wskaż współrzędną
\(y\) tak, aby 
wektor \(\vec{u} = (-4;-1;12)\) był
liniową kombinacją wektorów \(\vec{a}\) 
i \(\vec{b}\).}
{\(y = -2\)}{\(y = 1\)}{\(y = -1\)}{\(y = 3\)}{}{}}
\LANGcs{
\Question{
Jsou dány vektory \(\vec{a} = (1;y_{a};3)\),
\(\vec{b} = (2;-1;-2)\). Určete souřadnici
\(y_{a}\) tak, aby vektor
\(\vec{u} = (-4;-1;12)\) byl lineární
kombinací vektorů \(\vec{a},\ \vec{b}\).}
{\(y_{a} = -2\)}{\(y_{a} = 1\)}{\(y_{a} = -1\)}{\(y_{a} = 3\)}{}{}}
\LANGsk{
\Question{
Sú dané vektory \(\vec{a} = (1;y_{a};3)\),
\(\vec{b} = (2;-1;-2)\). Určte súradnicu
\(y_{a}\) tak, aby vektor
\(\vec{u} = (-4;-1;12)\) bol lineárnou
kombináciou vektorov \(\vec{a},\ \vec{b}\).}
{\(y_{a} = -2\)}{\(y_{a} = 1\)}{\(y_{a} = -1\)}{\(y_{a} = 3\)}{}{}}
\LANGes{
\Question{
Dados los vectores \(\vec{a} = (1;y;3)\),
\(\vec{b} = (2;-1;-2)\), halla la coordenada
\(y\) que garantiza que el vector \(\vec{u} = (-4;-1;12)\) es una combinación lineal de \(\vec{a}\) 
y \(\vec{b}\).}
{\(y = -2\)}{\(y = 1\)}{\(y = -1\)}{\(y = 3\)}{}{}}
\End

\Data\ID{9000100710}
\Updated{2020-07-15 03:07:37}
\Level{A}
\SubArea{Points and vectors}
\LANGen{
\Question{
Given points \(A = [-3;2]\) and
\(B = [1;y]\), find the values of
\(y\) which ensure the
length of the vector \(\overrightarrow{AB } \) 
equals \(5\).}
{\(y_{1} = -1\),
\(y_{2} = 5\)}{\(y_{1} = -1\),
\(y_{2} = 1\)}{\(y_{1} = 1\),
\(y_{2} = 5\)}{\(y_{1} = 5\),
\(y_{2} = -5\)}{}{}}
\LANGpl{
\Question{
Dane są punkty \(A = [-3;2]\) i
\(B = [1;y]\), określ wartość
\(y\) tak, aby
długość wektora \(\overrightarrow{AB } \) 
wynosiła \(5\).}
{\(y_{1} = -1\),
\(y_{2} = 5\)}{\(y_{1} = -1\),
\(y_{2} = 1\)}{\(y_{1} = 1\),
\(y_{2} = 5\)}{\(y_{1} = 5\),
\(y_{2} = -5\)}{}{}}
\LANGcs{
\Question{
Jsou dány body \(A = [-3;2]\) 
a \(B = [1;y]\). Určete
všechny hodnoty \(y\) 
tak, aby platilo \(|\overrightarrow{AB } | = 5\).}
{\(y = -1,\ y = 5\)}{\(y = -1,\ y = 1\)}{\(y = 1,\ y = 5\)}{\(y = 5,\ y = -5\)}{}{}}
\LANGsk{
\Question{
Sú dané body \(A = [-3;2]\) 
a \(B = [1;y]\). Určte
všetky hodnoty \(y\) 
tak, aby platilo \(|\overrightarrow{AB } | = 5\).}
{\(y = -1,\ y = 5\)}{\(y = -1,\ y = 1\)}{\(y = 1,\ y = 5\)}{\(y = 5,\ y = -5\)}{}{}}
\LANGes{
\Question{
Dados los puntos \(A = [-3;2]\) y
\(B = [1;y]\), halla los valores de
\(y\) para que la longitud del vector \(\overrightarrow{AB } \) 
equivalga a \(5\).}
{\(y_{1} = -1\),
\(y_{2} = 5\)}{\(y_{1} = -1\),
\(y_{2} = 1\)}{\(y_{1} = 1\),
\(y_{2} = 5\)}{\(y_{1} = 5\),
\(y_{2} = -5\)}{}{}}
\End

\Data\ID{9000101801}
\Updated{2020-07-15 03:07:37}
\Level{A}
\SubArea{Points and vectors}
\LANGen{
\Question{
Among vectors \(\vec{a} = (-1;2;0)\),
\(\vec{b} = (2;1;2)\),
\(\vec{c} = (1;3;0)\) and
\(\vec{d} = (-3;0;0)\) find a
pair of vectors with equal length.}
{\(\vec{b}\),
\(\vec{d}\)}{\(\vec{a}\),
\(\vec{c}\)}{\(\vec{a}\),
\(\vec{d}\)}{\(\vec{b}\),
\(\vec{c}\)}{}{}}
\LANGpl{
\Question{
Dane są wektory \(\vec{a} = (-1;2;0)\),
\(\vec{b} = (2;1;2)\),
\(\vec{c} = (1;3;0)\) i
\(\vec{d} = (-3;0;0)\) wskaż 
parę wektorów o tej samej długości.}
{\(\vec{b}\),
\(\vec{d}\)}{\(\vec{a}\),
\(\vec{c}\)}{\(\vec{a}\),
\(\vec{d}\)}{\(\vec{b}\),
\(\vec{c}\)}{}{}}
\LANGcs{
\Question{
Jsou dány vektory \(\vec{a} = (-1;2;0)\),
\(\vec{b} = (2;1;2)\),
\(\vec{c} = (1;3;0)\),
\(\vec{d} = (-3;0;0)\). Pro
kterou dvojici vektorů platí, že mají stejnou velikost?}
{\(\vec{b},\ \vec{d}\)}{\(\vec{a},\ \vec{c}\)}{\(\vec{a},\ \vec{d}\)}{\(\vec{b},\ \vec{c}\)}{}{}}
\LANGsk{
\Question{
Sú dané vektory \(\vec{a} = (-1;2;0)\),
\(\vec{b} = (2;1;2)\),
\(\vec{c} = (1;3;0)\),
\(\vec{d} = (-3;0;0)\). Pre
ktorú dvojicu vektorov platí, že majú rovnakú veľkosť?}
{\(\vec{b},\ \vec{d}\)}{\(\vec{a},\ \vec{c}\)}{\(\vec{a},\ \vec{d}\)}{\(\vec{b},\ \vec{c}\)}{}{}}
\LANGes{
\Question{
Dados los vectores \(\vec{a} = (-1;2;0)\),
\(\vec{b} = (2;1;2)\),
\(\vec{c} = (1;3;0)\) y
\(\vec{d} = (-3;0;0)\).  ¿Qué par de vectores tienen la misma longitud?}
{\(\vec{b}\),
\(\vec{d}\)}{\(\vec{a}\),
\(\vec{c}\)}{\(\vec{a}\),
\(\vec{d}\)}{\(\vec{b}\),
\(\vec{c}\)}{}{}}
\End

\Data\ID{9000101803}
\Updated{2020-07-15 03:07:37}
\Level{A}
\SubArea{Points and vectors}
\LANGen{
\Question{
In the following list identify a pair of points
\(C\),
\(D\) such that the
vector \(\overrightarrow{CD } \) is not
equal to the vector \(\overrightarrow{AB } \) 
where \(A = [1;3;-2]\) 
and \(B = [-2;4;3]\).}
{\(C = [1;-2;3],\ D = [-2;-1;-2]\)}{\(C = [6;1;-4],\ D = [3;2;1]\)}{\(C = [-3;5;7],\ D = [-6;6;12]\)}{\(C = [-3;8;14],\ D = [-6;9;19]\)}{}{}}
\LANGpl{
\Question{
Wskaż parę punktów
\(C\),
\(D\) tak, aby
wektor \(\overrightarrow{CD } \) nie był
równy wektorowi \(\overrightarrow{AB } \) 
jeśli \(A = [1;3;-2]\) 
i \(B = [-2;4;3]\).}
{\(C = [1;-2;3],\ D = [-2;-1;-2]\)}{\(C = [6;1;-4],\ D = [3;2;1]\)}{\(C = [-3;5;7],\ D = [-6;6;12]\)}{\(C = [-3;8;14],\ D = [-6;9;19]\)}{}{}}
\LANGcs{
\Question{
Jsou dány body \(A = [1;3;-2]\) 
a \(B = [-2;4;3]\). Vyberte
dvojici bodů \(C\),
\(D\) tak, aby se vektor \(\overrightarrow{CD } \) nerovnal vektoru \(\overrightarrow{AB } \).}
{\(C = [1;-2;3],\ D = [-2;-1;-2]\)}{\(C = [6;1;-4],\ D = [3;2;1]\)}{\(C = [-3;5;7],\ D = [-6;6;12]\)}{\(C = [-3;8;14],\ D = [-6;9;19]\)}{}{}}
\LANGsk{
\Question{
Sú dané body \(A = [1;3;-2]\) 
a \(B = [-2;4;3]\). Vyberte
dvojicu bodov \(C\),
\(D\) tak, aby orientovaná
úsečka \(\overrightarrow{CD } \) nebola
umiestnením vektora \(\overrightarrow{AB } \).}
{\(C = [1;-2;3],\ D = [-2;-1;-2]\)}{\(C = [6;1;-4],\ D = [3;2;1]\)}{\(C = [-3;5;7],\ D = [-6;6;12]\)}{\(C = [-3;8;14],\ D = [-6;9;19]\)}{}{}}
\LANGes{
\Question{
En la siguiente lista, identifica un par de puntos
\(C\),
\(D\) suponiendo que el vector \(\overrightarrow{CD } \) no equivale al vector
\(\overrightarrow{AB } \) 
donde \(A = [1;3;-2]\) 
y \(B = [-2;4;3]\).}
{\(C = [1;-2;3],\ D = [-2;-1;-2]\)}{\(C = [6;1;-4],\ D = [3;2;1]\)}{\(C = [-3;5;7],\ D = [-6;6;12]\)}{\(C = [-3;8;14],\ D = [-6;9;19]\)}{}{}}
\End

\Data\ID{9000101804}
\Updated{2020-07-15 03:07:37}
\Level{A}
\SubArea{Points and vectors}
\LANGen{
\Question{
In the following list identify a valid relation involving the vectors
\(\vec{a} = (2;-3)\),
\(\vec{b} = (1;3)\) and
\(\vec{c} = (5;-3)\).}
{\(\vec{c} = 2\vec{a} +\vec{ b}\)}{\(\vec{b} = \frac{1} 
{2}\vec{a} +\vec{ c}\)}{\(2\vec{a} +\vec{ b} +\vec{ c} =\vec{ o}\)}{\(\vec{a} = \frac{1} 
{2}\vec{b} +\vec{ c}\)}{}{}}
\LANGpl{
\Question{
Wskaż poprawną zależność między wektorami
\(\vec{a} = (2;-3)\),
\(\vec{b} = (1;3)\) i
\(\vec{c} = (5;-3)\).}
{\(\vec{c} = 2\vec{a} +\vec{ b}\)}{\(\vec{b} = \frac{1} 
{2}\vec{a} +\vec{ c}\)}{\(2\vec{a} +\vec{ b} +\vec{ c} =\vec{ o}\)}{\(\vec{a} = \frac{1} 
{2}\vec{b} +\vec{ c}\)}{}{}}
\LANGcs{
\Question{
Jsou dány vektory \(\vec{a} = (2;-3)\),
\(\vec{b} = (1;3)\),
\(\vec{c} = (5;-3)\).
Který z následujících vztahů mezi vektory je správný?}
{\(\vec{c} = 2\vec{a} +\vec{ b}\)}{\(\vec{b} = \frac{1} 
{2}\vec{a} +\vec{ c}\)}{\(2\vec{a} +\vec{ b} +\vec{ c} =\vec{ o}\)}{\(\vec{a} = \frac{1} 
{2}\vec{b} +\vec{ c}\)}{}{}}
\LANGsk{
\Question{
Sú dané vektory \(\vec{a} = (2;-3)\),
\(\vec{b} = (1;3)\),
\(\vec{c} = (5;-3)\).
Ktorý z nasledujúcich vzťahov medzi vektormi je správny?}
{\(\vec{c} = 2\vec{a} +\vec{ b}\)}{\(\vec{b} = \frac{1} 
{2}\vec{a} +\vec{ c}\)}{\(2\vec{a} +\vec{ b} +\vec{ c} =\vec{ o}\)}{\(\vec{a} = \frac{1} 
{2}\vec{b} +\vec{ c}\)}{}{}}
\LANGes{
\Question{
Dados lo vectores \(\vec{a} = (2;-3)\),
\(\vec{b} = (1;3)\),
\(\vec{c} = (5;-3)\).
¿Cuál de las siguientes relaciones entre los vectores es correcta?}
{\(\vec{c} = 2\vec{a} +\vec{ b}\)}{\(\vec{b} = \frac{1} 
{2}\vec{a} +\vec{ c}\)}{\(2\vec{a} +\vec{ b} +\vec{ c} =\vec{ o}\)}{\(\vec{a} = \frac{1} 
{2}\vec{b} +\vec{ c}\)}{}{}}
\End

\Data\ID{9000101809}
\Updated{2020-07-15 03:07:37}
\Level{A}
\SubArea{Points and vectors}
\LANGen{
\Question{
Given point \(A = [3;2]\),
find all the points \(X\) 
on the \(y\)-axis
such that \(|AX| = 5\).}
{\(X_{1} = [0;-2],\ X_{2} = [0;6]\)}{\(X_{1} = [0;-6],\ X_{2} = [0;2]\)}{\(X_{1} = [0;-6],\ X_{2} = [0;-2]\)}{\(X_{1} = [0;2],\ X_{2} = [0;6]\)}{}{}}
\LANGpl{
\Question{
Dany jest punkt \(A = [3;2]\),
wskaż wszystkie punkty 
leżące na osi \(y\),
jeśli \(|AX| = 5\).}
{\(X_{1} = [0;-2],\ X_{2} = [0;6]\)}{\(X_{1} = [0;-6],\ X_{2} = [0;2]\)}{\(X_{1} = [0;-6],\ X_{2} = [0;-2]\)}{\(X_{1} = [0;2],\ X_{2} = [0;6]\)}{}{}}
\LANGcs{
\Question{
Je dán bod \(A = [3;2]\). Vyberte
všechny body \(X\) 
ležící na ose \(y\), pro
které platí, že \(|AX| = 5\).}
{\(X_{1} = [0;-2],\ X_{2} = [0;6]\)}{\(X_{1} = [0;-6],\ X_{2} = [0;2]\)}{\(X_{1} = [0;-6],\ X_{2} = [0;-2]\)}{\(X_{1} = [0;2],\ X_{2} = [0;6]\)}{}{}}
\LANGsk{
\Question{
Je daný bod \(A = [3;2]\). Vyberte
všetky body \(X\) 
ležiace na osy \(y\), pre
ktoré platí, že \(|AX| = 5\).}
{\(X_{1} = [0;-2],\ X_{2} = [0;6]\)}{\(X_{1} = [0;-6],\ X_{2} = [0;2]\)}{\(X_{1} = [0;-6],\ X_{2} = [0;-2]\)}{\(X_{1} = [0;2],\ X_{2} = [0;6]\)}{}{}}
\LANGes{
\Question{
Dado el punto \(A = [3;2]\).
Halla todos los puntos \(X\) 
en el eje \(y\)
suponiendo que \(|AX| = 5\).}
{\(X_{1} = [0;-2],\ X_{2} = [0;6]\)}{\(X_{1} = [0;-6],\ X_{2} = [0;2]\)}{\(X_{1} = [0;-6],\ X_{2} = [0;-2]\)}{\(X_{1} = [0;2],\ X_{2} = [0;6]\)}{}{}}
\End

\Data\ID{9000101810}
\Updated{2020-07-15 03:07:37}
\Level{A}
\SubArea{Points and vectors}
\LANGen{
\Question{
Given points \(A = [1;2]\) 
and \(B = [4;4]\), find the
point \(X\) on the
\(x\)-axis such that
the distance from \(X\) 
to \(B\) is a double of
the distance from \(X\) 
to \(A\).
Find all solutions of the problem.}
{\(X_{1} = [2;0],\ X_{2} = [-2;0]\)}{\(X = [2;0]\)}{\(X = [8;0]\)}{\(X_{1} = [2;0],\ X_{2} = [-4;0]\)}{}{}}
\LANGpl{
\Question{
Dane są punkty \(A = [1;2]\) 
i \(B = [4;4]\), zaznacz wszystkie 
punkty \(X\) na osi
\(x\), tak, aby ich odległość 
od punktu \(B\) była dwa razy większa niż od punktu \(A\).}
{\(X_{1} = [2;0],\ X_{2} = [-2;0]\)}{\(X = [2;0]\)}{\(X = [8;0]\)}{\(X_{1} = [2;0],\ X_{2} = [-4;0]\)}{}{}}
\LANGcs{
\Question{
Jsou dány body \(A = [1;2]\) 
a \(B = [4;4]\). Vyberte
všechny body \(X\) 
ležící na ose \(x\),
pro které platí, že jejich vzdálenost od bodu
\(B\) je dvakrát větší
než od bodu \(A\).}
{\(X_{1} = [2;0],\ X_{2} = [-2;0]\)}{\(X = [2;0]\)}{\(X = [8;0]\)}{\(X_{1} = [2;0],\ X_{2} = [-4;0]\)}{}{}}
\LANGsk{
\Question{
Sú dané body \(A = [1;2]\) 
a \(B = [4;4]\). Vyberte
všetky body \(X\) 
na osy \(x\),
pre ktoré platí, že ich vzdialenosť od bodu
\(B\) je dvakrát väčšia
ako od bodu \(A\).}
{\(X_{1} = [2;0],\ X_{2} = [-2;0]\)}{\(X = [2;0]\)}{\(X = [8;0]\)}{\(X_{1} = [2;0],\ X_{2} = [-4;0]\)}{}{}}
\LANGes{
\Question{
Dados los puntos \(A = [1;2]\) 
y \(B = [4;4]\). Halla el punto \(X\) en el eje
\(x\) suponiendo que la distancia de
 \(X\) 
a \(B\) es un doble de la distancia de \(X\) 
a \(A\).
Determina todas las soluciones del problema.}
{\(X_{1} = [2;0],\ X_{2} = [-2;0]\)}{\(X = [2;0]\)}{\(X = [8;0]\)}{\(X_{1} = [2;0],\ X_{2} = [-4;0]\)}{}{}}
\End

\Data\ID{1003030603}
\Updated{2020-07-15 03:07:12}
\Level{B}
\SubArea{Points and vectors}
\LANGen{
\Question{
Let \( \overrightarrow{v}=(12;5) \). Find all the vectors \( \overrightarrow{u} \) that are perpendicular to the vector  \( \overrightarrow{v} \) and have the length of \( 26 \).}
{\( \overrightarrow{u_1}  =(10;-24);\  \overrightarrow{u_2}=(-10; 24) \)}{\( \overrightarrow{u}=(10;-24) \)}{\( \overrightarrow{u_1}=\frac12 (5;-12);\ \overrightarrow{u_2}=\frac12 (-5; 12) \)}{\( \overrightarrow{u_1}=26\cdot(5;-12);\ \overrightarrow{u_2}=26\cdot(-5; 12) \)}{}{}}
\LANGpl{
\Question{
Dany jest wektor \( \overrightarrow{v}=(12;5) \). Wskaż wszystkie wektory \( \overrightarrow{u} \) prostopadłe do wektora  \( \overrightarrow{v} \), których długość jest równa \( 26 \).}
{\( \overrightarrow{u_1}  =(10;-24);\  \overrightarrow{u_2}=(-10; 24) \)}{\( \overrightarrow{u}=(10;-24) \)}{\( \overrightarrow{u_1}=\frac12 (5;-12);\ \overrightarrow{u_2}=\frac12 (-5; 12) \)}{\( \overrightarrow{u_1}=26\cdot(5;-12);\ \overrightarrow{u_2}=26\cdot(-5; 12) \)}{}{}}
\LANGcs{
\Question{
Je dán vektor \( \overrightarrow{v}=(12;5) \). Najděte všechny takové vektory \( \overrightarrow{u} \), které jsou kolmé k vektoru \( \overrightarrow{v} \) a mají velikost \( 26 \).}
{\( \overrightarrow{u_1}  =(10;-24);\  \overrightarrow{u_2}=(-10; 24) \)}{\( \overrightarrow{u}=(10;-24) \)}{\( \overrightarrow{u_1}=\frac12 (5;-12);\ \overrightarrow{u_2}=\frac12 (-5; 12) \)}{\( \overrightarrow{u_1}=26\cdot(5;-12);\ \overrightarrow{u_2}=26\cdot(-5; 12) \)}{}{}}
\LANGsk{
\Question{
Je daný vektor \( \overrightarrow{v}=(12;5) \). Nájdite všetky také vektory \( \overrightarrow{u} \), ktoré sú kolmé k vektoru \( \overrightarrow{v} \) a majú veľkosť \( 26 \).}
{\( \overrightarrow{u_1}  =(10;-24);\  \overrightarrow{u_2}=(-10; 24) \)}{\( \overrightarrow{u}=(10;-24) \)}{\( \overrightarrow{u_1}=\frac12 (5;-12);\ \overrightarrow{u_2}=\frac12 (-5; 12) \)}{\( \overrightarrow{u_1}=26\cdot(5;-12);\ \overrightarrow{u_2}=26\cdot(-5; 12) \)}{}{}}
\LANGes{
\Question{
Sea\( \overrightarrow{v}=(12;5) \). Determina todos los vectores \( \overrightarrow{u} \) que son perpendiculares al vector  \( \overrightarrow{v} \) y tienen la longitud de \( 26 \).}
{\( \overrightarrow{u_1}  =(10;-24);\  \overrightarrow{u_2}=(-10; 24) \)}{\( \overrightarrow{u}=(10;-24) \)}{\( \overrightarrow{u_1}=\frac12 (5;-12);\ \overrightarrow{u_2}=\frac12 (-5; 12) \)}{\( \overrightarrow{u_1}=26\cdot(5;-12);\ \overrightarrow{u_2}=26\cdot(-5; 12) \)}{}{}}
\End

\Data\ID{1003030604}
\Updated{2020-07-15 03:07:12}
\Level{B}
\SubArea{Points and vectors}
\LANGen{
\Question{
Let \( \overrightarrow{a}=(2;- 3) \)  and \( \overrightarrow{b}=(3;-2) \).  Find all the vectors \( \overrightarrow{c} \) such that 
\[ \overrightarrow{a}\cdot\overrightarrow{c}=8\ \text{ and }\  \overrightarrow{b}\cdot\overrightarrow{c}=27. \]}
{\( \overrightarrow{c}=(13;6) \)}{\( \overrightarrow{c_1}=(13;6);\ \overrightarrow{c_2}=(-13;-6) \)}{\( \overrightarrow{c}=(13k;6k);\ k\in\mathbb{R} \)}{\( \overrightarrow{c}=(-13;-6) \)}{}{}}
\LANGpl{
\Question{
Podano \( \overrightarrow{a}=(2;- 3) \)  i \( \overrightarrow{b}=(3;-2) \).  Wskaż wszystkie wektory \( \overrightarrow{c} \) tak, aby
\[ \overrightarrow{a}\cdot\overrightarrow{c}=8\ \text{ i }\  \overrightarrow{b}\cdot\overrightarrow{c}=27. \]}
{\( \overrightarrow{c}=(13;6) \)}{\( \overrightarrow{c_1}=(13;6);\ \overrightarrow{c_2}=(-13;-6) \)}{\( \overrightarrow{c}=(13k;6k);\ k\in\mathbb{R} \)}{\( \overrightarrow{c}=(-13;-6) \)}{}{}}
\LANGcs{
\Question{
Jsou dány vektory \( \overrightarrow{a}=(2;- 3) \)  a \( \overrightarrow{b}=(3;-2) \).  Najděte všechny takové vektory \( \overrightarrow{c} \), pro které platí 
\[ \overrightarrow{a}\cdot\overrightarrow{c}=8\ \text{ a }\  \overrightarrow{b}\cdot\overrightarrow{c}=27. \]}
{\( \overrightarrow{c}=(13;6) \)}{\( \overrightarrow{c_1}=(13;6);\ \overrightarrow{c_2}=(-13;-6) \)}{\( \overrightarrow{c}=(13k;6k);\ k\in\mathbb{R} \)}{\( \overrightarrow{c}=(-13;-6) \)}{}{}}
\LANGsk{
\Question{
Sú dané vektory \( \overrightarrow{a}=(2;- 3) \)  a \( \overrightarrow{b}=(3;-2) \).  Nájdite všetky také vektory \( \overrightarrow{c} \), pre ktoré platí 
\[ \overrightarrow{a}\cdot\overrightarrow{c}=8\ \text{ a }\  \overrightarrow{b}\cdot\overrightarrow{c}=27. \]}
{\( \overrightarrow{c}=(13;6) \)}{\( \overrightarrow{c_1}=(13;6);\ \overrightarrow{c_2}=(-13;-6) \)}{\( \overrightarrow{c}=(13k;6k);\ k\in\mathbb{R} \)}{\( \overrightarrow{c}=(-13;-6) \)}{}{}}
\LANGes{
\Question{
Sean \( \overrightarrow{a}=(2;- 3) \)  y \( \overrightarrow{b}=(3;-2) \).  Determina todos los vectores \( \overrightarrow{c} \) suponiendo que
\[ \overrightarrow{a}\cdot\overrightarrow{c}=8\ \text{ y }\  \overrightarrow{b}\cdot\overrightarrow{c}=27. \]}
{\( \overrightarrow{c}=(13;6) \)}{\( \overrightarrow{c_1}=(13;6);\ \overrightarrow{c_2}=(-13;-6) \)}{\( \overrightarrow{c}=(13k;6k);\ k\in\mathbb{R} \)}{\( \overrightarrow{c}=(-13;-6) \)}{}{}}
\End

\Data\ID{1003030605}
\Updated{2020-07-15 03:07:12}
\Level{B}
\SubArea{Points and vectors}
\LANGen{
\Question{
Let \( \overrightarrow{a}=(3;-5) \) and \( \overrightarrow{b}=(6;-10) \). Find all the vectors  \( \overrightarrow{c} \) such that 
\[ \overrightarrow{a}\cdot\overrightarrow{c}=11\ \text{ and }\ \overrightarrow{b}\cdot\overrightarrow{c}=22\text{ .} \]}
{\( \overrightarrow{c}=(2+5k;-1+3k);\ k\in\mathbb{R} \)}{\( \overrightarrow{c}_1=(7;2);\ \overrightarrow{c}_2=(-7;-2) \)}{\( \overrightarrow{c}=(2k;-k);\ k\in\mathbb{R} \)}{\( \overrightarrow{c}_1=(2;-1);\ \overrightarrow{c}_2=(-2;1) \)}{}{}}
\LANGpl{
\Question{
Dane są wektory \( \overrightarrow{a}=(3;-5) \) i \( \overrightarrow{b}=(6;-10) \). Wskaż wszystkie wektory  \( \overrightarrow{c} \) tak, aby
\[ \overrightarrow{a}\cdot\overrightarrow{c}=11\ \text{ i }\ \overrightarrow{b}\cdot\overrightarrow{c}=22\text{ .} \]}
{\( \overrightarrow{c}=(2+5k;-1+3k);\ k\in\mathbb{R} \)}{\( \overrightarrow{c}_1=(7;2);\ \overrightarrow{c}_2=(-7;-2) \)}{\( \overrightarrow{c}=(2k;-k);\ k\in\mathbb{R} \)}{\( \overrightarrow{c}_1=(2;-1);\ \overrightarrow{c}_2=(-2;1) \)}{}{}}
\LANGcs{
\Question{
Jsou dány vektory \( \overrightarrow{a}=(3;-5) \) a \( \overrightarrow{b}=(6;-10) \). Najděte všechny vektory \( \overrightarrow{c} \), pro které platí 
\[ \overrightarrow{a}\cdot\overrightarrow{c}=11\ \text{ a }\ \overrightarrow{b}\cdot\overrightarrow{c}=22\text{ .} \]}
{\( \overrightarrow{c}=(2+5k;-1+3k);\ k\in\mathbb{R} \)}{\( \overrightarrow{c}_1=(7;2);\ \overrightarrow{c}_2=(-7;-2) \)}{\( \overrightarrow{c}=(2k;-k);\ k\in\mathbb{R} \)}{\( \overrightarrow{c}_1=(2;-1);\ \overrightarrow{c}_2=(-2;1) \)}{}{}}
\LANGsk{
\Question{
Sú dané vektory \( \overrightarrow{a}=(3;-5) \) a \( \overrightarrow{b}=(6;-10) \). Nájdite všetky vektory \( \overrightarrow{c} \), pre ktoré platí 
\[ \overrightarrow{a}\cdot\overrightarrow{c}=11\ \text{ a }\ \overrightarrow{b}\cdot\overrightarrow{c}=22\text{ .} \]}
{\( \overrightarrow{c}=(2+5k;-1+3k);\ k\in\mathbb{R} \)}{\( \overrightarrow{c}_1=(7;2);\ \overrightarrow{c}_2=(-7;-2) \)}{\( \overrightarrow{c}=(2k;-k);\ k\in\mathbb{R} \)}{\( \overrightarrow{c}_1=(2;-1);\ \overrightarrow{c}_2=(-2;1) \)}{}{}}
\LANGes{
\Question{
Sean \( \overrightarrow{a}=(3;-5) \) y \( \overrightarrow{b}=(6;-10) \). Determina todos los vectores  \( \overrightarrow{c} \) suponiendo que 
\[ \overrightarrow{a}\cdot\overrightarrow{c}=11\ \text{ y }\ \overrightarrow{b}\cdot\overrightarrow{c}=22\text{ .} \]}
{\( \overrightarrow{c}=(2+5k;-1+3k);\ k\in\mathbb{R} \)}{\( \overrightarrow{c}_1=(7;2);\ \overrightarrow{c}_2=(-7;-2) \)}{\( \overrightarrow{c}=(2k;-k);\ k\in\mathbb{R} \)}{\( \overrightarrow{c}_1=(2;-1);\ \overrightarrow{c}_2=(-2;1) \)}{}{}}
\End

\Data\ID{1003040210}
\Updated{2020-07-15 03:07:12}
\Level{B}
\SubArea{Points and vectors}
\LANGen{
\Question{
Given the points $A = [3;3;0]$ and $B = [0;3;3]$, specify all the points $C$ lying on the $y$-axis, such that $|\measuredangle ABC|=\frac{\pi}3$ holds.}
{$C_1=[0;0;0];\ C_2=[0;6;0]$}{$C_1=[0;3;0];\ C_2=[0;9;0]$}{$C_1=[0;-3;0];\ C_2=[0;3;0]$}{$C_1=[0;-6;0];\ C_2=[0;6;0]$}{}{}}
\LANGpl{
\Question{
Dane są wektory $A = [3;3;0]$ i $B = [0;3;3]$, określ wszystkie punkty $C$ leżące na osi $y$ tak, aby  $|\measuredangle ABC|=\frac{\pi}3$.}
{$C_1=[0;0;0];\ C_2=[0;6;0]$}{$C_1=[0;3;0];\ C_2=[0;9;0]$}{$C_1=[0;-3;0];\ C_2=[0;3;0]$}{$C_1=[0;-6;0];\ C_2=[0;6;0]$}{}{}}
\LANGcs{
\Question{
Jsou dány body $A = [3;3;0]$ a $B = [0;3;3]$. Určete souřadnice všech bodů $C$ ležících na ose $y$, pro které platí $|\measuredangle ABC|=\frac{\pi}3$.}
{$C_1=[0;0;0];\ C_2=[0;6;0]$}{$C_1=[0;3;0];\ C_2=[0;9;0]$}{$C_1=[0;-3;0];\ C_2=[0;3;0]$}{$C_1=[0;-6;0];\ C_2=[0;6;0]$}{}{}}
\LANGsk{
\Question{
Sú dané body $A = [3;3;0]$ a $B = [0;3;3]$. Určte súradnice všetkých bodov $C$ ležiacich na osi $y$, pre ktoré platí $|\measuredangle ABC|=\frac{\pi}3$.}
{$C_1=[0;0;0];\ C_2=[0;6;0]$}{$C_1=[0;3;0];\ C_2=[0;9;0]$}{$C_1=[0;-3;0];\ C_2=[0;3;0]$}{$C_1=[0;-6;0];\ C_2=[0;6;0]$}{}{}}
\LANGes{
\Question{
Dados los puntos $A = [3;3;0]$ y $B = [0;3;3]$. Determina las coordenadas de todos los puntos $C$ que se encuentran en el eje $y$, suponiendo que $|\measuredangle ABC|=\frac{\pi}3$.}
{$C_1=[0;0;0];\ C_2=[0;6;0]$}{$C_1=[0;3;0];\ C_2=[0;9;0]$}{$C_1=[0;-3;0];\ C_2=[0;3;0]$}{$C_1=[0;-6;0];\ C_2=[0;6;0]$}{}{}}
\End

\Data\ID{1103021001}
\Updated{2020-07-15 03:07:12}
\Level{B}
\SubArea{Points and vectors}
\IMG{1103021001_0.pdf}
\LANGen{
\Question{
Let \( ABCDEF \) be a regular hexagon with the centre \( S \) and the side of length \( 3\,\mathrm{cm}\).
The point \( G \) is the midpoint of the segment \( AB \).
The vectors  \( \overrightarrow{u} \), \( \overrightarrow{v} \), \( \overrightarrow{w} \), \( \overrightarrow{z} \) are indicated in the hexagon shown in the picture.
Find the dot product of: \( \overrightarrow{v}\cdot\overrightarrow{w} \), \( \overrightarrow{v}\cdot\overrightarrow{z} \) and \( \overrightarrow{v}\cdot\overrightarrow{u} \).}
{\( \overrightarrow{v}\cdot\overrightarrow{w}=9 \), \( \overrightarrow{v}\cdot\overrightarrow{z} = 0 \), \( \overrightarrow{v}\cdot\overrightarrow{u}=27 \)}{\( \overrightarrow{v}\cdot\overrightarrow{w}=9 \), \( \overrightarrow{v}\cdot\overrightarrow{z} = 0 \), \( \overrightarrow{v}\cdot\overrightarrow{u}=9\sqrt6 \)}{\( \overrightarrow{v}\cdot\overrightarrow{w}=\frac92 \), \( \overrightarrow{v}\cdot\overrightarrow{z} = 0 \), \( \overrightarrow{v}\cdot\overrightarrow{u}=9\sqrt6 \)}{\( \overrightarrow{v}\cdot\overrightarrow{w}=\frac92 \), \( \overrightarrow{v}\cdot\overrightarrow{z} = 1 \), \( \overrightarrow{v}\cdot\overrightarrow{u}=27 \)}{}{}}
\LANGpl{
\Question{
Dany jest prawidłowy sześciokąt \( ABCDEF \) o środku \( S \) i boku równym \( 3\,\mathrm{cm}\).
Punkt \( G \) to środek odcinka \( AB \).
Wektory w sześciokącie  \( \overrightarrow{u} \), \( \overrightarrow{v} \), \( \overrightarrow{w} \), \( \overrightarrow{z} \) wskazano na rysunku.
Oblicz iloczyn skalarny: \( \overrightarrow{v}\cdot\overrightarrow{w} \), \( \overrightarrow{v}\cdot\overrightarrow{z} \) and \( \overrightarrow{v}\cdot\overrightarrow{u} \).}
{\( \overrightarrow{v}\cdot\overrightarrow{w}=9 \), \( \overrightarrow{v}\cdot\overrightarrow{z} = 0 \), \( \overrightarrow{v}\cdot\overrightarrow{u}=27 \)}{\( \overrightarrow{v}\cdot\overrightarrow{w}=9 \), \( \overrightarrow{v}\cdot\overrightarrow{z} = 0 \), \( \overrightarrow{v}\cdot\overrightarrow{u}=9\sqrt6 \)}{\( \overrightarrow{v}\cdot\overrightarrow{w}=\frac92 \), \( \overrightarrow{v}\cdot\overrightarrow{z} = 0 \), \( \overrightarrow{v}\cdot\overrightarrow{u}=9\sqrt6 \)}{\( \overrightarrow{v}\cdot\overrightarrow{w}=\frac92 \), \( \overrightarrow{v}\cdot\overrightarrow{z} = 1 \), \( \overrightarrow{v}\cdot\overrightarrow{u}=27 \)}{}{}}
\LANGcs{
\Question{
Je dán pravidelný šestiúhelník \( ABCDEF \) se středem \( S \) a délkou strany \( 3\,\mathrm{cm}\).
Bod \( G \) je středem strany \( AB \).
V šestiúhelníku jsou vyznačeny vektory \( \overrightarrow{u} \), \( \overrightarrow{v} \), \( \overrightarrow{w} \), \( \overrightarrow{z} \).
Vypočtěte skalární součiny: \( \overrightarrow{v}\cdot\overrightarrow{w} \), \( \overrightarrow{v}\cdot\overrightarrow{z} \) a \( \overrightarrow{v}\cdot\overrightarrow{u} \).}
{\( \overrightarrow{v}\cdot\overrightarrow{w}=9 \), \( \overrightarrow{v}\cdot\overrightarrow{z} = 0 \), \( \overrightarrow{v}\cdot\overrightarrow{u}=27 \)}{\( \overrightarrow{v}\cdot\overrightarrow{w}=9 \), \( \overrightarrow{v}\cdot\overrightarrow{z} = 0 \), \( \overrightarrow{v}\cdot\overrightarrow{u}=9\sqrt6 \)}{\( \overrightarrow{v}\cdot\overrightarrow{w}=\frac92 \), \( \overrightarrow{v}\cdot\overrightarrow{z} = 0 \), \( \overrightarrow{v}\cdot\overrightarrow{u}=9\sqrt6 \)}{\( \overrightarrow{v}\cdot\overrightarrow{w}=\frac92 \), \( \overrightarrow{v}\cdot\overrightarrow{z} = 1 \), \( \overrightarrow{v}\cdot\overrightarrow{u}=27 \)}{}{}}
\LANGsk{
\Question{
Je daný pravidelný šesťuholník \( ABCDEF \) so stredom \( S \) a dĺžkou strany \( 3\,\mathrm{cm}\).
Bod \( G \) je stredom strany \( AB \).
V šesťuholníku sú vyznačené vektory \( \overrightarrow{u} \), \( \overrightarrow{v} \), \( \overrightarrow{w} \), \( \overrightarrow{z} \).
Vypočítajte skalárne súčiny: \( \overrightarrow{v}\cdot\overrightarrow{w} \), \( \overrightarrow{v}\cdot\overrightarrow{z} \) a \( \overrightarrow{v}\cdot\overrightarrow{u} \).}
{\( \overrightarrow{v}\cdot\overrightarrow{w}=9 \), \( \overrightarrow{v}\cdot\overrightarrow{z} = 0 \), \( \overrightarrow{v}\cdot\overrightarrow{u}=27 \)}{\( \overrightarrow{v}\cdot\overrightarrow{w}=9 \), \( \overrightarrow{v}\cdot\overrightarrow{z} = 0 \), \( \overrightarrow{v}\cdot\overrightarrow{u}=9\sqrt6 \)}{\( \overrightarrow{v}\cdot\overrightarrow{w}=\frac92 \), \( \overrightarrow{v}\cdot\overrightarrow{z} = 0 \), \( \overrightarrow{v}\cdot\overrightarrow{u}=9\sqrt6 \)}{\( \overrightarrow{v}\cdot\overrightarrow{w}=\frac92 \), \( \overrightarrow{v}\cdot\overrightarrow{z} = 1 \), \( \overrightarrow{v}\cdot\overrightarrow{u}=27 \)}{}{}}
\LANGes{
\Question{
Sea \( ABCDEF \)  un hexágono regular con el centro \( S \) y la longitud del lado \( 3\,\mathrm{cm}\).
El punto \( G \) es el centro del segmento \( AB \).
Los vectores  \( \overrightarrow{u} \), \( \overrightarrow{v} \), \( \overrightarrow{w} \), \( \overrightarrow{z} \) son marcados en el hexágono de la imagen.
Determina el producto escalar: \( \overrightarrow{v}\cdot\overrightarrow{w} \), \( \overrightarrow{v}\cdot\overrightarrow{z} \) and \( \overrightarrow{v}\cdot\overrightarrow{u} \).}
{\( \overrightarrow{v}\cdot\overrightarrow{w}=9 \), \( \overrightarrow{v}\cdot\overrightarrow{z} = 0 \), \( \overrightarrow{v}\cdot\overrightarrow{u}=27 \)}{\( \overrightarrow{v}\cdot\overrightarrow{w}=9 \), \( \overrightarrow{v}\cdot\overrightarrow{z} = 0 \), \( \overrightarrow{v}\cdot\overrightarrow{u}=9\sqrt6 \)}{\( \overrightarrow{v}\cdot\overrightarrow{w}=\frac92 \), \( \overrightarrow{v}\cdot\overrightarrow{z} = 0 \), \( \overrightarrow{v}\cdot\overrightarrow{u}=9\sqrt6 \)}{\( \overrightarrow{v}\cdot\overrightarrow{w}=\frac92 \), \( \overrightarrow{v}\cdot\overrightarrow{z} = 1 \), \( \overrightarrow{v}\cdot\overrightarrow{u}=27 \)}{}{}}
\End

\Data\ID{1103030501}
\Updated{2020-07-15 03:07:12}
\Level{B}
\SubArea{Points and vectors}
\IMG{1103030501.pdf}
\LANGen{
\Question{
The vectors \( \overrightarrow{u} \), \( \overrightarrow{v}\), \( \overrightarrow{w} \), \( \overrightarrow{z} \) are indicated in a cube shown in the figure. The cube  edge length is \( 1 \). Find the dot products of:
\[ \overrightarrow{v}\cdot\overrightarrow{z}\text{ ,}\ \ \overrightarrow{u}\cdot\overrightarrow{v} \text{ ,}\ \ \overrightarrow{w}\cdot\overrightarrow{u}\]}
{\( \overrightarrow{v}\cdot\overrightarrow{z}=1 \), \( \overrightarrow{u}\cdot\overrightarrow{v}=0 \), \( \overrightarrow{w}\cdot\overrightarrow{u}=1 \)}{\( \overrightarrow{v}\cdot\overrightarrow{z}=\frac{\sqrt2}2 \), \( \overrightarrow{u}\cdot\overrightarrow{v}=1 \), \( \overrightarrow{w}\cdot\overrightarrow{u}=\sqrt3 \)}{\( \overrightarrow{v}\cdot\overrightarrow{z}=\sqrt2 \), \( \overrightarrow{u}\cdot\overrightarrow{v}=0 \), \( \overrightarrow{w}\cdot\overrightarrow{u}=1 \)}{\( \overrightarrow{v}\cdot\overrightarrow{z}=1 \), \( \overrightarrow{u}\cdot\overrightarrow{v}=1 \), \( \overrightarrow{w}\cdot\overrightarrow{u}=\sqrt3 \)}{}{}}
\LANGpl{
\Question{
W sześcianie wskazano wektory \( \overrightarrow{u} \), \( \overrightarrow{v}\), \( \overrightarrow{w} \), \( \overrightarrow{z} \). Krawędź sześcianu jest równa \( 1 \). Oblicz iloczyn skalarny:
\[ \overrightarrow{v}\cdot\overrightarrow{z}\text{ ,}\ \ \overrightarrow{u}\cdot\overrightarrow{v} \text{ ,}\ \ \overrightarrow{w}\cdot\overrightarrow{u} \]}
{\( \overrightarrow{v}\cdot\overrightarrow{z}=1 \), \( \overrightarrow{u}\cdot\overrightarrow{v}=0 \), \( \overrightarrow{w}\cdot\overrightarrow{u}=1 \)}{\( \overrightarrow{v}\cdot\overrightarrow{z}=\frac{\sqrt2}2 \), \( \overrightarrow{u}\cdot\overrightarrow{v}=1 \), \( \overrightarrow{w}\cdot\overrightarrow{u}=\sqrt3 \)}{\( \overrightarrow{v}\cdot\overrightarrow{z}=\sqrt2 \), \( \overrightarrow{u}\cdot\overrightarrow{v}=0 \), \( \overrightarrow{w}\cdot\overrightarrow{u}=1 \)}{\( \overrightarrow{v}\cdot\overrightarrow{z}=1 \), \( \overrightarrow{u}\cdot\overrightarrow{v}=1 \), \( \overrightarrow{w}\cdot\overrightarrow{u}=\sqrt3 \)}{}{}}
\LANGcs{
\Question{
V krychli o délce hrany \( 1 \) jsou vyznačeny vektory \( \overrightarrow{u} \), \( \overrightarrow{v}\), \( \overrightarrow{w} \), \( \overrightarrow{z} \). Vypočtěte skalární součiny:
\[ \overrightarrow{v}\cdot\overrightarrow{z}\text{ ,}\ \ \overrightarrow{u}\cdot\overrightarrow{v} \text{ ,}\ \ \overrightarrow{w}\cdot\overrightarrow{u} \]}
{\( \overrightarrow{v}\cdot\overrightarrow{z}=1 \), \( \overrightarrow{u}\cdot\overrightarrow{v}=0 \), \( \overrightarrow{w}\cdot\overrightarrow{u}=1 \)}{\( \overrightarrow{v}\cdot\overrightarrow{z}=\frac{\sqrt2}2 \), \( \overrightarrow{u}\cdot\overrightarrow{v}=1 \), \( \overrightarrow{w}\cdot\overrightarrow{u}=\sqrt3 \)}{\( \overrightarrow{v}\cdot\overrightarrow{z}=\sqrt2 \), \( \overrightarrow{u}\cdot\overrightarrow{v}=0 \), \( \overrightarrow{w}\cdot\overrightarrow{u}=1 \)}{\( \overrightarrow{v}\cdot\overrightarrow{z}=1 \), \( \overrightarrow{u}\cdot\overrightarrow{v}=1 \), \( \overrightarrow{w}\cdot\overrightarrow{u}=\sqrt3 \)}{}{}}
\LANGsk{
\Question{
V kocke s dĺžkou hrany \( 1 \) sú vyznačené vektory \( \overrightarrow{u} \), \( \overrightarrow{v}\), \( \overrightarrow{w} \), \( \overrightarrow{z} \). Vypočítajte skalárne súčiny:
\[ \overrightarrow{v}\cdot\overrightarrow{z}\text{ ,}\ \ \overrightarrow{u}\cdot\overrightarrow{v} \text{ ,}\ \ \overrightarrow{w}\cdot\overrightarrow{u}\]}
{\( \overrightarrow{v}\cdot\overrightarrow{z}=1 \), \( \overrightarrow{u}\cdot\overrightarrow{v}=0 \), \( \overrightarrow{w}\cdot\overrightarrow{u}=1 \)}{\( \overrightarrow{v}\cdot\overrightarrow{z}=\frac{\sqrt2}2 \), \( \overrightarrow{u}\cdot\overrightarrow{v}=1 \), \( \overrightarrow{w}\cdot\overrightarrow{u}=\sqrt3 \)}{\( \overrightarrow{v}\cdot\overrightarrow{z}=\sqrt2 \), \( \overrightarrow{u}\cdot\overrightarrow{v}=0 \), \( \overrightarrow{w}\cdot\overrightarrow{u}=1 \)}{\( \overrightarrow{v}\cdot\overrightarrow{z}=1 \), \( \overrightarrow{u}\cdot\overrightarrow{v}=1 \), \( \overrightarrow{w}\cdot\overrightarrow{u}=\sqrt3 \)}{}{}}
\LANGes{
\Question{
Los vectores \( \overrightarrow{u} \), \( \overrightarrow{v}\), \( \overrightarrow{w} \), \( \overrightarrow{z} \) se indican en un cubo mostrado en la figura. La longitud de la arista del cubo es \( 1 \). Determina el producto escalar de:
\[ \overrightarrow{v}\cdot\overrightarrow{z}\text{ ,}\ \ \overrightarrow{u}\cdot\overrightarrow{v} \text{ ,}\ \ \overrightarrow{w}\cdot\overrightarrow{u}\]}
{\( \overrightarrow{v}\cdot\overrightarrow{z}=1 \), \( \overrightarrow{u}\cdot\overrightarrow{v}=0 \), \( \overrightarrow{w}\cdot\overrightarrow{u}=1 \)}{\( \overrightarrow{v}\cdot\overrightarrow{z}=\frac{\sqrt2}2 \), \( \overrightarrow{u}\cdot\overrightarrow{v}=1 \), \( \overrightarrow{w}\cdot\overrightarrow{u}=\sqrt3 \)}{\( \overrightarrow{v}\cdot\overrightarrow{z}=\sqrt2 \), \( \overrightarrow{u}\cdot\overrightarrow{v}=0 \), \( \overrightarrow{w}\cdot\overrightarrow{u}=1 \)}{\( \overrightarrow{v}\cdot\overrightarrow{z}=1 \), \( \overrightarrow{u}\cdot\overrightarrow{v}=1 \), \( \overrightarrow{w}\cdot\overrightarrow{u}=\sqrt3 \)}{}{}}
\End

\Data\ID{1103030502}
\Updated{2020-07-15 03:07:12}
\Level{B}
\SubArea{Points and vectors}
\IMG{1103030502.pdf}
\LANGen{
\Question{
Find the coordinates of the vectors \( \overrightarrow{u} \) and \( \overrightarrow{v} \)  given by the picture and evaluate their dot product.}
{\( \overrightarrow{u}=(-3;6);\ \ \overrightarrow{v} =(-9;-6);\ \  \overrightarrow{u}\cdot\overrightarrow{v} = -9 \)}{\( \overrightarrow{u}=(3;-6);\ \ \overrightarrow{v} =(9;6);\ \  \overrightarrow{u}\cdot\overrightarrow{v} = -9 \)}{\( \overrightarrow{u}=(-3;6);\ \ \overrightarrow{v} =(-9;-6);\ \  \overrightarrow{u}\cdot\overrightarrow{v} = 9 \)}{\( \overrightarrow{u}=(3;-6);\ \ \overrightarrow{v} =(9;6);\ \  \overrightarrow{u}\cdot\overrightarrow{v} = 0 \)}{}{}}
\LANGpl{
\Question{
Określ współrzędne wektorów \( \overrightarrow{u} \) i \( \overrightarrow{v} \)  podanych na rysunku i oblicz iloczyn skalarny.}
{\( \overrightarrow{u}=(-3;6);\ \ \overrightarrow{v} =(-9;-6);\ \  \overrightarrow{u}\cdot\overrightarrow{v} = -9 \)}{\( \overrightarrow{u}=(3;-6);\ \ \overrightarrow{v} =(9;6);\ \  \overrightarrow{u}\cdot\overrightarrow{v} = -9 \)}{\( \overrightarrow{u}=(-3;6);\ \ \overrightarrow{v} =(-9;-6);\ \  \overrightarrow{u}\cdot\overrightarrow{v} = 9 \)}{\( \overrightarrow{u}=(3;-6);\ \ \overrightarrow{v} =(9;6);\ \  \overrightarrow{u}\cdot\overrightarrow{v} = 0 \)}{}{}}
\LANGcs{
\Question{
V souřadném systému jsou dány vektory \( \overrightarrow{u} \) a \( \overrightarrow{v} \). Určete jejich souřadnice a vypočtěte jejich skalární součin.}
{\( \overrightarrow{u}=(-3;6);\ \ \overrightarrow{v} =(-9;-6);\ \  \overrightarrow{u}\cdot\overrightarrow{v} = -9 \)}{\( \overrightarrow{u}=(3;-6);\ \ \overrightarrow{v} =(9;6);\ \  \overrightarrow{u}\cdot\overrightarrow{v} = -9 \)}{\( \overrightarrow{u}=(-3;6);\ \ \overrightarrow{v} =(-9;-6);\ \  \overrightarrow{u}\cdot\overrightarrow{v} = 9 \)}{\( \overrightarrow{u}=(3;-6);\ \ \overrightarrow{v} =(9;6);\ \  \overrightarrow{u}\cdot\overrightarrow{v} = 0 \)}{}{}}
\LANGsk{
\Question{
V súradnicovom systéme sú dané vektory \( \overrightarrow{u} \) a \( \overrightarrow{v} \). Určte ich súradnice a vypočítajte ich skalárny súčin.}
{\( \overrightarrow{u}=(-3;6);\ \ \overrightarrow{v} =(-9;-6);\ \  \overrightarrow{u}\cdot\overrightarrow{v} = -9 \)}{\( \overrightarrow{u}=(3;-6);\ \ \overrightarrow{v} =(9;6);\ \  \overrightarrow{u}\cdot\overrightarrow{v} = -9 \)}{\( \overrightarrow{u}=(-3;6);\ \ \overrightarrow{v} =(-9;-6);\ \  \overrightarrow{u}\cdot\overrightarrow{v} = 9 \)}{\( \overrightarrow{u}=(3;-6);\ \ \overrightarrow{v} =(9;6);\ \  \overrightarrow{u}\cdot\overrightarrow{v} = 0 \)}{}{}}
\LANGes{
\Question{
Determina las coordenadas de los vectores \( \overrightarrow{u} \) y \( \overrightarrow{v} \)  dadas por la imagen y calcula su producto escalar.}
{\( \overrightarrow{u}=(-3;6);\ \ \overrightarrow{v} =(-9;-6);\ \  \overrightarrow{u}\cdot\overrightarrow{v} = -9 \)}{\( \overrightarrow{u}=(3;-6);\ \ \overrightarrow{v} =(9;6);\ \  \overrightarrow{u}\cdot\overrightarrow{v} = -9 \)}{\( \overrightarrow{u}=(-3;6);\ \ \overrightarrow{v} =(-9;-6);\ \  \overrightarrow{u}\cdot\overrightarrow{v} = 9 \)}{\( \overrightarrow{u}=(3;-6);\ \ \overrightarrow{v} =(9;6);\ \  \overrightarrow{u}\cdot\overrightarrow{v} = 0 \)}{}{}}
\End

\Data\ID{1103030503}
\Updated{2020-07-15 03:07:12}
\Level{B}
\SubArea{Points and vectors}
\IMG{1103030503.pdf}
\LANGen{
\Question{
Find the coordinates of the vectors \(  \overrightarrow{u} \) and \(  \overrightarrow{v} \) given by the picture and evaluate their dot product.}
{\( \overrightarrow{u}=(-8;-7;9);\ \ \overrightarrow{v} =(8;7;9);\ \  \overrightarrow{u}\cdot\overrightarrow{v} = -32 \)}{\( \overrightarrow{u}=(-8;-7;9);\ \ \overrightarrow{v} =(8;7;9);\ \  \overrightarrow{u}\cdot\overrightarrow{v} = 0 \)}{\( \overrightarrow{u}=(-8;-7;9);\ \ \overrightarrow{v} =(8;7;9);\ \  \overrightarrow{u}\cdot\overrightarrow{v} = (-64;-49;81) \)}{\( \overrightarrow{u}=(8;7;-9);\ \ \overrightarrow{v} =(-8;-7;-9);\ \  \overrightarrow{u}\cdot\overrightarrow{v} = (-64;-49;81) \)}{}{}}
\LANGpl{
\Question{
Określ współrzędne wektorów \(  \overrightarrow{u} \) i \(  \overrightarrow{v} \) podanych na rysunku i oblicz iloczyn skalarny.}
{\( \overrightarrow{u}=(-8;-7;9);\ \ \overrightarrow{v} =(8;7;9);\ \  \overrightarrow{u}\cdot\overrightarrow{v} = -32 \)}{\( \overrightarrow{u}=(-8;-7;9);\ \ \overrightarrow{v} =(8;7;9);\ \  \overrightarrow{u}\cdot\overrightarrow{v} = 0 \)}{\( \overrightarrow{u}=(-8;-7;9);\ \ \overrightarrow{v} =(8;7;9);\ \  \overrightarrow{u}\cdot\overrightarrow{v} = (-64;-49;81) \)}{\( \overrightarrow{u}=(8;7;-9);\ \ \overrightarrow{v} =(-8;-7;-9);\ \  \overrightarrow{u}\cdot\overrightarrow{v} = (-64;-49;81) \)}{}{}}
\LANGcs{
\Question{
Vektory \(  \overrightarrow{u} \) a \(  \overrightarrow{v} \) jsou znázorněny v souřadném systému. Určete jejich souřadnice a vypočtěte jejich skalární součin.}
{\( \overrightarrow{u}=(-8;-7;9);\ \ \overrightarrow{v} =(8;7;9);\ \  \overrightarrow{u}\cdot\overrightarrow{v} = -32 \)}{\( \overrightarrow{u}=(-8;-7;9);\ \ \overrightarrow{v} =(8;7;9);\ \  \overrightarrow{u}\cdot\overrightarrow{v} = 0 \)}{\( \overrightarrow{u}=(-8;-7;9);\ \ \overrightarrow{v} =(8;7;9);\ \  \overrightarrow{u}\cdot\overrightarrow{v} = (-64;-49;81) \)}{\( \overrightarrow{u}=(8;7;-9);\ \ \overrightarrow{v} =(-8;-7;-9);\ \  \overrightarrow{u}\cdot\overrightarrow{v} = (-64;-49;81) \)}{}{}}
\LANGsk{
\Question{
Vektory \(  \overrightarrow{u} \) a \(  \overrightarrow{v} \) sú znázornené v súradnicovom systému. Určte ich súradnice a vypočítajte ich skalárny súčin.}
{\( \overrightarrow{u}=(-8;-7;9);\ \ \overrightarrow{v} =(8;7;9);\ \  \overrightarrow{u}\cdot\overrightarrow{v} = -32 \)}{\( \overrightarrow{u}=(-8;-7;9);\ \ \overrightarrow{v} =(8;7;9);\ \  \overrightarrow{u}\cdot\overrightarrow{v} = 0 \)}{\( \overrightarrow{u}=(-8;-7;9);\ \ \overrightarrow{v} =(8;7;9);\ \  \overrightarrow{u}\cdot\overrightarrow{v} = (-64;-49;81) \)}{\( \overrightarrow{u}=(8;7;-9);\ \ \overrightarrow{v} =(-8;-7;-9);\ \  \overrightarrow{u}\cdot\overrightarrow{v} = (-64;-49;81) \)}{}{}}
\LANGes{
\Question{
Determina las coordenadas de los vectores \(  \overrightarrow{u} \) y \(  \overrightarrow{v} \) dadas por la imagen y calcula su producto escalar.}
{\( \overrightarrow{u}=(-8;-7;9);\ \ \overrightarrow{v} =(8;7;9);\ \  \overrightarrow{u}\cdot\overrightarrow{v} = -32 \)}{\( \overrightarrow{u}=(-8;-7;9);\ \ \overrightarrow{v} =(8;7;9);\ \  \overrightarrow{u}\cdot\overrightarrow{v} = 0 \)}{\( \overrightarrow{u}=(-8;-7;9);\ \ \overrightarrow{v} =(8;7;9);\ \  \overrightarrow{u}\cdot\overrightarrow{v} = (-64;-49;81) \)}{\( \overrightarrow{u}=(8;7;-9);\ \ \overrightarrow{v} =(-8;-7;-9);\ \  \overrightarrow{u}\cdot\overrightarrow{v} = (-64;-49;81) \)}{}{}}
\End

\Data\ID{1103030504}
\Updated{2020-07-15 03:07:12}
\Level{B}
\SubArea{Points and vectors}
\IMG{1103030504.pdf}
\LANGen{
\Question{
The vectors \( \overrightarrow{u} \)  and \( \overrightarrow{v} \) are given by the figure. Find cosine of the angle \(\varphi \) between \( \overrightarrow{u} \) and \( \overrightarrow{v} \).
Help: Use the dot product of the given vectors.}
{\( \cos\varphi=\frac{13\sqrt{10}}{50} \)}{\( \cos\varphi=\frac{970}{50} \)}{\( \cos\varphi=\frac{3\sqrt{10}}{10} \)}{\( \cos\varphi=\frac{\sqrt{10}}{5} \)}{}{}}
\LANGpl{
\Question{
W układzie współrzędnych podano wektory \( \overrightarrow{u} \)  i \( \overrightarrow{v} \). Oblicz cosinus kąta \(\varphi \) między \( \overrightarrow{u} \) i \( \overrightarrow{v} \).
Wskazówka: Użyj iloczynu skalarnego wektorów.}
{\( \cos\varphi=\frac{13\sqrt{10}}{50} \)}{\( \cos\varphi=\frac{970}{50} \)}{\( \cos\varphi=\frac{3\sqrt{10}}{10} \)}{\( \cos\varphi=\frac{\sqrt{10}}{5} \)}{}{}}
\LANGcs{
\Question{
V souřadném systému jsou znázorněny vektory \( \overrightarrow{u} \)  a \( \overrightarrow{v} \). Určete kosinus jejich odchylky \(\varphi \).
Nápověda: Užijte skalární součin vektorů.}
{\( \cos\varphi=\frac{13\sqrt{10}}{50} \)}{\( \cos\varphi=\frac{970}{50} \)}{\( \cos\varphi=\frac{3\sqrt{10}}{10} \)}{\( \cos\varphi=\frac{\sqrt{10}}{5} \)}{}{}}
\LANGsk{
\Question{
V súradnicovom systéme sú znázornené vektory \( \overrightarrow{u} \)  a \( \overrightarrow{v} \). Určte kosínus ich odchýlky \(\varphi \).
Nápoveda: Použite skalárny súčin vektorov.}
{\( \cos\varphi=\frac{13\sqrt{10}}{50} \)}{\( \cos\varphi=\frac{970}{50} \)}{\( \cos\varphi=\frac{3\sqrt{10}}{10} \)}{\( \cos\varphi=\frac{\sqrt{10}}{5} \)}{}{}}
\LANGes{
\Question{
Los vectores \( \overrightarrow{u} \)  y \( \overrightarrow{v} \) son mostrados en la imagen. Determina el coseno del ángulo \(\varphi \) entre \( \overrightarrow{u} \) y \( \overrightarrow{v} \).
Pista: Usa el producto escalar de los vectores indicados.}
{\( \cos\varphi=\frac{13\sqrt{10}}{50} \)}{\( \cos\varphi=\frac{970}{50} \)}{\( \cos\varphi=\frac{3\sqrt{10}}{10} \)}{\( \cos\varphi=\frac{\sqrt{10}}{5} \)}{}{}}
\End

\Data\ID{1103030505}
\Updated{2020-07-15 03:07:12}
\Level{B}
\SubArea{Points and vectors}
\IMG{1103030505.pdf}
\LANGen{
\Question{
The vectors \( \overrightarrow{u} \) and \( \overrightarrow{v} \) are given by the figure. Find cosine of the angle \( \varphi \) between \( \overrightarrow{u} \) and \( \overrightarrow{v} \).
Help: Use the dot product of the given vectors.}
{\( \cos\varphi=-\frac9{17} \)}{\( \cos\varphi=\frac9{17} \)}{\( \cos\varphi=\frac{\sqrt{17}}{2\sqrt{13}} \)}{\( \cos\varphi=-\frac{\sqrt{17}}{2\sqrt{13}} \)}{}{}}
\LANGpl{
\Question{
W układzie współrzędnych podano wektory \( \overrightarrow{u} \) i \( \overrightarrow{v} \). Oblicz cosinus kąta \( \varphi \) między \( \overrightarrow{u} \) i \( \overrightarrow{v} \).
Wskazówka: Użyj iloczynu skalarnego wektorów.}
{\( \cos\varphi=-\frac9{17} \)}{\( \cos\varphi=\frac9{17} \)}{\( \cos\varphi=\frac{\sqrt{17}}{2\sqrt{13}} \)}{\( \cos\varphi=-\frac{\sqrt{17}}{2\sqrt{13}} \)}{}{}}
\LANGcs{
\Question{
Vektory \( \overrightarrow{u} \) a \( \overrightarrow{v} \) jsou zadány v souřadném systému. Určete kosinus jejich odchylky \( \varphi \).
Nápověda: Užijte skalární součin daných vektorů.}
{\( \cos\varphi=-\frac9{17} \)}{\( \cos\varphi=\frac9{17} \)}{\( \cos\varphi=\frac{\sqrt{17}}{2\sqrt{13}} \)}{\( \cos\varphi=-\frac{\sqrt{17}}{2\sqrt{13}} \)}{}{}}
\LANGsk{
\Question{
Vektory \( \overrightarrow{u} \) a \( \overrightarrow{v} \) sú zadané v súradnicovom systéme. Určte kosínus ich odchýlky \( \varphi \).
Nápoveda: Použite skalárny súčin daných vektorov.}
{\( \cos\varphi=-\frac9{17} \)}{\( \cos\varphi=\frac9{17} \)}{\( \cos\varphi=\frac{\sqrt{17}}{2\sqrt{13}} \)}{\( \cos\varphi=-\frac{\sqrt{17}}{2\sqrt{13}} \)}{}{}}
\LANGes{
\Question{
Los vectores \( \overrightarrow{u} \) y \( \overrightarrow{v} \) son mostrados en la imagen. Determina el coseno del ángulo \( \varphi \) entre \( \overrightarrow{u} \) y \( \overrightarrow{v} \).
Pista: Usa el producto escalar de los vectores indicados.}
{\( \cos\varphi=-\frac9{17} \)}{\( \cos\varphi=\frac9{17} \)}{\( \cos\varphi=\frac{\sqrt{17}}{2\sqrt{13}} \)}{\( \cos\varphi=-\frac{\sqrt{17}}{2\sqrt{13}} \)}{}{}}
\End

\Data\ID{1103030601}
\Updated{2020-07-15 03:07:12}
\Level{B}
\SubArea{Points and vectors}
\IMG{1103030601.pdf}
\LANGen{
\Question{
In the cube \( ABCDEFGH \) find the angle \( \varphi \) between the vectors \( \overrightarrow{b}=\overrightarrow{EB} \) and  \( \overrightarrow{a}=\overrightarrow{AK} \), where \( K \) is the midpoint of \( HG \). Round \( \varphi \) to the nearest degree.
Help: Choose the appropriate coordinate system.}
{\( \varphi\,\dot{=}\,104^{\circ} \)}{\( \varphi\,\dot{=}\,76^{\circ} \)}{\( \varphi\,\dot{=}\,100^{\circ} \)}{\( \varphi\,\dot{=}\,80^{\circ} \)}{}{}}
\LANGpl{
\Question{
W sześcianie \( ABCDEFGH \) wyznacz kąt \( \varphi \) między wektorami \( \overrightarrow{b}=\overrightarrow{EB} \) i  \( \overrightarrow{a}=\overrightarrow{AK} \), gdzie \( K \) to środek \( HG \).  Zaokrągli  \( \varphi \) do pełnych stopni.
Wskazówka: Wybierz odpowiedni układ współrzędnych.}
{\( \varphi\,\dot{=}\,104^{\circ} \)}{\( \varphi\,\dot{=}\,76^{\circ} \)}{\( \varphi\,\dot{=}\,100^{\circ} \)}{\( \varphi\,\dot{=}\,80^{\circ} \)}{}{}}
\LANGcs{
\Question{
V krychli \( ABCDEFGH \) určete odchylku \( \varphi \) vektorů \( \overrightarrow{b}=\overrightarrow{EB} \) a  \( \overrightarrow{a}=\overrightarrow{AK} \), kde \( K \) je střed \( HG \). Zaokrouhlete hodnotu \( \varphi \) na celé stupně.
Nápověda: Řešte ve vhodně zvoleném souřadném systému.}
{\( \varphi\,\dot{=}\,104^{\circ} \)}{\( \varphi\,\dot{=}\,76^{\circ} \)}{\( \varphi\,\dot{=}\,100^{\circ} \)}{\( \varphi\,\dot{=}\,80^{\circ} \)}{}{}}
\LANGsk{
\Question{
V kocke \( ABCDEFGH \) určte odchýlku \( \varphi \) vektorov \( \overrightarrow{b}=\overrightarrow{EB} \) a  \( \overrightarrow{a}=\overrightarrow{AK} \), kde \( K \) je stred \( HG \). Zaokrúhlite hodnotu \( \varphi \) na celé stupne.
Nápoveda: Riešte vo vhodne zvolenom súradnicovom systéme.}
{\( \varphi\,\dot{=}\,104^{\circ} \)}{\( \varphi\,\dot{=}\,76^{\circ} \)}{\( \varphi\,\dot{=}\,100^{\circ} \)}{\( \varphi\,\dot{=}\,80^{\circ} \)}{}{}}
\LANGes{
\Question{
En el cubo \( ABCDEFGH \) determina el ángulo \( \varphi \) entre los vectores \( \overrightarrow{b}=\overrightarrow{EB} \) y  \( \overrightarrow{a}=\overrightarrow{AK} \), donde \( K \) es el centro de \( HG \). Redondea \( \varphi \) al grado más cercano.
Pista: Elige un sistema de coordenadas apropiado.}
{\( \varphi\,\dot{=}\,104^{\circ} \)}{\( \varphi\,\dot{=}\,76^{\circ} \)}{\( \varphi\,\dot{=}\,100^{\circ} \)}{\( \varphi\,\dot{=}\,80^{\circ} \)}{}{}}
\End

\Data\ID{1103030602}
\Updated{2020-07-15 03:07:12}
\Level{B}
\SubArea{Points and vectors}
\IMG{1103030602_0.pdf}
\LANGen{
\Question{
Let  \( ABCDV \) be a right pyramid with a square base, such that its opposite edges contain  a right angle (see the picture).
Specify the missing coordinate of the  apex \( V \).}
{\( m=2\sqrt2 \)}{\( m=-2\sqrt2 \)}{\( m=4\sqrt2 \)}{\( m=\sqrt2 \)}{}{}}
\LANGpl{
\Question{
Dany jest ostrosłup prawidłowy \( ABCDV \) o podstawie kwadratu, jego  przeciwne krawędzie boczne mają kąt prosty (patrz rysunek). Określ brakującą współrzędną wierzchołka \( V \).}
{\( m=2\sqrt2 \)}{\( m=-2\sqrt2 \)}{\( m=4\sqrt2 \)}{\( m=\sqrt2 \)}{}{}}
\LANGcs{
\Question{
Je dán pravidelný čtyřboký jehlan  \( ABCDV \), jehož protilehlé boční hrany svírají pravý úhel (viz obrázek).
Určete chybějící souřadnici vrcholu \( V \).}
{\( m=2\sqrt2 \)}{\( m=-2\sqrt2 \)}{\( m=4\sqrt2 \)}{\( m=\sqrt2 \)}{}{}}
\LANGsk{
\Question{
Je daný pravidelný štvorboký ihlan  \( ABCDV \), ktorého protiľahlé bočné hrany zvierajú pravý uhol (viď obrázok).
Určte chýbajúcu súradnicu vrcholu \( V \).}
{\( m=2\sqrt2 \)}{\( m=-2\sqrt2 \)}{\( m=4\sqrt2 \)}{\( m=\sqrt2 \)}{}{}}
\LANGes{
\Question{
Sea \( ABCDV \) una pirámide cuadrangular regular, cuyas aristas opuestas forman un ángulo recto (mira la imagen). Especifica la coordenada del vértice faltante. \( V \).}
{\( m=2\sqrt2 \)}{\( m=-2\sqrt2 \)}{\( m=4\sqrt2 \)}{\( m=\sqrt2 \)}{}{}}
\End

\Data\ID{1103040209}
\Updated{2020-07-15 03:07:12}
\Level{B}
\SubArea{Points and vectors}
\IMG{1103040209.pdf}
\LANGen{
\Question{
In the picture, there are indicated vectors $\vec{u}$ and $\vec{v}$ in three squares. Find the measure of an angle $\varphi$ between $\vec{u}$ and $\vec{v}$. Round $\varphi$ to the nearest degree.
Hint: Set up a coordinate system conveniently.}
{$\varphi\,\dot{=}\,8^{\circ}$}{$\varphi\,\dot{=}\,9^{\circ}$}{$\varphi\,\dot{=}\,10^{\circ}$}{$\varphi\,\dot{=}\,11^{\circ}$}{}{}}
\LANGpl{
\Question{
Na rysunku przedstawiono wektory  $\vec{u}$ and $\vec{v}$ w trzech kwadratach. Oblicz miarę kąta $\varphi$ pomiędzy $\vec{u}$ i $\vec{v}$. Zaokrągli $\varphi$ do pełnych stopni.
Wskazówka: utworzenie układu współrzędnych ułatwi zadanie.}
{$\varphi\,\dot{=}\,8^{\circ}$}{$\varphi\,\dot{=}\,9^{\circ}$}{$\varphi\,\dot{=}\,10^{\circ}$}{$\varphi\,\dot{=}\,11^{\circ}$}{}{}}
\LANGcs{
\Question{
V trojici čtverců na obrázku jsou vyznačeny vektory $\vec{u}$ a $\vec{v}$. Vypočtěte jejich odchylku $\varphi$ a zaokrouhlete ji na celé stupně.
Nápověda: Řešte ve vhodně zvoleném souřadném systému.}
{$\varphi\,\dot{=}\,8^{\circ}$}{$\varphi\,\dot{=}\,9^{\circ}$}{$\varphi\,\dot{=}\,10^{\circ}$}{$\varphi\,\dot{=}\,11^{\circ}$}{}{}}
\LANGsk{
\Question{
V trojici štvorcov na obrázku sú vyznačené vektory $\vec{u}$ a $\vec{v}$. Vypočítajte ich odchýlku $\varphi$ a zaokrúhlite ju na celé stupne.
Nápoveda: Riešte vo vhodne zvolenom súradnicovom systéme.}
{$\varphi\,\dot{=}\,8^{\circ}$}{$\varphi\,\dot{=}\,9^{\circ}$}{$\varphi\,\dot{=}\,10^{\circ}$}{$\varphi\,\dot{=}\,11^{\circ}$}{}{}}
\LANGes{
\Question{
En la imagen, se indican los vectores $\vec{u}$ y $\vec{v}$  en tres cuadrados. Calcula la desviación  $\varphi$ entre $\vec{u}$ y $\vec{v}$. Redondea $\varphi$ al grado más cercano.
Pista: Soluciona  en un sistema de coordenadas adecuadamente elegido.}
{$\varphi\,\dot{=}\,8^{\circ}$}{$\varphi\,\dot{=}\,9^{\circ}$}{$\varphi\,\dot{=}\,10^{\circ}$}{$\varphi\,\dot{=}\,11^{\circ}$}{}{}}
\End

\Data\ID{9000100706}
\Updated{2020-07-15 03:07:37}
\Level{B}
\SubArea{Points and vectors}
\LANGen{
\Question{
Given vectors \(\vec{a} = (-1;2;-3)\),
\(\vec{b} = (0;1;-1)\), find the
vector \(\vec{c}\) 
perpendicular to both of them.}
{\(\vec{c} = (-1;1;1)\)}{\(\vec{c} = (-3;0;1)\)}{\(\vec{c} = (2;4;2)\)}{\(\vec{c} = (-1;-1;1)\)}{}{}}
\LANGpl{
\Question{
Dane są wektory \(\vec{a} = (-1;2;-3)\),
\(\vec{b} = (0;1;-1)\), wskaż
wektor  \(\vec{c}\),
który jest prostopadły do obu wektorów.}
{\(\vec{c} = (-1;1;1)\)}{\(\vec{c} = (-3;0;1)\)}{\(\vec{c} = (2;4;2)\)}{\(\vec{c} = (-1;-1;1)\)}{}{}}
\LANGcs{
\Question{
Jsou dány vektory \(\vec{a} = (-1;2;-3)\),
\(\vec{b} = (0;1;-1)\). Vyberte
vektor \(\vec{c}\),
pro který platí, že je kolmý k oběma vektorům.}
{\(\vec{c} = (-1;1;1)\)}{\(\vec{c} = (-3;0;1)\)}{\(\vec{c} = (2;4;2)\)}{\(\vec{c} = (-1;-1;1)\)}{}{}}
\LANGsk{
\Question{
Sú dané vektory \(\vec{a} = (-1;2;-3)\),
\(\vec{b} = (0;1;-1)\). Vyberte
vektor \(\vec{c}\),
pre ktorý platí, že je kolmý k obom vektorom.}
{\(\vec{c} = (-1;1;1)\)}{\(\vec{c} = (-3;0;1)\)}{\(\vec{c} = (2;4;2)\)}{\(\vec{c} = (-1;-1;1)\)}{}{}}
\LANGes{
\Question{
Dados los vectores  \(\vec{a} = (-1;2;-3)\),
\(\vec{b} = (0;1;-1)\), halla el vector \(\vec{c}\) 
suponiendo que es perpendicular a  dos vectores dados.}
{\(\vec{c} = (-1;1;1)\)}{\(\vec{c} = (-3;0;1)\)}{\(\vec{c} = (2;4;2)\)}{\(\vec{c} = (-1;-1;1)\)}{}{}}
\End

\Data\ID{9000100707}
\Updated{2020-07-15 03:07:37}
\Level{B}
\SubArea{Points and vectors}
\LANGen{
\Question{
Consider points \(A = [-2;-1]\),
\(B = [1;y]\),
\(C = [3;-4]\). Find the coordinate
\(y\) which ensures
that the vectors \(\overrightarrow{AB } \) 
and \(\overrightarrow{AC } \) 
are perpendicular.}
{\(y = 4\)}{\(y = -4\)}{\(y = 0.8\)}{\(y = -0.8\)}{}{}}
\LANGpl{
\Question{
Dane są punkty \(A = [-2;-1]\),
\(B = [1;y]\),
\(C = [3;-4]\). Wskaż współrzędną 
\(y\) tak, aby
wektory \(\overrightarrow{AB } \) 
i \(\overrightarrow{AC } \) 
były prostopadłe.}
{\(y = 4\)}{\(y = -4\)}{\(y = 0.8\)}{\(y = -0.8\)}{}{}}
\LANGcs{
\Question{
V rovině jsou dány body \(A = [-2;-1]\),
\(B = [1;y_{B}]\),
\(C = [3;-4]\). Určete
souřadnici \(y_{B}\) tak,
aby platilo, že \(\overrightarrow{AB } \) 
\(\perp \) 
\(\overrightarrow{AC } \).}
{\(y_{B} = 4\)}{\(y_{B} = -4\)}{\(y_{B} = 0{,}8\)}{\(y_{B} = -0{,}8\)}{}{}}
\LANGsk{
\Question{
V rovine sú dané body \(A = [-2;-1]\),
\(B = [1;y_{B}]\),
\(C = [3;-4]\). Určte
súradnicu \(y_{B}\) tak,
aby platilo, že \(\overrightarrow{AB } \) 
\(\perp \) 
\(\overrightarrow{AC } \).}
{\(y_{B} = 4\)}{\(y_{B} = -4\)}{\(y_{B} = 0{,}8\)}{\(y_{B} = -0{,}8\)}{}{}}
\LANGes{
\Question{
Considera los puntos \(A = [-2;-1]\),
\(B = [1;y]\),
\(C = [3;-4]\). Halla la coordenada
\(y\) para que los vectores \(\overrightarrow{AB } \) 
y \(\overrightarrow{AC } \) 
sean perpendiculares.}
{\(y = 4\)}{\(y = -4\)}{\(y = 0.8\)}{\(y = -0.8\)}{}{}}
\End

\Data\ID{9000100708}
\Updated{2020-07-15 03:07:37}
\Level{B}
\SubArea{Points and vectors}
\LANGen{
\Question{
Given points \(A = [-2;-1]\),
\(B = [x;-3]\),
\(C = [4;-4]\), find the coordinate
\(x\) which ensures
that the vectors \(\overrightarrow{AB } \) 
and \(\overrightarrow{AC } \) 
are parallel.}
{\(x = 2\)}{\(x = 7\)}{\(x = -3\)}{\(x = 3\)}{}{}}
\LANGpl{
\Question{
Dane są punkty \(A = [-2;-1]\),
\(B = [x;-3]\),
\(C = [4;-4]\), wskaż współrzędną
\(x\) tak, aby
wektory \(\overrightarrow{AB } \) 
i \(\overrightarrow{AC } \) 
były równoległe.}
{\(x = 2\)}{\(x = 7\)}{\(x = -3\)}{\(x = 3\)}{}{}}
\LANGcs{
\Question{
V rovině jsou dány body \(A = [-2;-1]\),
\(B = [x_{B};-3]\),
\(C = [4;-4]\). Určete
souřadnici \(x_{B}\) tak,
aby platilo, že \(\overrightarrow{AB } ||\overrightarrow{AC } \).}
{\(x_{B} = 2\)}{\(x_{B} = 7\)}{\(x_{B} = -3\)}{\(x_{B} = 3\)}{}{}}
\LANGsk{
\Question{
V rovine sú dané body \(A = [-2;-1]\),
\(B = [x_{B};-3]\),
\(C = [4;-4]\). Určte
súradnicu \(x_{B}\) tak,
aby platilo, že \(\overrightarrow{AB } ||\overrightarrow{AC } \).}
{\(x_{B} = 2\)}{\(x_{B} = 7\)}{\(x_{B} = -3\)}{\(x_{B} = 3\)}{}{}}
\LANGes{
\Question{
Dados los puntos \(A = [-2;-1]\),
\(B = [x;-3]\),
\(C = [4;-4]\), halla la coordenada
\(x\) para que los vectores \(\overrightarrow{AB } \) 
y \(\overrightarrow{AC } \) 
sean paralelos.}
{\(x = 2\)}{\(x = 7\)}{\(x = -3\)}{\(x = 3\)}{}{}}
\End

\Data\ID{9000100709}
\Updated{2021-06-18 05:06:40}
\Level{B}
\SubArea{Points and vectors}
\LANGen{
\Question{
Find the value of the parameter \(z\) 
so that the vector \(\vec{w} = (8;2;z)\) 
is perpendicular to the vectors \(\vec{a} = (1;2;-3)\) 
and \(\vec{b} = (-1;2;1)\).}
{\(z = 4\)}{\(z = -12\)}{\(z = 2\)}{\(z = -4\)}{}{}}
\LANGpl{
\Question{
Wektor  \(\vec{w} = (8;2;z)\) 
jest prostopadły do wektorów  \(\vec{a} = (1;2;-3)\) 
i \(\vec{b} = (-1;2;1)\) jeśli:}
{\(z = 4\)}{\(z = -12\)}{\(z = 2\)}{\(z = -4\)}{}{}}
\LANGcs{
\Question{
Vektor \(\vec{w} = (8;2;z)\) je kolmý
na vektory \(\vec{a} = (1;2;-3),\ \vec{b} = (-1;2;1)\),
platí-li:}
{\(z = 4\)}{\(z = -12\)}{\(z = 2\)}{\(z = -4\)}{}{}}
\LANGsk{
\Question{
Vektor \(\vec{w} = (8;2;z)\) je kolmý
na vektory \(\vec{a} = (1;2;-3),\ \vec{b} = (-1;2;1)\),
ak platí:}
{\(z = 4\)}{\(z = -12\)}{\(z = 2\)}{\(z = -4\)}{}{}}
\LANGes{
\Question{
Halla el valor del parámetro \(z\) 
para que el vector \(\vec{w} = (8;2;z)\) 
sea perpendicular a los vectores \(\vec{a} = (1;2;-3)\) 
y \(\vec{b} = (-1;2;1)\).}
{\(z = 4\)}{\(z = -12\)}{\(z = 2\)}{\(z = -4\)}{}{}}
\End

\Data\ID{9000101802}
\Updated{2020-07-15 03:07:37}
\Level{B}
\SubArea{Points and vectors}
\LANGen{
\Question{
Among vectors \(\vec{u} = \left (- \frac{2}
{\sqrt{2}};2\sqrt{2}\right )\),
\(\vec{v} = (-5;10)\),
\(\vec{w} = (2.5;-5)\),
\(\vec{r} = (-3.5;6)\) 
find the vector which is not parallel to the vector
\(\vec{a} = (1;-2)\).}
{\(\vec{r}\)}{\(\vec{w}\)}{\(\vec{v}\)}{\(\vec{u}\)}{}{}}
\LANGpl{
\Question{
Dane są wektory \(\vec{u} = \left (- \frac{2}
{\sqrt{2}};2\sqrt{2}\right )\),
\(\vec{v} = (-5;10)\),
\(\vec{w} = (2.5;-5)\),
\(\vec{r} = (-3.5;6)\) 
wskaż wektor, który nie jest równoległy do wektora
\(\vec{a} = (1;-2)\).}
{\(\vec{r}\)}{\(\vec{w}\)}{\(\vec{v}\)}{\(\vec{u}\)}{}{}}
\LANGcs{
\Question{
Je dán vektor \(\vec{a} = (1;-2)\).
Který z vektorů \(\vec{u} = \left (- \frac{2}
{\sqrt{2}};2\sqrt{2}\right )\),
\(\vec{v} = (-5;10)\),
\(\vec{w} = (2{,}5;-5)\),
\(\vec{r} = (-3{,}5;6)\) není rovnoběžný
s vektorem \(\vec{a}\)?}
{\(\vec{r}\)}{\(\vec{w}\)}{\(\vec{v}\)}{\(\vec{u}\)}{}{}}
\LANGsk{
\Question{
Je daný vektor \(\vec{a} = (1;-2)\).
Ktorý z vektorov \(\vec{u} = \left (- \frac{2}
{\sqrt{2}};2\sqrt{2}\right )\),
\(\vec{v} = (-5;10)\),
\(\vec{w} = (2{,}5;-5)\),
\(\vec{r} = (-3{,}5;6)\) nie je rovnobežný
s vektorom \(\vec{a}\)?}
{\(\vec{r}\)}{\(\vec{w}\)}{\(\vec{v}\)}{\(\vec{u}\)}{}{}}
\LANGes{
\Question{
Dado el vector \(\vec{a} = (1;-2)\).
Cuál de los vectores \(\vec{u} = \left (- \frac{2}
{\sqrt{2}};2\sqrt{2}\right )\),
\(\vec{v} = (-5;10)\),
\(\vec{w} = (2{,}5;-5)\),
\(\vec{r} = (-3{,}5;6)\) no es paralelo al vector
 \(\vec{a}\)?}
{\(\vec{r}\)}{\(\vec{w}\)}{\(\vec{v}\)}{\(\vec{u}\)}{}{}}
\End

\Data\ID{9000101805}
\Updated{2020-07-15 03:07:37}
\Level{B}
\SubArea{Points and vectors}
\LANGen{
\Question{
Find the vector \(\vec{v}\) such that
the length of this vector is \(5\) 
and \(\vec{v}\) is perpendicular
to the vector \(\vec{u} = (-1;0.75)\).}
{\(\vec{v} = (3;4)\)}{\(\vec{v} = (3;-4)\)}{\(\vec{v} = (4;-3)\)}{\(\vec{v} = (5;0)\)}{}{}}
\LANGpl{
\Question{
Wskaż wektor \(\vec{v}\) tak, aby
jego długość wynosiła \(5\) 
i był prostopadły
do wektora \(\vec{u} = (-1;0.75)\).}
{\(\vec{v} = (3;4)\)}{\(\vec{v} = (3;-4)\)}{\(\vec{v} = (4;-3)\)}{\(\vec{v} = (5;0)\)}{}{}}
\LANGcs{
\Question{
Je dán vektor \(\vec{u} = (-1;0{,}75)\).
Vyberte vektor \(\vec{v}\),
pro který platí \(\vec{v} \perp \vec{ u}\) 
a \(|\vec{v}| = 5\).}
{\(\vec{v} = (3;4)\)}{\(\vec{v} = (3;-4)\)}{\(\vec{v} = (4;-3)\)}{\(\vec{v} = (5;0)\)}{}{}}
\LANGsk{
\Question{
Je daný vektor \(\vec{u} = (-1;0{,}75)\).
Vyberte vektor \(\vec{v}\),
pre ktorý platí \(\vec{v} \perp \vec{ u}\) 
a \(|\vec{v}| = 5\).}
{\(\vec{v} = (3;4)\)}{\(\vec{v} = (3;-4)\)}{\(\vec{v} = (4;-3)\)}{\(\vec{v} = (5;0)\)}{}{}}
\LANGes{
\Question{
Halla el vector \(\vec{v}\) suponiendo que
la longitud del vector es \(5\) 
y \(\vec{v}\) es perpendicular
al vector \(\vec{u} = (-1;0.75)\).}
{\(\vec{v} = (3;4)\)}{\(\vec{v} = (3;-4)\)}{\(\vec{v} = (4;-3)\)}{\(\vec{v} = (5;0)\)}{}{}}
\End

\Data\ID{9000101806}
\Updated{2020-07-15 03:07:37}
\Level{B}
\SubArea{Points and vectors}
\LANGen{
\Question{
Find the value of the parameter \(a\) 
which ensures that the vectors \(\vec{u} = (3;a;-2)\) 
and \(\vec{v} = (-6;4;a - 3)\) 
are perpendicular.}
{\(a = 6\)}{\(a = 12\)}{\(a = -6\)}{\(a = 3\)}{}{}}
\LANGpl{
\Question{
Wektory \(\vec{u} = (3;a;-2)\) 
i \(\vec{v} = (-6;4;a - 3)\) 
są prostopadłe jeśli:}
{\(a = 6\)}{\(a = 12\)}{\(a = -6\)}{\(a = 3\)}{}{}}
\LANGcs{
\Question{
Jsou dány vektory \(\vec{u} = (3;a;-2)\),
\(\vec{v} = (-6;4;a - 3)\). Pro které
\(a\in \mathbb{R}\) jsou
vektory \(\vec{u}\) 
a \(\vec{v}\) 
navzájem kolmé?}
{\(a = 6\)}{\(a = 12\)}{\(a = -6\)}{\(a = 3\)}{}{}}
\LANGsk{
\Question{
Sú dané vektory \(\vec{u} = (3;a;-2)\),
\(\vec{v} = (-6;4;a - 3)\). Pre ktoré
\(a\in \mathbb{R}\) sú
vektory \(\vec{u}\) 
a \(\vec{v}\) 
navzájom kolmé?}
{\(a = 6\)}{\(a = 12\)}{\(a = -6\)}{\(a = 3\)}{}{}}
\LANGes{
\Question{
Halla el valor del parámetro \(a\) 
para que los vectores \(\vec{u} = (3;a;-2)\) 
y \(\vec{v} = (-6;4;a - 3)\) 
sean perpendiculares.}
{\(a = 6\)}{\(a = 12\)}{\(a = -6\)}{\(a = 3\)}{}{}}
\End

\Data\ID{9000101807}
\Updated{2020-07-15 03:07:37}
\Level{B}
\SubArea{Points and vectors}
\LANGen{
\Question{
Given points \(A = [1;1]\),
\(B = [5;2]\) and
\(C = [8;7]\), find the
angle \(\measuredangle ABC\).}
{\(135^{\circ }\)}{\(26.5^{\circ }\)}{\(30^{\circ }\)}{\(60^{\circ }\)}{}{}}
\LANGpl{
\Question{
Dane są punkty \(A = [1;1]\),
\(B = [5;2]\) i
\(C = [8;7]\), kąt
 \(\measuredangle ABC\) jest równy:}
{\(135^{\circ }\)}{\(26.5^{\circ }\)}{\(30^{\circ }\)}{\(60^{\circ }\)}{}{}}
\LANGcs{
\Question{
V rovině jsou dány body \(A = [1;1]\),
\(B = [5;2]\),
\(C = [8;7]\). Velikost
úhlu \(ABC\) 
je rovna:}
{\(135^{\circ }\)}{\(26{,}5^{\circ }\)}{\(30^{\circ }\)}{\(60^{\circ }\)}{}{}}
\LANGsk{
\Question{
V rovine sú dané body \(A = [1;1]\),
\(B = [5;2]\),
\(C = [8;7]\). Veľkosť
uhla \(ABC\) 
je rovná:}
{\(135^{\circ }\)}{\(26{,}5^{\circ }\)}{\(30^{\circ }\)}{\(60^{\circ }\)}{}{}}
\LANGes{
\Question{
Dados los puntos \(A = [1;1]\),
\(B = [5;2]\) y
\(C = [8;7]\), halla el ángulo \(\measuredangle ABC\).}
{\(135^{\circ }\)}{\(26.5^{\circ }\)}{\(30^{\circ }\)}{\(60^{\circ }\)}{}{}}
\End

\Data\ID{9000101808}
\Updated{2020-07-15 03:07:37}
\Level{B}
\SubArea{Points and vectors}
\IMG{001018_08-figure0.pdf}
\LANGen{
\Question{
Consider a parallelogram \(ABCD\) 
with \(A = [1;3]\),
\(B = [2;-1]\) and
\(C = [5;1]\). Let
\(S\) be the center of
the diagonal \(BD\).
Find the vector \(\overrightarrow{AS } \).}
{\(\overrightarrow{AS } = (2;-1)\)}{\(\overrightarrow{AS } = (2;1)\)}{\(\overrightarrow{AS } = (1;3)\)}{\(\overrightarrow{AS } = (-2;1)\)}{}{}}
\LANGpl{
\Question{
Dany są punkty
 \(A = [1;3]\),
\(B = [2;-1]\) i
\(C = [5;1]\). Niech
\(S\) będzie środkiem
przekątnej \(BD\) czworokątu \(ABCD\), który jest  równoległobokiem wtedy i tylko wtedy, gdy wektor \(\overrightarrow{AS } \)  jest równy}
{\(\overrightarrow{AS } = (2;-1)\)}{\(\overrightarrow{AS } = (2;1)\)}{\(\overrightarrow{AS } = (1;3)\)}{\(\overrightarrow{AS } = (-2;1)\)}{}{}}
\LANGcs{
\Question{
Je dán rovnoběžník $ABCD$ s vrcholy \(A = [1;3]\),
\(B = [2;-1]\) a
\(C = [5;1]\). Najděte vektor $\overrightarrow{AS}$, kde \(S\) značí střed
úsečky \(BD\).}
{\(\overrightarrow{AS } = (2;-1)\)}{\(\overrightarrow{AS } = (2;1)\)}{\(\overrightarrow{AS } = (1;3)\)}{\(\overrightarrow{AS } = (-2;1)\)}{}{}}
\LANGsk{
\Question{
Je daný rovnobežník $ ABCD $ s vrcholmi \(A = [1; 3] \),
\(B = [2; -1] \) a
\(C = [5; 1] \). Nájdite vektor $ \overrightarrow{AS} $, kde \(S \) značí stred
úsečky \( BD \).}
{\(\overrightarrow{AS } = (2;-1)\)}{\(\overrightarrow{AS } = (2;1)\)}{\(\overrightarrow{AS } = (1;3)\)}{\(\overrightarrow{AS } = (-2;1)\)}{}{}}
\LANGes{
\Question{
Considera un paralelogramo \(ABCD\) 
con \(A = [1;3]\),
\(B = [2;-1]\) y
\(C = [5;1]\). Sea
\(S\) el centro de la
diagonal \(BD\).
Halla el vector \(\overrightarrow{AS } \).}
{\(\overrightarrow{AS } = (2;-1)\)}{\(\overrightarrow{AS } = (2;1)\)}{\(\overrightarrow{AS } = (1;3)\)}{\(\overrightarrow{AS } = (-2;1)\)}{}{}}
\End

\Data\ID{9000108701}
\Updated{2020-07-15 03:07:37}
\Level{B}
\SubArea{Points and vectors}
\LANGen{
\Question{
Find all vectors which are perpendicular to the vector
\(\vec{u} = (3;4)\) and have the
length equal to \(1\).}
{\(\left (\frac{4}
{5};-\frac{3}
{5}\right )\),
\(\left (-\frac{4}
{5}; \frac{3} 
{5}\right )\)}{\(\left (\frac{4}
{7};-\frac{3}
{7}\right )\),
\(\left (-\frac{4}
{7}; \frac{3} 
{7}\right )\)}{\(\left ( \frac{1}
{\sqrt{10}};- \frac{3}
{\sqrt{10}}\right )\),
\(\left (- \frac{1}
{\sqrt{10}}; \frac{3} 
{\sqrt{10}}\right )\)}{\(\left (\frac{4}
{5}; \frac{3} 
{5}\right )\),
\(\left (-\frac{4}
{5};-\frac{3}
{5}\right )\)}{}{}}
\LANGpl{
\Question{
Wskaż wszystkie wektory, które są prostopadłe do wektora
\(\vec{u} = (3;4)\), a ich długości wynosi 
 \(1\).}
{\(\left (\frac{4}
{5};-\frac{3}
{5}\right )\),
\(\left (-\frac{4}
{5}; \frac{3} 
{5}\right )\)}{\(\left (\frac{4}
{7};-\frac{3}
{7}\right )\),
\(\left (-\frac{4}
{7}; \frac{3} 
{7}\right )\)}{\(\left ( \frac{1}
{\sqrt{10}};- \frac{3}
{\sqrt{10}}\right )\),
\(\left (- \frac{1}
{\sqrt{10}}; \frac{3} 
{\sqrt{10}}\right )\)}{\(\left (\frac{4}
{5}; \frac{3} 
{5}\right )\),
\(\left (-\frac{4}
{5};-\frac{3}
{5}\right )\)}{}{}}
\LANGcs{
\Question{
Najděte všechny vektory, které mají velikost
\(1\) a jsou kolmé
k vektoru \(\vec{u} = (3;4)\).}
{\(\left (\frac{4}
{5};-\frac{3}
{5}\right )\),
\(\left (-\frac{4}
{5}; \frac{3} 
{5}\right )\)}{\(\left (\frac{4}
{7};-\frac{3}
{7}\right )\),
\(\left (-\frac{4}
{7}; \frac{3} 
{7}\right )\)}{\(\left ( \frac{1}
{\sqrt{10}};- \frac{3}
{\sqrt{10}}\right )\),
\(\left (- \frac{1}
{\sqrt{10}}; \frac{3} 
{\sqrt{10}}\right )\)}{\(\left (\frac{4}
{5}; \frac{3} 
{5}\right )\),
\(\left (-\frac{4}
{5};-\frac{3}
{5}\right )\)}{}{}}
\LANGsk{
\Question{
Nájdite všetky vektory, ktoré majú veľkosť
\(1\) a sú kolmé
k vektoru \(\vec{u} = (3;4)\).}
{\(\left (\frac{4}
{5};-\frac{3}
{5}\right )\),
\(\left (-\frac{4}
{5}; \frac{3} 
{5}\right )\)}{\(\left (\frac{4}
{7};-\frac{3}
{7}\right )\),
\(\left (-\frac{4}
{7}; \frac{3} 
{7}\right )\)}{\(\left ( \frac{1}
{\sqrt{10}};- \frac{3}
{\sqrt{10}}\right )\),
\(\left (- \frac{1}
{\sqrt{10}}; \frac{3} 
{\sqrt{10}}\right )\)}{\(\left (\frac{4}
{5}; \frac{3} 
{5}\right )\),
\(\left (-\frac{4}
{5};-\frac{3}
{5}\right )\)}{}{}}
\LANGes{
\Question{
Halla todos los vectores que son perpendiculares al vector
\(\vec{u} = (3;4)\) y tienen la longitud que equivale a \(1\).}
{\(\left (\frac{4}
{5};-\frac{3}
{5}\right )\),
\(\left (-\frac{4}
{5}; \frac{3} 
{5}\right )\)}{\(\left (\frac{4}
{7};-\frac{3}
{7}\right )\),
\(\left (-\frac{4}
{7}; \frac{3} 
{7}\right )\)}{\(\left ( \frac{1}
{\sqrt{10}};- \frac{3}
{\sqrt{10}}\right )\),
\(\left (- \frac{1}
{\sqrt{10}}; \frac{3} 
{\sqrt{10}}\right )\)}{\(\left (\frac{4}
{5}; \frac{3} 
{5}\right )\),
\(\left (-\frac{4}
{5};-\frac{3}
{5}\right )\)}{}{}}
\End

\Data\ID{9000108702}
\Updated{2020-07-15 03:07:37}
\Level{B}
\SubArea{Points and vectors}
\LANGen{
\Question{
A square has its center in \([1;4]\) 
and a vertex in \([-1;2]\).
Find the remaining vertices of the square.}
{\([3;6]\),
\([-1;6]\),
\([3;2]\)}{\([3;6]\),
\([-1;5]\),
\([3;1]\)}{\([3;6]\),
\([-2;6]\),
\([4;2]\)}{\([3;6]\),
\([-1;5]\),
\([3;2]\)}{}{}}
\LANGpl{
\Question{
Określ współrzędne pozostałych wierzchołków kwadratu, o wierzchołku   \([-1;2]\),
i  środku przekątnych   \([1;4]\).}
{\([3;6]\),
\([-1;6]\),
\([3;2]\)}{\([3;6]\),
\([-1;5]\),
\([3;1]\)}{\([3;6]\),
\([-2;6]\),
\([4;2]\)}{\([3;6]\),
\([-1;5]\),
\([3;2]\)}{}{}}
\LANGcs{
\Question{
Čtverec má jeden vrchol \([-1;2]\) a průsečík uhlopříček \([1;4]\). Určete souřadnice zbývajících vrcholů.}
{\([3;6]\),
\([-1;6]\),
\([3;2]\)}{\([3;6]\),
\([-1;5]\),
\([3;1]\)}{\([3;6]\),
\([-2;6]\),
\([4;2]\)}{\([3;6]\),
\([-1;5]\),
\([3;2]\)}{}{}}
\LANGsk{
\Question{
Štvorec má jeden vrchol \([- 1; 2] \) a priesečník uhlopriečok \( [1; 4] \). Určte súradnice zvyšných vrcholov.}
{\([3;6]\),
\([-1;6]\),
\([3;2]\)}{\([3;6]\),
\([-1;5]\),
\([3;1]\)}{\([3;6]\),
\([-2;6]\),
\([4;2]\)}{\([3;6]\),
\([-1;5]\),
\([3;2]\)}{}{}}
\LANGes{
\Question{
Un cuadrado tiene su centro en \([1;4]\) 
y un vértice en \([-1;2]\).
Halla otros vértices del cuadrado.}
{\([3;6]\),
\([-1;6]\),
\([3;2]\)}{\([3;6]\),
\([-1;5]\),
\([3;1]\)}{\([3;6]\),
\([-2;6]\),
\([4;2]\)}{\([3;6]\),
\([-1;5]\),
\([3;2]\)}{}{}}
\End

\Data\ID{9000108703}
\Updated{2020-07-15 03:07:37}
\Level{B}
\SubArea{Points and vectors}
\LANGen{
\Question{
Given points \(A = [1;3]\),
\(C = [4;3]\),
\(B = [x;2]\), find the value of the
parameter \(x\) which ensures
that the vector \(AB\) is
perpendicular to the vector \(AC\).}
{\(1\)}{\(2\)}{\(- 1\)}{\(0\)}{}{}}
\LANGpl{
\Question{
Dane są punkty \(A = [1;3]\),
\(C = [4;3]\),
\(B = [x;2]\), określ współrzędną 
 \(x\) tak, aby wektor
 \(AB\) był
prostopadły do wektora \(AC\).}
{\(1\)}{\(2\)}{\(- 1\)}{\(0\)}{}{}}
\LANGcs{
\Question{
Jsou dány body \(A = [1;3]\),
\(C = [4;3]\),
\(B = [x;2]\). Určete
souřadnici \(x\) tak,
aby byl vektor \(AB\) 
kolmý k vektoru \(AC\).}
{\(1\)}{\(2\)}{\(- 1\)}{\(0\)}{}{}}
\LANGsk{
\Question{
Sú dané body \(A = [1;3]\),
\(C = [4;3]\),
\(B = [x;2]\). Určte
súradnicu \(x\) tak,
aby bol vektor \(AB\) 
kolmý k vektoru \(AC\).}
{\(1\)}{\(2\)}{\(- 1\)}{\(0\)}{}{}}
\LANGes{
\Question{
Dados los puntos \(A = [1;3]\),
\(C = [4;3]\),
\(B = [x;2]\), halla el valor del parámetro \(x\) para que el vector \(AB\) sea
perpendicular al vector \(AC\).}
{\(1\)}{\(2\)}{\(- 1\)}{\(0\)}{}{}}
\End

\Data\ID{9000108704}
\Updated{2020-07-15 03:07:37}
\Level{B}
\SubArea{Points and vectors}
\LANGen{
\Question{
Consider a pair of vectors \(\vec{u} = (1;0;-1)\) 
and \(\vec{v} = (2;-1;1)\). Find all the
vectors \(\vec{w}\) which are
perpendicular to both \(\vec{u}\) 
and \(\vec{v}\) and
satisfy \(\left |\vec{w}\right | = 2\).}
{\(\vec{w} = \left (\frac{2\sqrt{11}}
 {11} ; \frac{6\sqrt{11}} 
 {11} ; \frac{2\sqrt{11}} 
 {11} \right )\),
\(\vec{w} = \left (-\frac{2\sqrt{11}}
 {11} ;-\frac{6\sqrt{11}}
 {11} ;-\frac{2\sqrt{11}}
 {11} \right )\)}{\(\vec{w} = (-1;-3;-1)\),
\(\vec{w} = (1;3;1)\)}{\(\vec{w} = \left (-\frac{1}
{2};-\frac{3}
{2};-\frac{1}
{2}\right )\),
\(\vec{w} = \left (\frac{1}
{2}; \frac{3} 
{2}; \frac{1} 
{2}\right )\)}{\(\vec{w} = \left (\frac{2\sqrt{2}}
 {3} ; \frac{3\sqrt{2}} 
 {2} ; \frac{2\sqrt{2}} 
 {3} \right )\),
\(\vec{w} = \left (-\frac{2\sqrt{2}}
 {3} ;-\frac{3\sqrt{2}}
 {2} ;-\frac{2\sqrt{2}}
 {3} \right )\)}{}{}}
\LANGpl{
\Question{
Dane są wektory \(\vec{u} = (1;0;-1)\) 
i \(\vec{v} = (2;-1;1)\). Znajdź wszystkie wektory
 \(\vec{w}\), które są prostopadłe
jednocześnie do wektora \(\vec{u}\) 
i \(\vec{v}\) oraz
 \(\left |\vec{w}\right | = 2\).}
{\(\vec{w} = \left (\frac{2\sqrt{11}}
 {11} ; \frac{6\sqrt{11}} 
 {11} ; \frac{2\sqrt{11}} 
 {11} \right )\),
\(\vec{w} = \left (-\frac{2\sqrt{11}}
 {11} ;-\frac{6\sqrt{11}}
 {11} ;-\frac{2\sqrt{11}}
 {11} \right )\)}{\(\vec{w} = (-1;-3;-1)\),
\(\vec{w} = (1;3;1)\)}{\(\vec{w} = \left (-\frac{1}
{2};-\frac{3}
{2};-\frac{1}
{2}\right )\),
\(\vec{w} = \left (\frac{1}
{2}; \frac{3} 
{2}; \frac{1} 
{2}\right )\)}{\(\vec{w} = \left (\frac{2\sqrt{2}}
 {3} ; \frac{3\sqrt{2}} 
 {2} ; \frac{2\sqrt{2}} 
 {3} \right )\),
\(\vec{w} = \left (-\frac{2\sqrt{2}}
 {3} ;-\frac{3\sqrt{2}}
 {2} ;-\frac{2\sqrt{2}}
 {3} \right )\)}{}{}}
\LANGcs{
\Question{
Jsou dány vektory \(\vec{u} = (1;0;-1)\) 
a \(\vec{v} = (2;-1;1)\). Najděte
všechny vektory \(\vec{w}\),
pro které platí \(\vec{w} \perp \vec{ u}\),
\(\vec{w} \perp \vec{ v}\) a
\(\left |\vec{w}\right | = 2\).}
{\(\vec{w} = \left (\frac{2\sqrt{11}}
 {11} ; \frac{6\sqrt{11}} 
 {11} ; \frac{2\sqrt{11}} 
 {11} \right )\),
\(\vec{w} = \left (-\frac{2\sqrt{11}}
 {11} ;-\frac{6\sqrt{11}}
 {11} ;-\frac{2\sqrt{11}}
 {11} \right )\)}{\(\vec{w} = (-1;-3;-1)\),
\(\vec{w} = (1;3;1)\)}{\(\vec{w} = \left (-\frac{1}
{2};-\frac{3}
{2};-\frac{1}
{2}\right )\),
\(\vec{w} = \left (\frac{1}
{2}; \frac{3} 
{2}; \frac{1} 
{2}\right )\)}{\(\vec{w} = \left (\frac{2\sqrt{2}}
 {3} ; \frac{3\sqrt{2}} 
 {2} ; \frac{2\sqrt{2}} 
 {3} \right )\),
\(\vec{w} = \left (-\frac{2\sqrt{2}}
 {3} ;-\frac{3\sqrt{2}}
 {2} ;-\frac{2\sqrt{2}}
 {3} \right )\)}{}{}}
\LANGsk{
\Question{
Sú dané vektory \(\vec{u} = (1;0;-1)\) 
a \(\vec{v} = (2;-1;1)\). Nájdite
všetky vektory \(\vec{w}\),
pre ktoré platí \(\vec{w} \perp \vec{ u}\),
\(\vec{w} \perp \vec{ v}\) a
\(\left |\vec{w}\right | = 2\).}
{\(\vec{w} = \left (\frac{2\sqrt{11}}
 {11} ; \frac{6\sqrt{11}} 
 {11} ; \frac{2\sqrt{11}} 
 {11} \right )\),
\(\vec{w} = \left (-\frac{2\sqrt{11}}
 {11} ;-\frac{6\sqrt{11}}
 {11} ;-\frac{2\sqrt{11}}
 {11} \right )\)}{\(\vec{w} = (-1;-3;-1)\),
\(\vec{w} = (1;3;1)\)}{\(\vec{w} = \left (-\frac{1}
{2};-\frac{3}
{2};-\frac{1}
{2}\right )\),
\(\vec{w} = \left (\frac{1}
{2}; \frac{3} 
{2}; \frac{1} 
{2}\right )\)}{\(\vec{w} = \left (\frac{2\sqrt{2}}
 {3} ; \frac{3\sqrt{2}} 
 {2} ; \frac{2\sqrt{2}} 
 {3} \right )\),
\(\vec{w} = \left (-\frac{2\sqrt{2}}
 {3} ;-\frac{3\sqrt{2}}
 {2} ;-\frac{2\sqrt{2}}
 {3} \right )\)}{}{}}
\LANGes{
\Question{
Considera un par de los vectores \(\vec{u} = (1;0;-1)\) 
y \(\vec{v} = (2;-1;1)\). Halla todos los vectores \(\vec{w}\) que son
perpendiculares  a los vectores \(\vec{u}\) 
y \(\vec{v}\) suponiendo que
\(\left |\vec{w}\right | = 2\).}
{\(\vec{w} = \left (\frac{2\sqrt{11}}
 {11} ; \frac{6\sqrt{11}} 
 {11} ; \frac{2\sqrt{11}} 
 {11} \right )\),
\(\vec{w} = \left (-\frac{2\sqrt{11}}
 {11} ;-\frac{6\sqrt{11}}
 {11} ;-\frac{2\sqrt{11}}
 {11} \right )\)}{\(\vec{w} = (-1;-3;-1)\),
\(\vec{w} = (1;3;1)\)}{\(\vec{w} = \left (-\frac{1}
{2};-\frac{3}
{2};-\frac{1}
{2}\right )\),
\(\vec{w} = \left (\frac{1}
{2}; \frac{3} 
{2}; \frac{1} 
{2}\right )\)}{\(\vec{w} = \left (\frac{2\sqrt{2}}
 {3} ; \frac{3\sqrt{2}} 
 {2} ; \frac{2\sqrt{2}} 
 {3} \right )\),
\(\vec{w} = \left (-\frac{2\sqrt{2}}
 {3} ;-\frac{3\sqrt{2}}
 {2} ;-\frac{2\sqrt{2}}
 {3} \right )\)}{}{}}
\End

\Data\ID{9000108705}
\Updated{2020-07-15 03:07:37}
\Level{B}
\SubArea{Points and vectors}
\LANGen{
\Question{
Given vectors \(\vec{u} = (1;y;3)\) and
\(\vec{v} = (1;2;1)\), find the value of
the parameter \(y\) 
which ensures that the vectors are perpendicular.}
{\(- 2\)}{\(1\)}{\(2\)}{\(0\)}{}{}}
\LANGpl{
\Question{
Dane są wektory \(\vec{u} = (1;y;3)\) i
\(\vec{v} = (1;2;1)\), określ współrzędną
 \(y\) 
tak, aby wektory były wzajemnie prostopadłe.}
{\(- 2\)}{\(1\)}{\(2\)}{\(0\)}{}{}}
\LANGcs{
\Question{
Jsou dány vektory \(\vec{u} = (1;y;3)\) 
a \(\vec{v} = (1;2;1)\). Určete
souřadnici \(y\) 
tak, aby byly zadané vektory navzájem kolmé.}
{\(- 2\)}{\(1\)}{\(2\)}{\(0\)}{}{}}
\LANGsk{
\Question{
Sú dané vektory \(\vec{u} = (1;y;3)\) 
a \(\vec{v} = (1;2;1)\). Určte
súradnicu \(y\) 
tak, aby boli zadané vektory navzájom kolmé.}
{\(- 2\)}{\(1\)}{\(2\)}{\(0\)}{}{}}
\LANGes{
\Question{
Dados los vectores \(\vec{u} = (1;y;3)\) y
\(\vec{v} = (1;2;1)\), halla el valor del parámetro \(y\) 
para que los vectores sean perpendiculares.}
{\(- 2\)}{\(1\)}{\(2\)}{\(0\)}{}{}}
\End

\Data\ID{9000108706}
\Updated{2020-07-15 03:07:37}
\Level{B}
\SubArea{Points and vectors}
\LANGen{
\Question{
Find all vectors which are parallel to the vector
\(\vec{u} = (3;-1)\) and have the
length equal to \(1\).}
{\(\left (\frac{3\sqrt{10}}
 {10} ;-\frac{\sqrt{10}}
 {10} \right )\),
\(\left (-\frac{3\sqrt{10}}
 {10} ; \frac{\sqrt{10}} 
 {10} \right )\)}{\((0;-1)\),
\((0;1)\)}{\((-3;1)\),
\((3;-1)\)}{\(\left (\frac{3}
{4};-\frac{1}
{4}\right )\),
\(\left (-\frac{3}
{4}; \frac{1} 
{4}\right )\)}{}{}}
\LANGpl{
\Question{
Znajdź wszystkie wektory równoległe do wektora 
\(\vec{u} = (3;-1)\) o długości równej
 \(1\).}
{\(\left (\frac{3\sqrt{10}}
 {10} ;-\frac{\sqrt{10}}
 {10} \right )\),
\(\left (-\frac{3\sqrt{10}}
 {10} ; \frac{\sqrt{10}} 
 {10} \right )\)}{\((0;-1)\),
\((0;1)\)}{\((-3;1)\),
\((3;-1)\)}{\(\left (\frac{3}
{4};-\frac{1}
{4}\right )\),
\(\left (-\frac{3}
{4}; \frac{1} 
{4}\right )\)}{}{}}
\LANGcs{
\Question{
Najděte všechny vektory rovnoběžné s vektorem
\(\vec{u} = (3;-1)\), které mají
velikost \(1\).}
{\(\left (\frac{3\sqrt{10}}
 {10} ;-\frac{\sqrt{10}}
 {10} \right )\),
\(\left (-\frac{3\sqrt{10}}
 {10} ; \frac{\sqrt{10}} 
 {10} \right )\)}{\((0;-1)\),
\((0;1)\)}{\((-3;1)\),
\((3;-1)\)}{\(\left (\frac{3}
{4};-\frac{1}
{4}\right )\),
\(\left (-\frac{3}
{4}; \frac{1} 
{4}\right )\)}{}{}}
\LANGsk{
\Question{
Nájdite všetky vektory rovnobežné s vektorom
\(\vec{u} = (3;-1)\), ktoré majú
veľkosť \(1\).}
{\(\left (\frac{3\sqrt{10}}
 {10} ;-\frac{\sqrt{10}}
 {10} \right )\),
\(\left (-\frac{3\sqrt{10}}
 {10} ; \frac{\sqrt{10}} 
 {10} \right )\)}{\((0;-1)\),
\((0;1)\)}{\((-3;1)\),
\((3;-1)\)}{\(\left (\frac{3}
{4};-\frac{1}
{4}\right )\),
\(\left (-\frac{3}
{4}; \frac{1} 
{4}\right )\)}{}{}}
\LANGes{
\Question{
Halla todos los vectores que son paralelos al vector
\(\vec{u} = (3;-1)\) y tienen la longitud que equivale a \(1\).}
{\(\left (\frac{3\sqrt{10}}
 {10} ;-\frac{\sqrt{10}}
 {10} \right )\),
\(\left (-\frac{3\sqrt{10}}
 {10} ; \frac{\sqrt{10}} 
 {10} \right )\)}{\((0;-1)\),
\((0;1)\)}{\((-3;1)\),
\((3;-1)\)}{\(\left (\frac{3}
{4};-\frac{1}
{4}\right )\),
\(\left (-\frac{3}
{4}; \frac{1} 
{4}\right )\)}{}{}}
\End

\Data\ID{9000108801}
\Updated{2020-07-15 03:07:37}
\Level{B}
\SubArea{Points and vectors}
\LANGen{
\Question{
Find the angle between the vector \(\vec{u} = (1;\sqrt{2})\) 
and \(\vec{v} = (3;-1)\).
Round to the nearest degree.}
{\(73^{\circ }\)}{\(42^{\circ }\)}{\(57^{\circ }\)}{\(64^{\circ }\)}{}{}}
\LANGpl{
\Question{
Określ kąt pomiędzy wektorem \(\vec{u} = (1;\sqrt{2})\) 
i \(\vec{v} = (3;-1)\).
Wynik zaokrągli do pełnych stopni.}
{\(73^{\circ }\)}{\(42^{\circ }\)}{\(57^{\circ }\)}{\(64^{\circ }\)}{}{}}
\LANGcs{
\Question{
Vypočítejte odchylku vektorů \(\vec{u} = (1;\sqrt{2})\) 
a \(\vec{v} = (3;-1)\).
Zaokrouhlete na celé stupně.}
{\(73^{\circ }\)}{\(42^{\circ }\)}{\(57^{\circ }\)}{\(64^{\circ }\)}{}{}}
\LANGsk{
\Question{
Vypočítajte odchýlku vektorov \(\vec{u} = (1;\sqrt{2})\) 
a \(\vec{v} = (3;-1)\).
Zaokrúhlite na celé stupne.}
{\(73^{\circ }\)}{\(42^{\circ }\)}{\(57^{\circ }\)}{\(64^{\circ }\)}{}{}}
\LANGes{
\Question{
Halla el ángulo entre los vectores \(\vec{u} = (1;\sqrt{2})\) 
y \(\vec{v} = (3;-1)\).
Redondea al grado más cercano.}
{\(73^{\circ }\)}{\(42^{\circ }\)}{\(57^{\circ }\)}{\(64^{\circ }\)}{}{}}
\End

\Data\ID{9000108802}
\Updated{2020-07-15 03:07:37}
\Level{B}
\SubArea{Points and vectors}
\LANGen{
\Question{
Given the points \(A = [1;2]\),
\(B = [2;6]\) and
\(C = [3;-1]\), find the interior
angles of the triangle \(ABC\).
Round to the nearest degree.}
{\(22^{\circ }\),
\(26^{\circ }\),
\(132^{\circ }\)}{\(26^{\circ }\),
\(45^{\circ }\),
\(109^{\circ }\)}{\(22^{\circ }\),
\(48^{\circ }\),
\(110^{\circ }\)}{\(17^{\circ }\),
\(31^{\circ }\),
\(132^{\circ }\)}{}{}}
\LANGpl{
\Question{
Dane są punkty \(A = [1;2]\),
\(B = [2;6]\) i
\(C = [3;-1]\), określ kąty wewnętrzne trójkąta  \(ABC\).
Wynik zaokrągli do pełnych stopni.}
{\(22^{\circ }\),
\(26^{\circ }\),
\(132^{\circ }\)}{\(26^{\circ }\),
\(45^{\circ }\),
\(109^{\circ }\)}{\(22^{\circ }\),
\(48^{\circ }\),
\(110^{\circ }\)}{\(17^{\circ }\),
\(31^{\circ }\),
\(132^{\circ }\)}{}{}}
\LANGcs{
\Question{
Určete velikost vnitřních úhlů trojúhelníku
\(ABC\), je-li
\(A = [1;2]\),
\(B = [2;6]\),
\(C = [3;-1]\).
Zaokrouhlete na celé stupně.}
{\(22^{\circ }\),
\(26^{\circ }\),
\(132^{\circ }\)}{\(26^{\circ }\),
\(45^{\circ }\),
\(109^{\circ }\)}{\(22^{\circ }\),
\(48^{\circ }\),
\(110^{\circ }\)}{\(17^{\circ }\),
\(31^{\circ }\),
\(132^{\circ }\)}{}{}}
\LANGsk{
\Question{
Určte veľkosť vnútorných uhlov trojuholníka
\(ABC\), ak
\(A = [1;2]\),
\(B = [2;6]\),
\(C = [3;-1]\).
Zaokrúhlite na celé stupne.}
{\(22^{\circ }\),
\(26^{\circ }\),
\(132^{\circ }\)}{\(26^{\circ }\),
\(45^{\circ }\),
\(109^{\circ }\)}{\(22^{\circ }\),
\(48^{\circ }\),
\(110^{\circ }\)}{\(17^{\circ }\),
\(31^{\circ }\),
\(132^{\circ }\)}{}{}}
\LANGes{
\Question{
Dados los puntos \(A = [1;2]\),
\(B = [2;6]\) y
\(C = [3;-1]\), halla los ángulos interiores del triángulo \(ABC\).
Redondea al grado más cercano.}
{\(22^{\circ }\),
\(26^{\circ }\),
\(132^{\circ }\)}{\(26^{\circ }\),
\(45^{\circ }\),
\(109^{\circ }\)}{\(22^{\circ }\),
\(48^{\circ }\),
\(110^{\circ }\)}{\(17^{\circ }\),
\(31^{\circ }\),
\(132^{\circ }\)}{}{}}
\End

\Data\ID{9000108803}
\Updated{2020-07-15 03:07:37}
\Level{B}
\SubArea{Points and vectors}
\LANGen{
\Question{
Consider the vector \(\vec{u} = (\sqrt{3};1)\).
Find the vector \(\vec{w}\) 
such that \(\left |\vec{w}\right | = 4\) and the
angle between \(\vec{u}\) 
and \(\vec{w}\) is
\(60^{\circ }\). Find
all solutions.}
{\(\vec{w} = (0;4)\),
\(\vec{w} = (2\sqrt{3};-2)\)}{\(\vec{w} = (0;-4)\),
\(\vec{w} = (\sqrt{7};-3)\)}{\(\vec{w} = (0;4)\),
\(\vec{w} = (\sqrt{7};3)\)}{\(\vec{w} = (\sqrt{5};\sqrt{11})\),
\(\vec{w} = (2\sqrt{3};-2)\)}{}{}}
\LANGpl{
\Question{
Dany jest wektor \(\vec{u} = (\sqrt{3};1)\).
Wyznacz wszystkie wektory \(\vec{w}\) 
zakładając, że  \(\left |\vec{w}\right | = 4\), a miara kąta pomiędzy
 \(\vec{u}\) 
i \(\vec{w}\) wynosi
\(60^{\circ }\).}
{\(\vec{w} = (0;4)\),
\(\vec{w} = (2\sqrt{3};-2)\)}{\(\vec{w} = (0;-4)\),
\(\vec{w} = (\sqrt{7};-3)\)}{\(\vec{w} = (0;4)\),
\(\vec{w} = (\sqrt{7};3)\)}{\(\vec{w} = (\sqrt{5};\sqrt{11})\),
\(\vec{w} = (2\sqrt{3};-2)\)}{}{}}
\LANGcs{
\Question{
Je dán vektor \(\vec{u} = (\sqrt{3};1)\). Najděte
všechny vektory \(\vec{w}\) 
takové, že \(\left |\vec{w}\right | = 4\) a odchylka vektorů \(\vec{u}\),
\(\vec{w}\) je
\(60^{\circ }\).}
{\(\vec{w} = (0;4)\),
\(\vec{w} = (2\sqrt{3};-2)\)}{\(\vec{w} = (0;-4)\),
\(\vec{w} = (\sqrt{7};-3)\)}{\(\vec{w} = (0;4)\),
\(\vec{w} = (\sqrt{7};3)\)}{\(\vec{w} = (\sqrt{5};\sqrt{11})\),
\(\vec{w} = (2\sqrt{3};-2)\)}{}{}}
\LANGsk{
\Question{
Je daný vektor \(\vec{u} = (\sqrt{3};1)\). Nájdite
všetky vektory \(\vec{w}\) 
také, že \(\left |\vec{w}\right | = 4\) a
odchýlka vektorov \(\vec{u}\),
\(\vec{w}\) je
\(60^{\circ }\).}
{\(\vec{w} = (0;4)\),
\(\vec{w} = (2\sqrt{3};-2)\)}{\(\vec{w} = (0;-4)\),
\(\vec{w} = (\sqrt{7};-3)\)}{\(\vec{w} = (0;4)\),
\(\vec{w} = (\sqrt{7};3)\)}{\(\vec{w} = (\sqrt{5};\sqrt{11})\),
\(\vec{w} = (2\sqrt{3};-2)\)}{}{}}
\LANGes{
\Question{
Considera el vector \(\vec{u} = (\sqrt{3};1)\).
Halla el vector \(\vec{w}\) 
suponiendo que \(\left |\vec{w}\right | = 4\) y el
ángulo entre \(\vec{u}\) 
y \(\vec{w}\) es
\(60^{\circ }\). Halla todas las soluciones.}
{\(\vec{w} = (0;4)\),
\(\vec{w} = (2\sqrt{3};-2)\)}{\(\vec{w} = (0;-4)\),
\(\vec{w} = (\sqrt{7};-3)\)}{\(\vec{w} = (0;4)\),
\(\vec{w} = (\sqrt{7};3)\)}{\(\vec{w} = (\sqrt{5};\sqrt{11})\),
\(\vec{w} = (2\sqrt{3};-2)\)}{}{}}
\End

\Data\ID{9000108804}
\Updated{2020-07-15 03:07:37}
\Level{B}
\SubArea{Points and vectors}
\LANGen{
\Question{
The point \(A = [3;2]\) is rotated
about the center \(B = [1;1]\) 
by \(60^{\circ }\). Find
the coordinate of its final position. Consider both clockwise and counterclockwise
direction.}
{\(\left [2\pm \frac{\sqrt{3}}
 {2} ; \frac{3} 
{2} \mp \sqrt{3}\right ]\)}{\(\left [1\pm \frac{\sqrt{3}}
 {2} ; \frac{1} 
{2} \mp \sqrt{3}\right ]\)}{\(\left [2\pm \frac{\sqrt{2}}
 {2} ; \frac{3} 
{2} \mp \sqrt{2}\right ]\)}{\(\left [1\pm \frac{\sqrt{2}}
 {2} ; \frac{1} 
{2} \mp \sqrt{2}\right ]\)}{}{}}
\LANGpl{
\Question{
Wyznacz współrzędne punktu \(A^{\prime}\), który 
powstaje przez obrót punktu  \(A = [3;2]\) wokół \(B = [1;1]\) 
o kąt  \(60^{\circ }\). Weź pod uwagę zarówno dodatni, jak i ujemny kąt.}
{\(\left [2\pm \frac{\sqrt{3}}
 {2} ; \frac{3} 
{2} \mp \sqrt{3}\right ]\)}{\(\left [1\pm \frac{\sqrt{3}}
 {2} ; \frac{1} 
{2} \mp \sqrt{3}\right ]\)}{\(\left [2\pm \frac{\sqrt{2}}
 {2} ; \frac{3} 
{2} \mp \sqrt{2}\right ]\)}{\(\left [1\pm \frac{\sqrt{2}}
 {2} ; \frac{1} 
{2} \mp \sqrt{2}\right ]\)}{}{}}
\LANGcs{
\Question{
Určete body, které vzniknou rotací bodu $A=[3;2]$ okolo bodu $B=[1;1]$ o $60^{\circ}$. Uvažujte rotaci v kladném i záporném smyslu.}
{\(\left [2\pm \frac{\sqrt{3}}
 {2} ; \frac{3} 
{2} \mp \sqrt{3}\right ]\)}{\(\left [1\pm \frac{\sqrt{3}}
 {2} ; \frac{1} 
{2} \mp \sqrt{3}\right ]\)}{\(\left [2\pm \frac{\sqrt{2}}
 {2} ; \frac{3} 
{2} \mp \sqrt{2}\right ]\)}{\( \left [1\pm \frac{\sqrt{2}}
 {2} ; \frac{1} 
{2} \mp \sqrt{2}\right ]\)}{}{}}
\LANGsk{
\Question{
Určite body, ktoré vzniknú rotáciou bodu $ A = [3; 2] $ okolo bodu $ B = [1; 1] $ o $ 60^{\circ} $. Uvažujte rotáciu v kladnom i zápornom zmysle.}
{\(\left [2\pm \frac{\sqrt{3}}
 {2} ; \frac{3} 
{2} \mp \sqrt{3}\right ]\)}{\(\left [1\pm \frac{\sqrt{3}}
 {2} ; \frac{1} 
{2} \mp \sqrt{3}\right ]\)}{\(\left [2\pm \frac{\sqrt{2}}
 {2} ; \frac{3} 
{2} \mp \sqrt{2}\right ]\)}{\(\left [1\pm \frac{\sqrt{2}}
 {2} ; \frac{1} 
{2} \mp \sqrt{2}\right ]\)}{}{}}
\LANGes{
\Question{
Determina los puntos que surgen gracias a la rotación del punto $A=[3;2]$ alrededor del punto $B=[1;1]$ en $60^{\circ}$. Considera la rotación positiva y negativa.}
{\(\left [2\pm \frac{\sqrt{3}}
 {2} ; \frac{3} 
{2} \mp \sqrt{3}\right ]\)}{\(\left [1\pm \frac{\sqrt{3}}
 {2} ; \frac{1} 
{2} \mp \sqrt{3}\right ]\)}{\(\left [2\pm \frac{\sqrt{2}}
 {2} ; \frac{3} 
{2} \mp \sqrt{2}\right ]\)}{\(\left [1\pm \frac{\sqrt{2}}
 {2} ; \frac{1} 
{2} \mp \sqrt{2}\right ]\)}{}{}}
\End

\Data\ID{9000108805}
\Updated{2020-07-15 03:07:37}
\Level{B}
\SubArea{Points and vectors}
\LANGen{
\Question{
Find the angle between \(\vec{u} = (1;-2;3)\) 
and \(\vec{v} = (-1;0;2)\).
Round to the nearest degree.}
{\(53^{\circ }\)}{\(27^{\circ }\)}{\(60^{\circ }\)}{\(46^{\circ }\)}{}{}}
\LANGpl{
\Question{
Oblicz kąt pomiędzy \(\vec{u} = (1;-2;3)\) 
i \(\vec{v} = (-1;0;2)\).
Wynik zaokrągli do pełnych stopni.}
{\(53^{\circ }\)}{\(27^{\circ }\)}{\(60^{\circ }\)}{\(46^{\circ }\)}{}{}}
\LANGcs{
\Question{
Vypočítejte odchylku vektorů \(\vec{u} = (1;-2;3)\) 
a \(\vec{v} = (-1;0;2)\).
Zaokrouhlete na celé stupně.}
{\(53^{\circ }\)}{\(27^{\circ }\)}{\(60^{\circ }\)}{\(46^{\circ }\)}{}{}}
\LANGsk{
\Question{
Vypočítajte odchýlku vektorov \(\vec{u} = (1;-2;3)\) 
a \(\vec{v} = (-1;0;2)\).
Zaokrúhlite na celé stupne.}
{\(53^{\circ }\)}{\(27^{\circ }\)}{\(60^{\circ }\)}{\(46^{\circ }\)}{}{}}
\LANGes{
\Question{
Halla el ángulo entre \(\vec{u} = (1;-2;3)\) 
y \(\vec{v} = (-1;0;2)\).
Redondea al grado más cercano.}
{\(53^{\circ }\)}{\(27^{\circ }\)}{\(60^{\circ }\)}{\(46^{\circ }\)}{}{}}
\End

\Data\ID{9000108806}
\Updated{2020-07-15 03:07:37}
\Level{B}
\SubArea{Points and vectors}
\LANGen{
\Question{
Find the coordinate $y$ which ensures that the vectors \(\vec{u} = (-6;y;3)\) and
\(\vec{v} = (12;4;4)\) are perpendicular.}
{\(15\)}{\(12\)}{\(\sqrt{5}\)}{\(\frac{5}
{3}\)}{}{}}
\LANGpl{
\Question{
Wyznacz  \(y\) tak, aby 
wektory \(\vec{u} = (-6;y;3)\) i
\(\vec{v} = (12;4;4)\) były wzajemnie prostopadłe.}
{\(15\)}{\(12\)}{\(\sqrt{5}\)}{\(\frac{5}
{3}\)}{}{}}
\LANGcs{
\Question{
Doplňte souřadnici \(y\) 
tak, aby byly vektory \(\vec{u} = (-6;y;3)\) 
a \(\vec{v} = (12;4;4)\) 
navzájem kolmé.}
{\(15\)}{\(12\)}{\(\sqrt{5}\)}{\(\frac{5}
{3}\)}{}{}}
\LANGsk{
\Question{
Doplňte súradnicu \(y\) 
tak, aby boli vektory \(\vec{u} = (-6;y;3)\) 
a \(\vec{v} = (12;4;4)\) 
navzájom kolmé.}
{\(15\)}{\(12\)}{\(\sqrt{5}\)}{\(\frac{5}
{3}\)}{}{}}
\LANGes{
\Question{
Halla la coordenada $y$ para que los vectores \(\vec{u} = (-6;y;3)\) y \(\vec{v} = (12;4;4)\) sean perpendiculares.}
{\(15\)}{\(12\)}{\(\sqrt{5}\)}{\(\frac{5}
{3}\)}{}{}}
\End

\Data\ID{9000108807}
\Updated{2020-07-15 03:07:37}
\Level{B}
\SubArea{Points and vectors}
\LANGen{
\Question{
Find the angle between the median \(t_{c}\) 
and side \(c\) in
the triangle \(ABC\) 
for \(A = [1;2]\),
\(B = [7;-2]\) and
\(C = [6;1]\).
Round to the nearest degree. Hint: In geometry, the median \(t_{c}\) 
of the triangle \(ABC\) is the line
segment joining the vertex \(C\) 
to the midpoint of the opposing side.}
{\(60^{\circ }\)}{\(50^{\circ }\)}{\(43^{\circ }\)}{\(71^{\circ }\)}{}{}}
\LANGpl{
\Question{
Wyznacz kąt pomiędzy punktem przecięcia środkowych \(t_{c}\), 
a bokiem  \(c\), w
trójkącie\(ABC\), 
jeśli \(A = [1;2]\),
\(B = [7;-2]\) i
\(C = [6;1]\). Wynik zaokrągli do pełnych stopni.}
{\(60^{\circ }\)}{\(50^{\circ }\)}{\(43^{\circ }\)}{\(71^{\circ }\)}{}{}}
\LANGcs{
\Question{
Zjistěte odchylku těžnice \(t_{c}\) 
a strany \(c\) 
trojúhelníku \(ABC\),
je-li \(A = [1;2]\),
\(B = [7;-2]\),
\(C = [6;1]\).
Zaokrouhlete na celé stupně.}
{\(60^{\circ }\)}{\(50^{\circ }\)}{\(43^{\circ }\)}{\(71^{\circ }\)}{}{}}
\LANGsk{
\Question{
Zistite odchýlku ťažnice \(t_{c}\) 
a strany \(c\) 
trojuholníka \(ABC\),
ak \(A = [1;2]\),
\(B = [7;-2]\),
\(C = [6;1]\).
Zaokrúhlite na celé stupne.}
{\(60^{\circ }\)}{\(50^{\circ }\)}{\(43^{\circ }\)}{\(71^{\circ }\)}{}{}}
\LANGes{
\Question{
Halla el ángulo entre la mediana \(t_{c}\) 
y el lado \(c\) en el triángulo \(ABC\) 
para  \(A = [1;2]\),
\(B = [7;-2]\) y
\(C = [6;1]\).
Redondea al grado más cercano. Pista: En la geometría, la mediana \(t_{c}\) 
del triángulo \(ABC\) es el segmento que une el vértice \(C\)  con el punto medio de su lado opuesto.}
{\(60^{\circ }\)}{\(50^{\circ }\)}{\(43^{\circ }\)}{\(71^{\circ }\)}{}{}}
\End

\Data\ID{9000108808}
\Updated{2020-07-15 03:07:37}
\Level{B}
\SubArea{Points and vectors}
\LANGen{
\Question{
Find the angle between the altitude \(v_{c}\) 
and side \(b\) in
the triangle \(ABC\) 
for \(A = [1;2]\),
\(B = [7;-2]\) and
\(C = [6;1]\).
Round to the nearest degree. Hint: In geometry, the altitude \(v_{c}\) 
of the triangle \(ABC\) is the line
segment through the vertex \(C\) 
and perpendicular to the line containing the opposite side of the triangle.}
{\(68^{\circ }\)}{\(75^{\circ }\)}{\(44^{\circ }\)}{\(61^{\circ }\)}{}{}}
\LANGpl{
\Question{
Oblicz kąt pomiędzy wysokością  \(v_{c}\), 
a bokiem \(b\) w
trójkącie \(ABC\) 
dla \(A = [1;2]\),
\(B = [7;-2]\) i
\(C = [6;1]\).
Wynik zaokrągli do pełnych stopni.}
{\(68^{\circ }\)}{\(75^{\circ }\)}{\(44^{\circ }\)}{\(61^{\circ }\)}{}{}}
\LANGcs{
\Question{
Zjistěte odchylku výšky \(v_{c}\) 
a strany \(b\) 
trojúhelníku \(ABC\),
je-li \(A = [1;2]\),
\(B = [7;-2]\),
\(C = [6;1]\).
Zaokrouhlete na celé stupně.}
{\(68^{\circ }\)}{\(75^{\circ }\)}{\(44^{\circ }\)}{\(61^{\circ }\)}{}{}}
\LANGsk{
\Question{
Zistite odchýlku výšky \(v_{c}\) 
a strany \(b\) 
trojuholníka \(ABC\),
ak \(A = [1;2]\),
\(B = [7;-2]\),
\(C = [6;1]\).
Zaokrúhlite na celé stupne.}
{\(68^{\circ }\)}{\(75^{\circ }\)}{\(44^{\circ }\)}{\(61^{\circ }\)}{}{}}
\LANGes{
\Question{
Halla el ángulo entre la altura \(v_{c}\) 
y el lado \(b\) en
el triángulo \(ABC\) 
para \(A = [1;2]\),
\(B = [7;-2]\) y
\(C = [6;1]\).
Redondea al grado más cercano. Pista: En geometría, la altura \(v_{c}\) 
del triángulo \(ABC\) es el segmento perpendicular a un lado que va desde el vértice \(C\)  a este lado.}
{\(68^{\circ }\)}{\(75^{\circ }\)}{\(44^{\circ }\)}{\(61^{\circ }\)}{}{}}
\End

\Data\ID{1003020901}
\Updated{2020-07-15 03:07:12}
\Level{C}
\SubArea{Points and vectors}
\LANGen{
\Question{
Let there be vectors: \(\overrightarrow{a}=(1;3;-1)\), \(\overrightarrow{b}=(0;3;1)\), \(\overrightarrow{c}=(-1;2;2)\). Find \(\overrightarrow{a}\times\overrightarrow{b}\) and \(\left(\overrightarrow{a}\times\overrightarrow{b}\right)\cdot\overrightarrow{c}\).}
{\(\overrightarrow{a}\times\overrightarrow{b}=(6;-1;3); \left(\overrightarrow{a}\times\overrightarrow{b}\right)\cdot\overrightarrow{c}=-2\)}{\(\overrightarrow{a}\times\overrightarrow{b}=8; \left(\overrightarrow{a}\times\overrightarrow{b}\right)\cdot\overrightarrow{c}=(-8,16,16)\)}{\(\overrightarrow{a}\times\overrightarrow{b}=(-6;1;-3); \left(\overrightarrow{a}\times\overrightarrow{b}\right)\cdot\overrightarrow{c}=2\)}{\(\overrightarrow{a}\times\overrightarrow{b}=\sqrt{46}; \left(\overrightarrow{a}\times\overrightarrow{b}\right)\cdot\overrightarrow{c}=2\)}{}{}}
\LANGpl{
\Question{
Dane są wektory: \(\overrightarrow{a}=(1;3;-1)\), \(\overrightarrow{b}=(0;3;1)\), \(\overrightarrow{c}=(-1;2;2)\). Oblicz  \(\overrightarrow{a}\times\overrightarrow{b}\) i \(\left(\overrightarrow{a}\times\overrightarrow{b}\right)\cdot\overrightarrow{c}\).}
{\(\overrightarrow{a}\times\overrightarrow{b}=(6;-1;3); \left(\overrightarrow{a}\times\overrightarrow{b}\right)\cdot\overrightarrow{c}=-2\)}{\(\overrightarrow{a}\times\overrightarrow{b}=8; \left(\overrightarrow{a}\times\overrightarrow{b}\right)\cdot\overrightarrow{c}=(-8,16,16)\)}{\(\overrightarrow{a}\times\overrightarrow{b}=(-6;1;-3); \left(\overrightarrow{a}\times\overrightarrow{b}\right)\cdot\overrightarrow{c}=2\)}{\(\overrightarrow{a}\times\overrightarrow{b}=\sqrt{46}; \left(\overrightarrow{a}\times\overrightarrow{b}\right)\cdot\overrightarrow{c}=2\)}{}{}}
\LANGcs{
\Question{
Jsou dány vektory: \(\overrightarrow{a}=(1;3;-1)\), \(\overrightarrow{b}=(0;3;1)\), \(\overrightarrow{c}=(-1;2;2)\). Vypočtěte \(\overrightarrow{a}\times\overrightarrow{b}\) a \(\left(\overrightarrow{a}\times\overrightarrow{b}\right)\cdot\overrightarrow{c}\).}
{\(\overrightarrow{a}\times\overrightarrow{b}=(6;-1;3); \left(\overrightarrow{a}\times\overrightarrow{b}\right)\cdot\overrightarrow{c}=-2\)}{\(\overrightarrow{a}\times\overrightarrow{b}=8; \left(\overrightarrow{a}\times\overrightarrow{b}\right)\cdot\overrightarrow{c}=(-8,16,16)\)}{\(\overrightarrow{a}\times\overrightarrow{b}=(-6;1;-3); \left(\overrightarrow{a}\times\overrightarrow{b}\right)\cdot\overrightarrow{c}=2\)}{\(\overrightarrow{a}\times\overrightarrow{b}=\sqrt{46}; \left(\overrightarrow{a}\times\overrightarrow{b}\right)\cdot\overrightarrow{c}=2\)}{}{}}
\LANGsk{
\Question{
Sú dané vektory: \(\overrightarrow{a}=(1;3;-1)\), \(\overrightarrow{b}=(0;3;1)\), \(\overrightarrow{c}=(-1;2;2)\). Nájdite \(\overrightarrow{a}\times\overrightarrow{b}\) and \(\left(\overrightarrow{a}\times\overrightarrow{b}\right)\cdot\overrightarrow{c}\).}
{\(\overrightarrow{a}\times\overrightarrow{b}=(6;-1;3); \left(\overrightarrow{a}\times\overrightarrow{b}\right)\cdot\overrightarrow{c}=-2\)}{\(\overrightarrow{a}\times\overrightarrow{b}=8; \left(\overrightarrow{a}\times\overrightarrow{b}\right)\cdot\overrightarrow{c}=(-8,16,16)\)}{\(\overrightarrow{a}\times\overrightarrow{b}=(-6;1;-3); \left(\overrightarrow{a}\times\overrightarrow{b}\right)\cdot\overrightarrow{c}=2\)}{\(\overrightarrow{a}\times\overrightarrow{b}=\sqrt{46}; \left(\overrightarrow{a}\times\overrightarrow{b}\right)\cdot\overrightarrow{c}=2\)}{}{}}
\LANGes{
\Question{
Dados los vectores: \(\overrightarrow{a}=(1;3;-1)\), \(\overrightarrow{b}=(0;3;1)\), \(\overrightarrow{c}=(-1;2;2)\). Determina \(\overrightarrow{a}\times\overrightarrow{b}\) y \(\left(\overrightarrow{a}\times\overrightarrow{b}\right)\cdot\overrightarrow{c}\).}
{\(\overrightarrow{a}\times\overrightarrow{b}=(6;-1;3); \left(\overrightarrow{a}\times\overrightarrow{b}\right)\cdot\overrightarrow{c}=-2\)}{\(\overrightarrow{a}\times\overrightarrow{b}=8; \left(\overrightarrow{a}\times\overrightarrow{b}\right)\cdot\overrightarrow{c}=(-8,16,16)\)}{\(\overrightarrow{a}\times\overrightarrow{b}=(-6;1;-3); \left(\overrightarrow{a}\times\overrightarrow{b}\right)\cdot\overrightarrow{c}=2\)}{\(\overrightarrow{a}\times\overrightarrow{b}=\sqrt{46}; \left(\overrightarrow{a}\times\overrightarrow{b}\right)\cdot\overrightarrow{c}=2\)}{}{}}
\End

\Data\ID{1003040201}
\Updated{2020-07-15 03:07:12}
\Level{C}
\SubArea{Points and vectors}
\LANGen{
\Question{
We are given the vectors $\vec{a}=(-1; 2;3)$, $\vec{b}=(3; 1; -2)$ and $\vec{c}=(1; 2;-1)$. Find the coordinates of a vector $\vec{v}$, such that $\vec{v}$ is perpendicular to both vectors $\vec{a}$ and $\vec{b}$, while $\vec{v}\cdot\vec{c}=12$ holds.}
{$\vec{v}=(-6;6;-6)$}{$\vec{v}=(6;-6;6)$}{$\vec{v}=(-7;7;-7)$}{$\vec{v}=(7;-7;7)$}{}{}}
\LANGpl{
\Question{
Dane są wektory $\vec{a}=(-1; 2;3)$, $\vec{b}=(3; 1; -2)$ i $\vec{c}=(1; 2;-1)$. Wyznacz współrzędne wektora $\vec{v}$ tak, aby $\vec{v}$ był prostopadły do obu wektorów $\vec{a}$ i $\vec{b}$, jeśli $\vec{v}\cdot\vec{c}=12$.}
{$\vec{v}=(-6;6;-6)$}{$\vec{v}=(6;-6;6)$}{$\vec{v}=(-7;7;-7)$}{$\vec{v}=(7;-7;7)$}{}{}}
\LANGcs{
\Question{
Jsou dány vektory $\vec{a}=(-1; 2;3)$, $\vec{b}=(3; 1; -2)$ a $\vec{c}=(1; 2;-1)$. Určete souřadnice vektoru $\vec{v}$, který je kolmý k oběma daným vektorům $\vec{a}$ a $\vec{b}$, přičemž platí $\vec{v}\cdot\vec{c}=12$.}
{$\vec{v}=(-6;6;-6)$}{$\vec{v}=(6;-6;6)$}{$\vec{v}=(-7;7;-7)$}{$\vec{v}=(7;-7;7)$}{}{}}
\LANGsk{
\Question{
Sú dané vektory $\vec{a}=(-1; 2;3)$, $\vec{b}=(3; 1; -2)$ a $\vec{c}=(1; 2;-1)$. Určte súradnice vektora $\vec{v}$, ktorý je kolmý k obom daným vektorom $\vec{a}$ a $\vec{b}$, pričom platí $\vec{v}\cdot\vec{c}=12$.}
{$\vec{v}=(-6;6;-6)$}{$\vec{v}=(6;-6;6)$}{$\vec{v}=(-7;7;-7)$}{$\vec{v}=(7;-7;7)$}{}{}}
\LANGes{
\Question{
Dados los vectores $\vec{a}=(-1; 2;3)$, $\vec{b}=(3; 1; -2)$ y $\vec{c}=(1; 2;-1)$. Determina las coordenadas de un vector $\vec{v}$, suponiendo que $\vec{v}$ es perpendicular a ambos vectores $\vec{a}$ y $\vec{b}$, mientras $\vec{v}\cdot\vec{c}=12$ .}
{$\vec{v}=(-6;6;-6)$}{$\vec{v}=(6;-6;6)$}{$\vec{v}=(-7;7;-7)$}{$\vec{v}=(7;-7;7)$}{}{}}
\End

\Data\ID{1003040205}
\Updated{2021-06-18 04:06:29}
\Level{C}
\SubArea{Points and vectors}
\LANGen{
\Question{
We are given the vectors $\vec{a}=(1;-2;-2)$, $\vec{b}=(0;1;3)$  and $\vec{c}=(1;-1;0)$. Find the mixed product $(\vec{a}\times\vec{b})\cdot\vec{c}$.}
{$-1$}{$(1;-2;-2)$}{The mixed product is not defined.}{$(-8;8;0)$}{}{}}
\LANGpl{
\Question{
Dane są wektory $\vec{a}=(1;-2;-2)$, $\vec{b}=(0;1;3)$  i $\vec{c}=(1;-1;0)$. Wyznacz $(\vec{a}\times\vec{b})\cdot\vec{c}$.}
{$(\vec{a}\times\vec{b})\cdot\vec{c}=-1$}{$(\vec{a}\times\vec{b})\cdot\vec{c}=(1;-2;-2)$}{$(\vec{a}\times\vec{b})\cdot\vec{c}$ jest nieokreślone}{$(\vec{a}\times\vec{b})\cdot\vec{c}=(-8;8;0)$}{}{}}
\LANGcs{
\Question{
Jsou dány vektory $\vec{a}=(1;-2;-2)$, $\vec{b}=(0;1;3)$  a $\vec{c}=(1;-1;0)$. Vypočtěte smíšený součin $(\vec{a}\times\vec{b})\cdot\vec{c}$.}
{$-1$}{$(1;-2;-2)$}{Smíšený součin není definován.}{$(-8;8;0)$}{}{}}
\LANGsk{
\Question{
Sú dané vektory $\vec{a}=(1;-2;-2)$, $\vec{b}=(0;1;3)$  a $\vec{c}=(1;-1;0)$. Vypočítajte zmiešaný súčin $(\vec{a}\times\vec{b})\cdot\vec{c}$.}
{$(\vec{a}\times\vec{b})\cdot\vec{c}=-1$}{$(\vec{a}\times\vec{b})\cdot\vec{c}=(1;-2;-2)$}{$(\vec{a}\times\vec{b})\cdot\vec{c}$ is not defined}{$(\vec{a}\times\vec{b})\cdot\vec{c}=(-8;8;0)$}{}{}}
\LANGes{
\Question{
Dados los vectores $\vec{a}=(1;-2;-2)$, $\vec{b}=(0;1;3)$  y $\vec{c}=(1;-1;0)$. Determina el producto mixto $(\vec{a}\times\vec{b})\cdot\vec{c}$.}
{$(\vec{a}\times\vec{b})\cdot\vec{c}=-1$}{$(\vec{a}\times\vec{b})\cdot\vec{c}=(1;-2;-2)$}{$(\vec{a}\times\vec{b})\cdot\vec{c}$ no es definido}{$(\vec{a}\times\vec{b})\cdot\vec{c}=(-8;8;0)$}{}{}}
\End

\Data\ID{1003040207}
\Updated{2020-07-15 03:07:12}
\Level{C}
\SubArea{Points and vectors}
\LANGen{
\Question{
Given the points $A = [2;0;3]$ and $B = [-1;2;0]$, specify all the points $C$ lying on the $z$-axis, such that the area of the triangle $ABC$ is $2\sqrt2$.
Hint: Use a cross product of vectors.}
{$C_1=[0;0;1];\ C_2=\left[0;0;\frac{29}{13}\right]$}{$C_1=[0;0;1];\ C_2=\left[0;0;-1\right]$}{$C_1=[0;0;-1];\ C_2=\left[0;0;\frac{13}{29}\right]$}{$C_1=[0;0;-1];\ C_2=\left[0;0;\frac{29}{13}\right]$}{}{}}
\LANGpl{
\Question{
Dane są wektory $A = [2;0;3]$ i $B = [-1;2;0]$, określ wszystkie punkty $C$ leżące na osi $z$ tak, aby  powierzchnia trójkąta $ABC$ była równa $2\sqrt2$.
Wskazówka: Użyj iloczynu wektorowego.}
{$C_1=[0;0;1];\ C_2=\left[0;0;\frac{29}{13}\right]$}{$C_1=[0;0;1];\ C_2=\left[0;0;-1\right]$}{$C_1=[0;0;-1];\ C_2=\left[0;0;\frac{13}{29}\right]$}{$C_1=[0;0;-1];\ C_2=\left[0;0;\frac{29}{13}\right]$}{}{}}
\LANGcs{
\Question{
Jsou dány body $A = [2;0;3]$ a $B = [-1;2;0]$. Určete souřadnice všech takových bodů $C$ ležících na ose $z$, aby obsah trojúhelníku $ABC$ byl $2\sqrt2$.
Nápověda: Užijte vektorový součin vektorů.}
{$C_1=[0;0;1];\ C_2=\left[0;0;\frac{29}{13}\right]$}{$C_1=[0;0;1];\ C_2=\left[0;0;-1\right]$}{$C_1=[0;0;-1];\ C_2=\left[0;0;\frac{13}{29}\right]$}{$C_1=[0;0;-1];\ C_2=\left[0;0;\frac{29}{13}\right]$}{}{}}
\LANGsk{
\Question{
Sú dané body $A = [2;0;3]$ a $B = [-1;2;0]$. Určte súradnice všetkých takých bodov $C$ ležiacich na osi $z$, aby obsah trojuholníka $ABC$ bol $2\sqrt2$.
Nápoveda: Použite vektorový súčin vektorov.}
{$C_1=[0;0;1];\ C_2=\left[0;0;\frac{29}{13}\right]$}{$C_1=[0;0;1];\ C_2=\left[0;0;-1\right]$}{$C_1=[0;0;-1];\ C_2=\left[0;0;\frac{13}{29}\right]$}{$C_1=[0;0;-1];\ C_2=\left[0;0;\frac{29}{13}\right]$}{}{}}
\LANGes{
\Question{
Dados los puntos $A = [2;0;3]$ y $B = [-1;2;0]$. Determina las coordenadas de todos los puntos $C$ que se encuentran en el eje $z$, para que el área del triángulo $ABC$ sea $2\sqrt2$.  

Pista: Usa un producto vectorial.}
{$C_1=[0;0;1];\ C_2=\left[0;0;\frac{29}{13}\right]$}{$C_1=[0;0;1];\ C_2=\left[0;0;-1\right]$}{$C_1=[0;0;-1];\ C_2=\left[0;0;\frac{13}{29}\right]$}{$C_1=[0;0;-1];\ C_2=\left[0;0;\frac{29}{13}\right]$}{}{}}
\End

\Data\ID{1103040202}
\Updated{2020-07-15 03:07:12}
\Level{C}
\SubArea{Points and vectors}
\IMG{1103040202.pdf}
\LANGen{
\Question{
We are given the points $A = [1 ; -2 ; -3]$, $B = [4 ; 1 ; -1]$, $D = [-3 ; 3 ; 1]$ and $E = [2 ; 0 ; 5]$ (see the picture).
Find the volume of the pyramid $ABCDE$  with the parallelogram base $ABCD$ and the apex $E$.}
{$V=\frac{178}3$}{$V=\frac{89}3$}{$V=178$}{$V=89$}{}{}}
\LANGpl{
\Question{
Dane są punkty $A = [1 ; -2 ; -3]$, $B = [4 ; 1 ; -1]$, $D = [-3 ; 3 ; 1]$ i $E = [2 ; 0 ; 5]$ (spójrz na rysunek).
Oblicz objętość ostrosłupa $ABCDE$, którego podstawą jest równoległobok  $ABCD$ i wierzchołku $E$.}
{$V=\frac{178}3$}{$V=\frac{89}3$}{$V=178$}{$V=89$}{}{}}
\LANGcs{
\Question{
Jsou dány body $A = [1 ; -2 ; -3]$, $B = [4 ; 1 ; -1]$, $D = [-3 ; 3 ; 1]$ a $E = [2 ; 0 ; 5]$ (viz obrázek).
Vypočtěte objem jehlanu $ABCDE$  s rovnoběžníkovou podstavou $ABCD$ a vrcholem $E$.}
{$V=\frac{178}3$}{$V=\frac{89}3$}{$V=178$}{$V=89$}{}{}}
\LANGsk{
\Question{
Sú dané body $A = [1 ; -2 ; -3]$, $B = [4 ; 1 ; -1]$, $D = [-3 ; 3 ; 1]$ a $E = [2 ; 0 ; 5]$ (viď obrázok).
Vypočítajte objem ihlanu $ABCDE$  s rovnobežníkovou podstavou $ABCD$ a vrcholom $E$.}
{$V=\frac{178}3$}{$V=\frac{89}3$}{$V=178$}{$V=89$}{}{}}
\LANGes{
\Question{
Dados los puntos $A = [1 ; -2 ; -3]$, $B = [4 ; 1 ; -1]$, $D = [-3 ; 3 ; 1]$ y $E = [2 ; 0 ; 5]$ (mira la imagen).
Determina el volumen de la pirámide $ABCDE$  con la base del paralelogramo $ABCD$ y el vértice $E$.}
{$V=\frac{178}3$}{$V=\frac{89}3$}{$V=178$}{$V=89$}{}{}}
\End

\Data\ID{1103040203}
\Updated{2020-07-15 03:07:12}
\Level{C}
\SubArea{Points and vectors}
\IMG{1103040203.pdf}
\LANGen{
\Question{
We are given the points  $A = [1 ; -2 ; 3]$, $B = [1 ; -2 ; -1]$, $C = [6 ; 10 ; -1]$ and $D = [4 ; -2 ; 3]$.
Find the volume of the tetrahedron $ABCD$ shown in the picture.}
{$V=24$}{$V=48$}{$V=72$}{$V=16$}{}{}}
\LANGpl{
\Question{
Dane są punkty  $A = [1 ; -2 ; 3]$, $B = [1 ; -2 ; -1]$, $C = [6 ; 10 ; -1]$ and $D = [4 ; -2 ; 3]$.
Oblicz objętość czworościanu $ABCD$ pokazanego na rysunku.}
{$V=24$}{$V=48$}{$V=72$}{$V=16$}{}{}}
\LANGcs{
\Question{
Jsou dány body  $A = [1 ; -2 ; 3]$, $B = [1 ; -2 ; -1]$, $C = [6 ; 10 ; -1]$ a $D = [4 ; -2 ; 3]$.
Vypočtěte objem čtyřstěnu $ABCD$ znázorněného na obrázku.}
{$V=24$}{$V=48$}{$V=72$}{$V=16$}{}{}}
\LANGsk{
\Question{
Sú dané body  $A = [1 ; -2 ; 3]$, $B = [1 ; -2 ; -1]$, $C = [6 ; 10 ; -1]$ a $D = [4 ; -2 ; 3]$.
Vypočítajte objem štvorstena $ABCD$ znázorneného na obrázku.}
{$V=24$}{$V=48$}{$V=72$}{$V=16$}{}{}}
\LANGes{
\Question{
Dados los puntos $A = [1 ; -2 ; 3]$, $B = [1 ; -2 ; -1]$, $C = [6 ; 10 ; -1]$ y $D = [4 ; -2 ; 3]$.
Calcula el volumen del tetraedro $ABCD$ mostrado en la imagen.}
{$V=24$}{$V=48$}{$V=72$}{$V=16$}{}{}}
\End

\Data\ID{1103040204}
\Updated{2020-07-15 03:07:12}
\Level{C}
\SubArea{Points and vectors}
\IMG{1103040204.pdf}
\LANGen{
\Question{
We are given the points $A = [1;2;1]$, $B = [7;3;0]$, $C = [-1;5;2]$ and $D = [1;0;6]$.
Find the volume of the triangular prism $ABCDEF$ shown in the picture.}
{$V=54$}{$V=108$}{$V=36$}{$V=56$}{}{}}
\LANGpl{
\Question{
Dane są punkty $A = [1;2;1]$, $B = [7;3;0]$, $C = [-1;5;2]$ i $D = [1;0;6]$.
Oblicz objętość  graniastosłupa trójkątnego $ABCDEF$ przedstawionego na rysunku.}
{$V=54$}{$V=108$}{$V=36$}{$V=56$}{}{}}
\LANGcs{
\Question{
Jsou dány body $A = [1;2;1]$, $B = [7;3;0]$, $C = [-1;5;2]$ a $D = [1;0;6]$.
Vypočtěte objem trojbokého hranolu $ABCDEF$ znázorněného na obrázku.}
{$V=54$}{$V=108$}{$V=36$}{$V=56$}{}{}}
\LANGsk{
\Question{
Sú dané body $A = [1;2;1]$, $B = [7;3;0]$, $C = [-1;5;2]$ a $D = [1;0;6]$.
Vypočítajte objem trojbokého hranolu $ABCDEF$ znázorneného na obrázku.}
{$V=54$}{$V=108$}{$V=36$}{$V=56$}{}{}}
\LANGes{
\Question{
Dados los puntos $A = [1;2;1]$, $B = [7;3;0]$, $C = [-1;5;2]$ and $D = [1;0;6]$.
Calcula el volumen del prisma triangular  $ABCDEF$ mostrado en la imagen.}
{$V=54$}{$V=108$}{$V=36$}{$V=56$}{}{}}
\End

\Data\ID{1103040206}
\Updated{2020-07-15 03:07:12}
\Level{C}
\SubArea{Points and vectors}
\IMG{1103040206.pdf}
\LANGen{
\Question{
Given the points $A = [1;5]$ and $B = [-4;2]$, specify all the points $C$ lying on the $x$-axis, such that the area of the triangle $ABC$ is $14$.
Hint: Use a cross product of vectors.}
{$C_1=[2;0];\ C_2=\left[-\frac{50}3;0\right]$}{$C_1=[1;0];\ C_2=\left[-\frac{47}3;0\right]$}{$C_1=[2;0];\ C_2=\left[-\frac{47}3;0\right]$}{$C_1=[1;0];\ C_2=\left[-\frac{50}3;0\right]$}{}{}}
\LANGpl{
\Question{
Dane są punkty $A = [1;5]$ i $B = [-4;2]$, określ wszystkie punkty $C$ leżące na osi $x$ tak, aby powierzchnia trójkąta $ABC$ wynosiła $14$.
Wskazówka: Użyj iloczynu wektorowego.}
{$C_1=[2;0];\ C_2=\left[-\frac{50}3;0\right]$}{$C_1=[1;0];\ C_2=\left[-\frac{47}3;0\right]$}{$C_1=[2;0];\ C_2=\left[-\frac{47}3;0\right]$}{$C_1=[1;0];\ C_2=\left[-\frac{50}3;0\right]$}{}{}}
\LANGcs{
\Question{
Jsou dány body $A = [1;5]$ a $B = [-4;2]$. Určete souřadnice všech takových bodů $C$ ležících na ose $x$, aby obsah trojúhelníka $ABC$ byl $14$.
Nápověda: Užijte vektorový součin vektorů.}
{$C_1=[2;0];\ C_2=\left[-\frac{50}3;0\right]$}{$C_1=[1;0];\ C_2=\left[-\frac{47}3;0\right]$}{$C_1=[2;0];\ C_2=\left[-\frac{47}3;0\right]$}{$C_1=[1;0];\ C_2=\left[-\frac{50}3;0\right]$}{}{}}
\LANGsk{
\Question{
Sú dané body $A = [1;5]$ a $B = [-4;2]$. Určte súradnice všetkých takých bodov $C$ ležiacich na osi $x$, aby obsah trojuholníka $ABC$ byl $14$.
Nápoveda: Použite vektorový súčin vektorov.}
{$C_1=[2;0];\ C_2=\left[-\frac{50}3;0\right]$}{$C_1=[1;0];\ C_2=\left[-\frac{47}3;0\right]$}{$C_1=[2;0];\ C_2=\left[-\frac{47}3;0\right]$}{$C_1=[1;0];\ C_2=\left[-\frac{50}3;0\right]$}{}{}}
\LANGes{
\Question{
Dados los puntos $A = [1;5]$ y $B = [-4;2]$, specifica todos los puntos $C$ que se encuentran en el eje $x$, suponiendo que la área del triángulo  $ABC$ es $14$.
Pista: Usa un producto vectorial.}
{$C_1=[2;0];\ C_2=\left[-\frac{50}3;0\right]$}{$C_1=[1;0];\ C_2=\left[-\frac{47}3;0\right]$}{$C_1=[2;0];\ C_2=\left[-\frac{47}3;0\right]$}{$C_1=[1;0];\ C_2=\left[-\frac{50}3;0\right]$}{}{}}
\End

\Data\ID{1103040208}
\Updated{2020-07-15 03:07:12}
\Level{C}
\SubArea{Points and vectors}
\IMG{1103040208_0.pdf}
\LANGen{
\Question{
We are given the points $A = [4;5;-1]$, $B = [-2;-1;2]$, $C = [-1;-3;0]$ and  $D = [0;m;2]$.  Find the missing coordinate of the point $D$ such that the point $D$ lies in the plane determined by the points $A$, $B$ and $C$.
Hint: Use a linear combination of vectors shown in the picture or use their mixed product.}
{$m=3$}{$m=-3$}{$m=1$}{$m$ does not exist}{}{}}
\LANGpl{
\Question{
Dane są punkty $A = [4;5;-1]$, $B = [-2;-1;2]$, $C = [-1;-3;0]$ i  $D = [0;m;2]$.  Wyznacz brakujące współrzędne punktu $D$ tak, aby punkt $D$ leżał na płaszczyźnie określonej przez punkty $A$, $B$ i $C$.
Wskazówka: Użyj liniowej kombinacji wektorów pokazanych na obrazku lub użyj ich iloczynu mieszanego.}
{$m=3$}{$m=-3$}{$m=1$}{$m$ does not exist}{}{}}
\LANGcs{
\Question{
Jsou dány body $A = [4;5;-1]$, $B = [-2;-1;2]$, $C = [-1;-3;0]$ a  $D = [0;m;2]$.  Určete chybějící souřadnici bodu $D$ tak, aby bod $D$ ležel v rovině určené body $A$, $B$ a $C$.
Nápověda: Užijte lineární kombinaci vektorů vyznačených na obrázku nebo užijte jejich smíšený součin.}
{$m=3$}{$m=-3$}{$m=1$}{$m$ does not exist}{}{}}
\LANGsk{
\Question{
Sú dané body $A = [4;5;-1]$, $B = [-2;-1;2]$, $C = [-1;-3;0]$ a  $D = [0;m;2]$.  Určte chýbajúcu súradnicu bodu $D$ tak, aby bod $D$ ležal v rovine určenej bodmi $A$, $B$ a $C$.
Nápoveda: Použite lineárnu kombináciu vektorov vyznačených na obrázku alebo použite ich zmiešaný súčin.}
{$m=3$}{$m=-3$}{$m=1$}{$m$ neexistuje}{}{}}
\LANGes{
\Question{
Dados los puntos $A = [4;5;-1]$, $B = [-2;-1;2]$, $C = [-1;-3;0]$ y  $D = [0;m;2]$.  Halla la coordenada faltante de un punto $D$ suponiendo que el punto $D$ se encuentra en el plano determinado por los puntos $A$, $B$ y $C$.
Pista: Usa una combinación lineal de vectores que está marcada en la imagen o usa su producto mixto.}
{$m=3$}{$m=-3$}{$m=1$}{$m$ does not exist}{}{}}
\End

\Data\ID{9000108707}
\Updated{2020-07-15 03:07:37}
\Level{C}
\SubArea{Points and vectors}
\IMG{001087_07-figure0.pdf}
\LANGen{
\Question{
Find the area of the parallelogram \(ABCD\) 
where \(A = [1;3;2]\),
\(B = [3;4;6]\),
\(D = [2;5;8]\).}
{\(\sqrt{77}\)}{\(8\)}{\(2\sqrt{13}\)}{\(\sqrt{94}\)}{}{}}
\LANGpl{
\Question{
Oblicz pole równoległoboku \(ABCD\) jeśli
 \(A = [1;3;2]\),
\(B = [3;4;6]\),
\(D = [2;5;8]\).}
{\(\sqrt{77}\)}{\(8\)}{\(2\sqrt{13}\)}{\(\sqrt{94}\)}{}{}}
\LANGcs{
\Question{
Vypočítejte obsah rovnoběžníku
\(ABCD\), je-li
\(A = [1;3;2]\),
\(B = [3;4;6]\),
\(D = [2;5;8]\).}
{\(\sqrt{77}\)}{\(8\)}{\(2\sqrt{13}\)}{\(\sqrt{94}\)}{}{}}
\LANGsk{
\Question{
Vypočítajte obsah rovnobežníka
\(ABCD\), ak je
\(A = [1;3;2]\),
\(B = [3;4;6]\),
\(D = [2;5;8]\).}
{\(\sqrt{77}\)}{\(8\)}{\(2\sqrt{13}\)}{\(\sqrt{94}\)}{}{}}
\LANGes{
\Question{
Halla la área de un paralelogramo \(ABCD\) 
donde \(A = [1;3;2]\),
\(B = [3;4;6]\),
\(D = [2;5;8]\).}
{\(\sqrt{77}\)}{\(8\)}{\(2\sqrt{13}\)}{\(\sqrt{94}\)}{}{}}
\End

\Data\ID{9000108708}
\Updated{2020-07-15 03:07:37}
\Level{C}
\SubArea{Points and vectors}
\IMG{001087_08-figure0.pdf}
\LANGen{
\Question{
Find the volume of the parallelepiped \(ABCDEFGH\) 
where \(A = [1;0;0]\),
\(B = [2;0;0]\),
\(D = [3;-2;0]\),
\(E = [2;1;5]\).}
{\(10\)}{\(12\)}{\(15\)}{\(20\)}{}{}}
\LANGpl{
\Question{
Oblicz objętość równoległościanu \(ABCDEFGH\), gdzie \(A = [1;0;0]\),
\(B = [2;0;0]\),
\(D = [3;-2;0]\),
\(E = [2;1;5]\).}
{\(10\)}{\(12\)}{\(15\)}{\(20\)}{}{}}
\LANGcs{
\Question{
Určete objem rovnoběžnostěnu \(ABCDEFGH\),
je-li \(A = [1;0;0]\),
\(B = [2;0;0]\),
\(D = [3;-2;0]\),
\(E = [2;1;5]\).}
{\(10\)}{\(12\)}{\(15\)}{\(20\)}{}{}}
\LANGsk{
\Question{
Určte objem rovnobežnostena \(ABCDEFGH\),
ak \(A = [1;0;0]\),
\(B = [2;0;0]\),
\(D = [3;-2;0]\),
\(E = [2;1;5]\).}
{\(10\)}{\(12\)}{\(15\)}{\(20\)}{}{}}
\LANGes{
\Question{
Calcula el volumen de un paralelepípedo\(ABCDEFGH\) 
donde \(A = [1;0;0]\),
\(B = [2;0;0]\),
\(D = [3;-2;0]\),
\(E = [2;1;5]\).}
{\(10\)}{\(12\)}{\(15\)}{\(20\)}{}{}}
\End
